\section{Exact precision computation and polynomial compilation}\label{sec:compiler}
Encode a signed integer by
\[
 \operatorname{enc}(n)=\begin{cases}2n,&n\ge0,\\-2n-1,&n<0.\end{cases}
\]
The five natural inputs decode in the fixed order $(n,\kappa,a,l,j)$.
The decision predicate rejects unless
$n\notin\Cset$, $\kappa\ne0$, $a,l\ge1$, and $j=na$.
For nonzero multiples of $P_0$, $z=-x/y$ is well-defined. Set
\[
 z([l]P_0)=u_l+v_ls,\qquad\epsilon_l=|u_l|+|v_l|,
 \qquad r(a,j)=\tfrac12\Tr_{K_0/\Q}\frac{z([j]P_0)}{z([a]P_0)}.
\]
With the fixed constant $A_{\rm ht}^{(0)}=H_6(S_{\rm ht})$, its final
tests are
\[
 |r(a,j)-n|<\epsilon_l,\qquad
 4\epsilon_l<4^{-A_{\rm ht}^{(0)}|\kappa F(n)|^{24}}.
\]
These are strict tests on exactly computable rationals and integers.

The reference program S21 uses exact quadratic-field arithmetic and capped
integer operations. The ideal register intermediate representation S22
computes the same predicate on every five-natural tuple. Its independent
all-input argument covers normalized signed arithmetic, origin and
zero-ordinate cases, equal-abscissa addition, and every early rejection.
Nine exact polynomial identities verify the nonexceptional projective
formulas and preservation of $Y^2=X^3+AXZ^4$. Scalar loops decrease by
halving; the precision loop divides by four; factorial iteration is bounded.
This is a proof about the ideal intermediate representation, rather than
an assertion that a budgeted simulator always finishes.

For example, if $\epsilon=U/V>0$ with natural $U,V$ and integer $B\ge0$,
\[
 4\epsilon<4^{-B}\quad\Longleftrightarrow\quad
 4^B\le\left\lfloor\frac{V-1}{4U}\right\rfloor.
\]
This permits a finite cap. For positive $L$, $\sigma(x)=\min(x,L+1)$
commutes with saturated addition and multiplication. Powers and factorials
can be computed with this saturation. Capping an exponent is safe for
base at least two because an exponent above $L$ already makes the value
exceed $L$. Base zero, base one and exponent zero are separate cases.
This proves equality with the fixed-height expression without requiring
its full evaluation when the available precision is too small.

\subsection{The natural polynomial and the power interface}
The saved precision polynomial has $1\,305\,811$ natural variables,
the first five being inputs. Its SHA256 is
\begin{center}\small\ttfamily
f3e163f5c60b972c98b48647274f956dd7d78428e27ddc0543574a5f18120b8c.
\end{center}
Source S23 specializes a fixed universal natural halting polynomial at
the code of a wrapper. The wrapper inverts four Cantor pairs and then
runs the unchanged precision machine. Its pairing equations are
\[
 2t_i=(c_i+t_{i+1})(c_i+t_{i+1}+1)+2t_{i+1}\quad(0\le i<4),
 \qquad t_4=c_4.
\]
Their triangular intervals prove unique natural unpairing. The new
polynomial sums the square of the old halting polynomial and the four
squared differences above.

The universal compiler packs a finite run of a fixed $m$-state
increment/decrement/branch program into base-$B$ integers. Acceptance
means reaching the unique final halt with all registers zero; the fixed
interpreter includes a finite cleanup stage. Choose
\[
 c>0,\quad K=2^c,\quad B=4K>m,\quad T=B^q,\quad
 F_q=1+B+\cdots+B^q,\quad D=(K-1)F_q,
\]
and require every input to be below $K$. For natural integers, say
$X$ is masked by $Y$ when every set binary bit of $X$ occurs in $Y$.
Require each register history $R_j$ to be masked by $D$ and satisfy
$R_j<T$. Each state history $S_i$ is masked by $F_q$, and
\[
 \sum_{i=0}^{m-1}S_i=F_q,\qquad S_{\rm halt}=T.
\]
Thus register digits are below $K$, terminal register digits are zero,
and state digits are zero or one. Since $B>m$, the state sum has no
carries and selects exactly one state at each time.

The nonzero mask $N_j$ is determined by
\[
 K N_j=(R_j+D)\mathbin{\mathrm{AND}}(K F_q).
\]
Indeed each digit $r+K-1<2K<B$ has its $K$-bit set exactly when
$r>0$. For a decrement/branch instruction at state $i$ testing
register $j$, impose
\[
 V_i=S_i\mathbin{\mathrm{AND}}N_j,\qquad S_i=V_i+W_i.
\]
These are the nonzero and zero branch histories. Form the packed
increment indicator $A_j$ from states incrementing $j$, the decrement
indicator $C_j$ from the appropriate $V_i$, and each incoming-state
indicator $I_i$ from the states or branches directed to $i$.
They have zero-one digits, since there is just one active instruction
and branch at each time. The transition equations are
\[
 R_j+BC_j=\operatorname{initial}_j+BR_j+BA_j,\qquad
 S_i=\mathbf1_{i=0}+BI_i.
\]
Every register coefficient on either side is at most $K<B$;
uniqueness of base-$B$ digits gives the initial conditions and each
exact transition. The possible extra final slot vanishes because
halt has no outgoing instruction and the terminal registers are zero.
For $q=0$ the same equations require the initial state already to be
halt and the inputs to vanish. Conversely, any finite accepting run
fits after choosing $c$ sufficiently large, and supplies all these
masks and equations.

\subsection{The literal alpha and power formulas}
The alpha system below is \texttt{Exp\_Matrices.alpha\_equations},
with equivalence \texttt{alpha\_equivalence}; the power system is
\texttt{Exponential\_Relation.exp\_repr} and its polynomial expansion
\texttt{exp\_dioph} in the pinned AFP source \cite{afp}.
Their variables range over natural numbers. The alpha recurrence itself
is integer-valued; the bounds below justify its natural-valued uses.
Define $U_n=\alpha_b(n)$ by
\[
 U_0=0,\qquad U_1=1,\qquad U_{n+2}=bU_{n+1}-U_n\quad(n\ge0).
\]
For $b\ge3$ this recurrence gives
$0\le U_n<U_{n+1}$ and $U_n\ge n$. For $n\ge1$, the matrix
\[
 M=\begin{pmatrix}b&-1\\1&0\end{pmatrix},\qquad
 M^n=\begin{pmatrix}U_{n+1}&-U_n\\U_n&-U_{n-1}\end{pmatrix}
\]
has determinant one. Thus adjacent terms are coprime and
$U_n^2-bU_nU_{n-1}+U_{n-1}^2=1$.
Parameter congruence follows directly from the recurrence:
$b\equiv b'\pmod d$ implies
$\alpha_b(n)\equiv\alpha_{b'}(n)\pmod d$.
In particular $\alpha_2(n)=n$. This is parameter congruence,
not merely congruence between indices.

Every natural solution of $x^2+y^2=1+bxy$, $b>2$, is an unordered
adjacent pair. To prove this, order it as $X>Y$. If $Y=0$, $X=1$.
Otherwise $Z=bY-X$ satisfies $XZ=Y^2-1$, hence $0\le Z<Y$,
and $(Y,Z)$ satisfies the same equation. Descent reaches $(1,0)$;
reversing the steps generates exactly the adjacent pairs.
No ordering assumption on the original pair is needed.

For $k>0$, reduction of $M^k$ modulo $U_k$ is the invertible scalar
$U_{k+1}I$. Dividing an index by $k$ and using strict growth proves
$U_k\mid U_m$ if and only if $k\mid m$. More precisely, for $m>0$,
\[
                    U_k^2\mid U_m\quad\Longleftrightarrow\quad
                         kU_k\mid m.                           \tag{C.1}
\]
For $k=0$ both sides are false. For $k>0$ write $m=kq$ and expand
$M^{kq}=(U_kM-U_{k-1}I)^q$ modulo $U_k^2$. Its lower-left entry is
$qU_k(-U_{k-1})^{q-1}$, so coprimality proves (C.1), including $U_k=1$.

For $m>0$ put $v=U_{m+1}-U_{m-1}=\operatorname{tr}(M^m)$.
Cayley--Hamilton gives $M^{2m}\equiv-I\pmod v$.
Extend the recurrence by $U_{-j}=-U_j$. Every $n\ge0$ has
$n=2lm+\varepsilon j$ with $l\ge0$, $0\le j\le m$ and
$\varepsilon\in\{1,-1\}$; then
$U_n\equiv(-1)^l\varepsilon U_j\pmod v$.
The value sign need not equal the index sign. If $b\ge4$,
$v=bU_m-2U_{m-1}>2U_m$.

The actual alpha formula has natural witnesses $r,s,t,u,v,w,x,y$
and the following fifteen conditions:
\[
\begin{gathered}
 b>3,\quad u^2+t^2=1+but,\quad s^2+r^2=1+bsr,\quad r<s,\quad
 u^2\mid s,\quad v+2r=bs,\\
 w\equiv b\pmod v,\quad w\equiv2\pmod u,\quad w>2,\quad
 x^2+y^2=1+wxy,\\
 2X<u,\quad2X<v,\quad x\equiv X\pmod v,\quad
 2N<u,\quad x\equiv N\pmod u.
\end{gathered}                                                   \tag{C.2}
\]
A strict inequality $a<b$ is $b=a+1+d$; divisibility is $b=ad$;
congruence is $A+Mu=B+Mv$, with fresh natural witnesses.
At zero modulus the last equation says equality, so there is no
exception hidden in this replacement.

For soundness, the unordered Pell result gives $u=U_k$, $k>0$,
and $x=\alpha_w(n)$. The ordered pair gives $s=U_m,r=U_{m-1}$,
$m>0$. By (C.1), $u\mid m$. Parameter congruence gives
$x\equiv U_n\pmod v$ and $x\equiv n\pmod u$.
Reflect $n$ to $j\in[0,m]$. Since $2X<v$ and $2U_j<v$,
the congruence $X\equiv\pm U_j\pmod v$ forces $X=U_j$,
with the negative case forcing both zero.
Then $j\le X$, $2j<u$ and $2N<u$ force $N=j$ from
$N\equiv\pm j\pmod u$. Thus $X=U_N$.

Conversely choose an odd $u=U_k>2U_N$. Such terms are unbounded,
since adjacent terms are coprime and the sequence grows.
Set $t=U_{k+1}$, $m=ku$, $s=U_m$, $r=U_{m-1}$ and $v=bs-2r$.
Equation (C.1) gives $u^2\mid s$, and the required strict bounds hold.
The moduli $u,v$ are coprime: a common divisor is odd, divides $s$
and $2r$, hence divides both adjacent terms. CRT gives
$w\equiv2\pmod u$, $w\equiv b\pmod v$ with $w>2$.
Take $x=\alpha_w(N)$, $y=\alpha_w(N+1)$.
Parameter congruence verifies every remaining condition.
This includes $N=X=0$.

The literal power formula separates $Q=0,R=0,P=1$ and
$Q=0,R>0,P=0$. For $Q\ge1$ it imposes
\[
 a_0=\alpha_{Q+4}(R+1),\quad b=a_0+Q^2+2,\quad
 m+Q^2+1=bQ,\quad P<m,
\]
\[
 a_1=\alpha_b(R),\quad a_2=\alpha_b(R+1),\quad
 P+ba_1\equiv Qa_1+a_2\pmod m.                                  \tag{C.3}
\]
Induction gives $a_0\ge(Q+3)^R$, so $m>Q^R$ and $m>0$.
Let $E_R=Q U_R-bU_R+U_{R+1}$. The recurrence and
$bQ\equiv Q^2+1\pmod m$ give $E_0=1$ and
$E_{R+1}\equiv QE_R\pmod m$. Thus (C.3) says $P\equiv Q^R\pmod m$.
Both values lie in $[0,m)$, forcing equality. Conversely equality
supplies every witness. This includes $R=0$, $Q=1$ and $0^0=1$.

These formulas make the binomial and bit operations polynomial.
For $B=2^{n+1}$, the equation
\[
 (1+B)^n=B^k(Bq+C)+r,\quad C<B,\quad r<B^k
\]
selects exactly $\binom nk$, since all its base-$B$ coefficients
are below $B$, including zero for $k>n$.
The identity
$(1+t)^n=\prod_{i:\,\text{bit}_i(n)=1}(1+t^{2^i})$ over $\F_2$
makes odd binomial coefficients equivalent to bit inclusion.
Bit disjointness of $a,b$ is equivalent to inclusion of $a$ in $a+b$:
at the least shared bit there is no incoming carry and that bit
is cleared. Finally $x=z+a,y=z+b$ with $z,a,b$ pairwise bit-disjoint
says exactly $z=x\mathbin{\rm AND}y$.

\subsection{Interpreter, lowering and five-input specialization}
The universal program numbering is $e+1=2^w(2D+1)$.
At $w=0$ it halts. Otherwise the four base-$2^w$ digits at positions
$4p,\ldots,4p+3$ give the opcode, register and two targets at state $p$.
Opcode one increments, opcode two decrements or takes its zero branch;
every other opcode halts. Missing digits are zero, so decoding is total.

The interpreter stores a finite tape as
$0$ or $\operatorname{cons}(h,t)=2^h(2t+1)$.
Repeated halving decodes a nonempty head and tail; the tail is smaller
than the code. To update a register, save the finite reversed prefix,
replace the entry and restore the prefix. Absent entries are zero.
The reserved jump register remains zero. Copies restore their source;
aliased addition snapshots its operand before overwriting it.
Division maintains $x=qy+r$ and subtracts only when $r\ge y$;
the interpreter's divisors are at least two.
These invariants prove that one interpreter loop executes exactly one
decoded instruction, with finite work. A simulated halt enters cleanup
and clears all registers. An infinite simulation never enters cleanup.
Thus its all-zero halt is equivalent to virtual halting, without
requiring the virtual final tape itself to be zero.

The precision IR uses normalized signed pairs for integer arithmetic
and pairs over the basis $1,\sqrt2$ for field arithmetic.
Norm zero implies the field element is zero: after clearing
denominators, a primitive solution of $a^2=2b^2$ would have both
entries even. Every division is guarded by that exact test.
The Jacobian point formulas satisfy the curve equation and the affine
chord or tangent formulas after multiplying denominators.
Equal abscissas split into doubling and opposite-point branches;
zero ordinate doubles to the identity. The fixed projective point
is the same affine $P_0$. These are polynomial identities, independent
of the generator theorem.

The scalar loop halves its remaining count. The precision loop divides
its remainder by four. The capped arithmetic computes exactly
$\min(x,L+1)$; zero bases, unit bases and zero exponents are separate.
For bases at least two an exponent exceeding $L$ already exceeds the
cap. Thus the ideal IR terminates and agrees with the stated rational
predicate on all five natural inputs, including every early rejection.

Each primitive lowering uses fresh auxiliary registers, snapshots
overlapping sources, and preserves the live-register map at the end
of each finite block. Saturating decrement gives natural subtraction,
and repeated addition gives multiplication. All variable divisions
are guarded. Induction on completed loop iterations proves the same
claim for loops. Acceptance enters cleanup; rejection enters a
permanent self-loop. Hence the primitive machine halts exactly on
acceptance, although the ideal decision program is total.

For $\operatorname{pair}(a,b)=(a+b)(a+b+1)/2+b$, unpairing maintains
$\text{remainder}=\text{input}-s(s+1)/2$ and subtracts $s+1$
while $s<\text{remainder}$. At exit
$0\le\text{remainder}\le s$, so $b=\text{remainder}$, $a=s-b$.
Four repetitions put the five codes into registers zero through four.
Fresh scratch lies above the old register range and is cleared;
old work registers remain zero. Relocating every target by the prefix
length therefore preserves the precision machine.

The polynomial specialization fixes the universal program input to
this wrapper code, maps the old packed input to variable five, and
maps every old variable $i\ge2$ and private-block start by $i\mapsto i+7$.
The four pairing variables occupy five through eight, disjoint from
the five inputs. The appended four squared equations force their
unique Cantor values. Thus, for the exact saved polynomial $C$,
\[
 \forall\mathbf c\in\N^5,\qquad
 \exists\mathbf z\in\N^{1305806}\ C(\mathbf c,\mathbf z)=0
 \quad\Longleftrightarrow\quad
        \operatorname{accepts\_naturals}(\mathbf c)=\mathrm{True}. \tag{C.4}
\]
The natural compiler proof has no height, restricted-converse or
rational-model premise. Its use of the height expression concerns
only computation of that fixed integer expression.

\subsection{Translation to point codes}
The typing order is independent of the natural predicate. First,
\texttt{member} forces a canonical code in $D$, and raw addition on two
canonical elliptic inputs forces their canonical sum, including the
origin and two-torsion cases. The fixed scalar graph gives $[m]Q$ on
$D$. In \texttt{multiply}, the three root membership calls are forced
conjuncts inside \texttt{height\_certificate}. Thus the four scalar shifts
are in $D$ before they are passed, unchanged, to the 27 slope calls and
the six test-group calls within each slope. This supplies the hypotheses
of every raw addition in those groups. Only after globally typed
multiplication is established do the four square products and three
successive additions justify \texttt{natural}. Their outputs are canonical
by the raw-addition postcondition, so no separate output membership
assumption is needed. The signed-encoding and exceptional branches start
their addition chains at this natural code, the literal origin, or the
fixed signed generator. An inactive product branch is not assumed true.
The fixed generator equations are conjuncts outside the internal products
in all three templates that use them. These observations give the base
cases for the source-DAG induction below.
Apply (C.4) using the exact signed index model. Induction on the actual
source DAG gives each node the code of its signed integer value.
Natural guards constrain source variables, not intermediate sums or
products. Raw addition receives canonical codes; globally typed
multiplication justifies the natural guards. Fixed generator equations
and scalar templates supply constants.

The power-value recipe has three signed code inputs, seventy-nine
private natural witnesses, and an output code for the value of the
power-template polynomial. That value is not independently required
to vanish. Its saved root equality is $[17751579,15]$.
The precision parent has $54056$ explicit variables and $15845$
disjoint private blocks of seventy-nine variables. No parent variable
node accesses a private coordinate. Replacing each block by its
private power-value call therefore preserves the full existential
assignment. Only the final parent root is equated to the identity.

At each primitive formula, equality is a difference, conjunction is
a sum of squares, and disjunction is a product. A topological induction,
including the globally conjoined fixed equations, proves equality of
polynomial values under substitution, not merely equisatisfiability
of an unrelated formula. This is the required link to the saved AST.

\section{A variable-shift relation for ordinary values}\label{sec:bridge}
We prove the ordinary-value deduction from the exact index model,
the compiled precision relation, the height theorem and $\mathsf{PC}$.
These are supplied by Sections~\ref{sec:model}, \ref{sec:compiler},
\ref{sec:height} and Theorem~\ref{hyp:PC}, respectively. Capital letters below denote five-coordinate point codes;
their lower-case indices are ordinary integers. The external variable
$t$ is any rational number.

\subsection{Approximation for every fixed shift}
Define $\Close(N,\mathcal K,t)$ by choosing $A,L,J$, linking all five
signed codes to precision inputs in the order $(n,\kappa,a,l,j)$, and
imposing $|t-r(a,j)|<\epsilon_l$ on those same coordinates. Put
\[
 B(n,\kappa)=H_6(S_{\rm ht})|\kappa F(n)|^{24}.
\]
The precision contract gives
\begin{align}
 \Close(\iota(n),\iota(\kappa),t)&\Longrightarrow
                  |t-n|<e^{-B(n,\kappa)},\label{eq:close-sound}\\
 \Close(\iota(n),\iota(\kappa),n)&\quad\text{for every }
                  n\notin\Cset,\ \kappa\ne0.\label{eq:close-complete}
\end{align}
The first follows from $|t-n|<2\epsilon_l<4^{-B}/2<e^{-B}$.
For the second, positive multiples of $P_0$ approach the identity
simultaneously in the compact group $E(\mathbf R)\times E(\mathbf R)$.
This follows, for instance, by partitioning the compact closure of the
cyclic subgroup into arbitrarily small neighborhoods and subtracting
two sufficiently close multiples. If the resulting positive differences
were bounded, a fixed nonzero multiple would be the identity; infinite
order excludes this. Choose $l$ first to make the positive $\epsilon_l$
small enough. Since the derivative of $[n]$ at the identity is $n$
\cite[IV.2.3(a), p.~122]{silverman},
$r(a,na)\to n$ along positive multiples $a$ approaching it. Choose such
an $a$ and put $j=na$. This works for every already fixed nonzero shift
$\kappa$, regardless of its size.

\subsection{The six local tests}
For $N=\iota(n)$, $\mathcal K=\iota(\kappa)$, form the code for
$\kappa n$, choose $H=\iota(h)$, $U=\iota(s_0)$, and put
$V_0=H-\mathcal K$, $V_1=U-\iota(\kappa n)$.
Require $\kappa\ne0$, $H,V_0\ne O$. Define $w(Q)=1/x(Q)$ for
$Q\ne O$, $w(O)=0$. For $t\notin\Cset$ use
\begin{align*}
 D_1&=w(H),&D_2&=w(U),&D_3&=\sigma w(V_0),&D_4&=\sigma w(V_1),\\
 D_5&=z(U)/z(H)-t,&D_6&=\sigma(z(V_1)/z(V_0))-t.
\end{align*}
Here $\sigma$ is conjugation in $K_0$. Choose $\eta,q\in\Q$ and impose
\begin{align*}
 \eta^2&=1,&R_0&=\eta2730^6F(t),&a_0&=1-R_0q^2\ne0,\\
 R&=R_0a_0^2,&C_j&=D_j/F(t),&W&=\prod_{j=1}^6(1+RC_j^2)\ne0,\\
 \Theta&=\Tr_{K_0/\Q}(W^{-1})\ne0,&B_{\rm pole}&=R/\Theta^2,
 &\Phi(B_{\rm pole})&.
\end{align*}
Call this relation $\Test(N,\mathcal K,t)$. All quotients are guarded
polynomially. There is one common denominator $F(t)$, one product $W$,
one trace and one pole instance.

\begin{lemma}[Divisibility]\label{lem:test}
If $\Test(\iota(n),\iota(\kappa),t)$ holds, then
$M_{S_{\rm ht}}(t)\mid\kappa F(n)$.
\end{lemma}
\begin{proof}
Let $p\mid M_{S_{\rm ht}}(t)$. Exactly one $c\in\Cset$ has
$\delta=v_p(t-c)>0$ odd. The contact differences are units, so
$v_p(F(t))=\delta$; the field is unramified at $p$. Put $b=v_p(a_0)$.
If $q=0$, then $b=0$. Otherwise $\delta+2v_p(q)$ is odd and
nonzero, so $b=\min(0,\delta+2v_p(q))\le0$. In either case
$\alpha=v_p(R)=\delta+2b$ is odd. At each place above $p$,
\[
 v(1+RC_j^2)=\min(0,\alpha+2(v(D_j)-\delta)).
\]
Let $\lambda=\max_{v\mid p}v(W)\le0$. Then
$v_p(\Theta)\ge-\lambda\ge0$. The pole input has odd valuation
$\alpha-2v_p(\Theta)$, which is positive by Lemma~\ref{lem:pole}.
Therefore $\alpha>0$ and $\lambda> -\alpha/2$.

At a place attaining $\lambda$, each nonpositive factor valuation is
at least $\lambda$. Consequently
\[
 v(D_j)\ge\delta-3\alpha/4=\delta/4-3b/2>0
 \quad(1\le j\le6).
\]
Integral valuations are at least one, at the same place for all tests.
The first four put $H,U$ in one formal group and $V_0,V_1$ in the
conjugate formal group. At the identity we set both $z(O)=0$ and
$\log(O)=0$; thus the cases $U=O$ and $V_1=O$ obey the same ratio
identities below. Their denominators are nonzero by the guards on
$H$ and $V_0$. For a nonzero point in either formal group,
\[
 \log(Q)=z(Q)(1+e_Q),\qquad v(e_Q)\ge1.
\]
Here $p>324$ and $v(p)=1$. By
\cite[IV.5.5 and IV.6.4, pp.~129--132]{silverman},
the logarithm is an injective isometry on $v(z)\ge1$.
The integral invariant differential gives in relative degree $r-1$
the bound $(r-1)v(z)-v_p(r)\ge1$ for every $r\ge2$.
Thus the guarded nonzero points $H,V_0$ have nonzero logarithms.
The fifth test first makes $z(U)/z(H)$ integral and congruent to
$t$; multiplication by the units $1+e_Q$ preserves that congruence.
The relation $[s_0]H=[h]U$ then gives
$\log(U)/\log(H)=s_0/h$. Likewise the sixth gives
$(s_0-\kappa n)/(h-\kappa)\equiv t$.
Multiplying by the integer denominators and subtracting yields
$v_p(\kappa(n-t))\ge1$. Since $t\equiv c$, we get
$p\mid\kappa(n-c)$, hence $p\mid\kappa F(n)$. The squarefree
product divides as asserted. No unit assumption on $\kappa$ was used.
\end{proof}

\begin{lemma}[One suitable shift]\label{lem:test-complete}
For every $n\in\Z\setminus\Cset$ there is
$\kappa\ne0$ with $\Test(\iota(n),\iota(\kappa),n)$.
\end{lemma}
\begin{proof}
Choose the sign so $R_0=\eta2730^6F(n)>0$ and let $S'$ be its finite
prime support. Each sufficiently small local formal neighborhood has
finite index in the compact local group
\cite[VII.6.3, p.~200]{silverman}. A positive common multiple
of the quotient orders gives $\kappa$ with $[\kappa]P_0$ in every
chosen neighborhood. Take
\[
 H=[2\kappa]P_0,\quad U=[2\kappa n]P_0,\quad
 V_0=[\kappa]P_0,\quad V_1=[\kappa n]P_0.
\]
The formal parameter and the derivative of $[n]$ make each $C_j$
integral above $S'$ after sufficiently shrinking the neighborhoods.
Infinite order preserves the nonzero guards.

Choose rational $q$ integral at $S'$ and arbitrarily close to
$1/\sqrt{R_0}$ in $\mathbf R$, without equality. Denominators
$R_0m+1$ with nearest integer numerators give such choices, adjusting
by one if equality occurs. At $S'$, $a_0$ is a unit and
$W\equiv1$, $\Theta\equiv2$ modulo $p^{v_p(R_0)}$.
Integral trace justifies this also at the ramified place two.
The pole input has valuation at least four at two and at least six
at the other primes of $S_0$. Outside $S'$, its valuation is
$2v_p(a_0)-2v_p(\Theta)$ and is even, so its reduced denominator is
a square. In both real embeddings, $R$ tends to zero positively,
$W\to1$ and $\Theta\to2$. All guards hold and
$0<B_{\rm pole}<1/2$ for a sufficiently close choice.
Theorem~\ref{hyp:PC} supplies the pole witnesses.
\end{proof}

\subsection{Exact ordinary value}
Define $\Val(N,t)$ by membership of $N$ and the disjunction
\[
 \bigvee_{c\in\Cset}(N=\iota(c)\wedge t=c)
 \quad\vee\quad
 \exists\mathcal K\,
 (\Close(N,\mathcal K,t)\wedge\Test(N,\mathcal K,t)).
\]
The same five coordinates $\mathcal K$ are used in both calls.
In an accepting main branch $n\notin\Cset$, $\kappa\ne0$, and
Lemma~\ref{lem:test} and the height estimate imply
\[
 h(t)\le H_6(S_{\rm ht})|\kappa F(n)|^{24}=B(n,\kappa).
\]
For $t=a/b\ne n$ in reduced form, $|a-nb|\ge1$, hence
\[
 |t-n|\ge1/b\ge e^{-h(t)}\ge e^{-B(n,\kappa)},
\]
contradicting~\eqref{eq:close-sound}. Thus $t=n$.
Conversely choose $\kappa$ by Lemma~\ref{lem:test-complete}, then use
\eqref{eq:close-complete} for that same shift. Private witness blocks
are disjoint, so the two calls can be satisfied simultaneously.
The six remaining integers have exactly their displayed branches.

Therefore $\Val(\iota(n),t)$ is exactly the graph $t=n$.
Existentially quantifying $N$ gives ordinary $\Z$ in $\Q$.
S01 and S16 multiply the exception and main-branch polynomials and
conjoin membership by a sum of squares. The literal substitutions in
Section~\ref{sec:syntax} and Appendix~\ref{app:literal} identify this
polynomial with the saved root $P_*$. This proves the zero-set part of
Lemma~\ref{lem:assembly}.

% End 04-computation-value.tex


% Begin 05-bounds-appendices.tex
\section{The cutoff and polynomial size}\label{sec:size}
The literal constants in S03 define the catalogue product $H_{\rm cat}$,
its maximum $H_{\max}$, $N_s=P_sH_{\rm cat}$ with
$P_s=\prod_{p\le324}p$, the least formal multiple $M_0$, and $c_*=4M_0$.
They take
\[
 R=\max(D,H_{\max},c_*,10^{12}),\quad
 A_{\rm ht}=2\cdot10^{13}\max(1,C_{\rm cm},Q_R^{17}),\quad
 T=(48A_{\rm ht})^{13}.
\]
The pole clearing constant is $\ell_0=12P_s^{2916}H_{\rm cat}^{732}$.
The symbol $F_0=324!$ is distinct from the six-contact polynomial $F(u)$.

\begin{proposition}[Literal cutoff]\label{prop:cutoff}
The fixed prescription satisfies
\[
 R=D,\qquad T=(96\cdot10^{13})^{13}D^{2210G_0^2D^3}
            <2^{\,2^{\,2^{9013}}}.
\]
This is a statement about the prescribed integer, independently of its
semantic role in excluding false multiplication triples.
\end{proposition}
\begin{proof}
We give the bounds of S25. Exact small calculations give
$10^5<W<2^{512}$ and $2^{2000}<F_0<2^{2250}$.
For $q\ge10^5$, the unconditional progression estimate
\cite[Theorem 1.2, (1.10)--(1.11), Corollary 6.17 with $A=1$]{bmor}
applies for $x\ge\exp(0.03\sqrt q(\log q)^3)$, including the possible
exceptional-zero case. At $X=2^{W^3}$ its error is at most
$X/(160\log X)$ and the hypothesis follows from
$0.03\sqrt W(\log W)^3<(\log2)W^3$.
Hence $\vartheta(X;W,a)>X/(2W)$, giving more than
$X/(2W^4)>22$ primes in each unit class. No prime at most 324
belongs to such a class. Thus
\[
 H_{\max}<X,\quad H_{\rm cat}\le X^{22W},\quad
 \log_2N_s<2^{2054}.
\]
For every odd prime dividing $N_s$, take the actual reduction order of
$P_0$ at each place above it, and let their product be $m_0$.
All preceding positive multiples stay in the dyadic formal subgroup,
since $P_0$ starts there and the subgroup is closed under addition.
A nonidentity point there has negative abscissa valuation, so its
reciprocal is defined; the constructor accepts the identity directly.
Thus every earlier search step is defined without a nontorsion premise.
At the bad place $s\equiv3\pmod7$, the point reduces to $(2,6)$
on $y^2=x^3$ and is nonsingular because $2y\not\equiv0$.
The other place $s\equiv4\pmod7$ has good reduction.
Take the actual order of the reduced point in the nonsingular
reduction group, using
\cite[VII.2.1, p.~188]{silverman}.
At a good place the residue degree is at most two, so Hasse's
bound \cite[V.1.1, p.~138]{silverman} gives order at most
$(p+1)^2<p^3$ for odd $p$. The singular smooth group over
$\F_7$ has order seven. There are at most two places over each
rational prime, so $m_0\le N_s^6$.
If $[N_s^4m_0]P_0$ is the identity, the search has already succeeded.
Otherwise its reciprocal coordinate is defined and the following
estimates apply. At each odd place, $[m_0]P_0$ is formal, and multiplication by
$N_s^4$ gives
\[
 v(1/x([N_s^4m_0]P_0))\ge2+8v_p(N_s).
\]
At two normalize $v_s(s)=1$, so $v_s(2)=2$. Initially
$v_s(x)=-6$, $v_s(y)=-9$, and $v_s(z)=3>2$.
The log isometry therefore gives
$v_s(z([2]Q))=v_s(z(Q))+2$ on this subgroup, and hence
\[
 v_s(1/x([N_s^4m_0]P_0))\ge6+16v_2(N_s).
\]
Writing this reciprocal as $a+bs$, extraction of $a$ and $b$
uses division by $2$ and $2s$ and loses at most three units of
$v_s$. Thus both rational coefficients lie in
$N_s^4\Z_{(N_s)}$, where denominators in $\Z_{(N_s)}$ are
prime to $N_s$. This proves the required divisibility and guard. Therefore the literal search
succeeds by $N_s^4m_0$ and $M_0\le N_s^{10}$.
This uses actual finite orders, not numerical upper bounds incorrectly
treated as multiples of those orders.

Consequently $c_*<2^{2^{2059}}$, whereas
$2^{2^{8008}}<D<2^{2^{9010}}$, so $R=D$.
Next $C_{\rm cm}<2^{50000D^2}$ and
\[
 \log_2Q_D^{17}=170G_0^2D^3\log_2D>50000D^2.
\]
The maximum is therefore $Q_D^{17}$, giving the identity for $T$.
Finally $G_0<2^{37}$, $E_0<D^3$, $\log_2Q_D<D^5$,
$\log_2A_{\rm ht}<D^7$, and $\log_2T<D^8<2^{2^{9013}}$.
\end{proof}

\subsection{Quantitative syntax bounds}
Use coefficient $\ell^1$ norm for integer polynomials and the weighted
pair norm $\|a+bs\|=\|a\|_1+2\|b\|_1$. They are submultiplicative.
The division-polynomial recurrences give
\[
 \deg f_n\le(n+1)^3,\qquad\|f_n\|\le2^{100(n+1)^3}.
\]
For the even recurrence the exponent slack is
$4m^3-24m-27\ge9$ for $m\ge3$; for the odd recurrence it is at least
ten for $m\ge2$. This absorbs the fixed factors and addition.

The following facts have been checked for the actual saved expression:
the absolute values of all literal fixed constants and all occurring
division-polynomial indices are at most $T$; the specified syntax-height
traversal has root height at most 256; each model recipe has at most
$2^{20}$ constraints and each composition recipe has at most $2^{20}$
terms. The check accounts for all 23 templates, seven recipes, private
witness blocks, 173 scalar constant nodes and 50 finite constructors. Its exact
artifact identity and independent replay are recorded in the source map.
Together with the constructor bounds below, this discharges the
quantitative part of $\mathsf B$. The all-input arithmetic interpretation of the same constituents is the
separate substitution induction of Section~\ref{sec:compiler}.

Put $B_*=T+1$. Coordinates of $\iota(n)$ for
$|n|\le T$ have absolute numerator and positive denominator at most $2^{K_*}$, with
$K_*=2^{30}(T+1)^3$. Evaluate the division polynomials at $P_0$ and use
the fixed coordinate formulas; at most eight quadratic-field operations
are needed, each taking an available height bound $H$ to at most $3H^4$.
There are at most $8(T+1)^3$ guard factors. The resulting bounds for
each finite constructor are
\[
 \deg\le B_*^4,\qquad\log_2\max(1,\|\cdot\|_1)\le B_*^8.
\]

Give variables syntax height zero, constants height one, and add one
for each arithmetic operation or finite constructor. A template call
adds one plus the callee's intrinsic height and the maximum argument
height. The substitution inequalities
\[
 \deg P(Q)\le\deg P\max(1,\max_i\deg Q_i),\qquad
 \|P(Q)\|_1\le\|P\|_1\max(1,\max_i\|Q_i\|_1)^{\deg P}
\]
give degree and logarithmic norm at most $B_*^{8h}$ at positive height
$h$, with the bounded sums in the actual recipes treated as in S26.
This general call rule applies to polynomial arguments. In each outer
composition recipe the arguments are coordinate variables, so the call
is only a renaming, even when coordinates repeat. For any such map
$\sigma$,
\[
 \deg P(Y_{\sigma(i)})\le\deg P,\qquad
 \|P(Y_{\sigma(i)})\|_1
 =\sum_\beta\left|\sum_{\alpha\mapsto\beta}c_\alpha\right|
 \le\sum_\alpha|c_\alpha|=\|P\|_1.
\]
Thus these calls retain the callee's majorant height before charging
the sum-of-squares or product connective. The seven recipe heights are
$228,233,235,251,252,254,256$. A rule charging an extra unit for each
of these mere renamings would describe a looser syntax-depth bound,
not this proved degree-and-norm majorant.
For the model wrappers use the safe height bound four for the
generator equations and root differences; the callee contribution
dominates this bound. The literal height-two claim for those equations
in S26 is not used.
This height is a majorant for the compact constructor syntax, not for
an expanded monomial tree. The height-256 bound yields
\[
 \deg P_*\le(T+1)^{2048},\qquad
 |[X^\alpha]P_*|\le2^{(T+1)^{2048}}.
\]

Since $T$ is an integer, $T+1\le U=2^{2^{2^{9013}}}$. Hence
\[
 (T+1)^{2048}\le U^{2048}=2^{2^{2^{9013}+11}}
                     <2^{2^{2^{9014}}}.
\]
As $9014<65536=\tw4$, the degree is less than $\tw7$ and the
coefficient magnitude less than $\tw8$. The root recipe declares
$292\,890\,231\,904$ witnesses. Together with the preceding estimates,
these complete the bounds for the same $P_*$ in
Theorem~\ref{thm:target}.

\section{Consequences of the integer definition}
The proof of Theorem~\ref{thm:target} combines the following results.
\begin{center}\small
\begin{tabular}{p{.20\textwidth}p{.43\textwidth}p{.27\textwidth}}
\toprule
Result & Mathematical inputs & Proof location\\\midrule
Index group & Two-isogeny descent, reductions, finite height search &
Section~\ref{sec:model}\\
Six-contact bound & Quaternionic model, tame pullback, compatible-isogeny
height theorem, theta comparison & Section~\ref{sec:height}\\
Pole completeness & Simultaneous descent, restricted converse, forward
rank-one theorem & Section~\ref{sec:pole}\\
Multiplication & Index group, height bound, selected primes,
pole completeness, finite guard & Section~\ref{sec:mult}\\
Precision relation & Decision program, lowering, power and packed histories &
Section~\ref{sec:compiler}\\
Ordinary value & Shared-shift Close and Test, height and rational separation &
Section~\ref{sec:bridge}\\
Bounded $P_*$ & Literal substitutions and quantitative syntax &
Sections~\ref{sec:syntax}, \ref{sec:size}; Appendix~\ref{app:literal}\\\bottomrule
\end{tabular}
\end{center}

\subsection{Three corollaries}\label{sec:four-corollaries}
We apply the exact equality between the rational zero projection of
$P_*$ and ordinary $\Z$ proved in Theorem~\ref{thm:target}.

\begin{corollary}[Negation of Mazur's finite-component conjecture]
\label{cor:mazur}
Mazur's
finite-connected-components conjecture for the real-topological closure
of rational points, Conjecture~3 of \cite{mazur1992,mazur1995}, is
false. Its stronger Conjecture~2 is false as well.
\end{corollary}
\begin{proof}
We recall Mazur's projection argument, originating in
\cite{mazur1992} and repeated in \cite[p.~166]{mazur1995}.
Conjecture~3 in \cite[p.~165]{mazur1995} asserts that the closure
of the rational points of any variety over $\Q$ in its real locus has
finitely many connected components. In that source a variety means a
reduced scheme of finite type; irreducibility is not required.
Put $r=292\,890\,231\,904$ and
\[
 V=\bigl(\{P_*(t,w_1,\ldots,w_r)=0\}\subset\mathbf A^{r+1}_{\Q}\bigr)_{\rm red},
 \qquad C=\overline{V(\Q)}\subset V(\mathbf R).
\]
Reduction changes neither the rational nor the real points. Let
$\pi:V(\mathbf R)\to\mathbf R$ be projection to $t$.
Theorem~\ref{thm:target} gives $\pi(V(\Q))=\Z$. Since $\Z$ is closed in $\mathbf R$
and $\pi$ is continuous, $C\subset\pi^{-1}(\Z)$, so $\pi(C)\subset\Z$.
Every integer already occurs on $V(\Q)$, hence $\pi(C)=\Z$.
The image of any connected component of $C$ is a connected subset of
$\Z$, therefore a singleton. Points projecting to distinct integers
belong to distinct components. Thus $C$ has infinitely many connected
components, contradicting the stated conjecture. Conjecture~2 in
the same source would give the homotopy type of a finite simplicial
complex and hence only finitely many connected components, so it
implies Conjecture~3 and is contradicted too.
\end{proof}

\begin{corollary}[A universal rational Diophantine polynomial]
\label{cor:rational-universal}
There exist $m\ge1$ and a fixed polynomial
$U\in\Z[e,n,x_1,\ldots,x_m]$ such that, for every $e,n\in\N$,
\[
 \exists\mathbf x\in\Q^m\ U(e,n,\mathbf x)=0
 \quad\Longleftrightarrow\quad
 \text{machine $e$ halts on input $n$}.
\]
Here $e$ and $n$ are ordinary integer parameters, with $\N$ embedded
in $\Q$, and are not elliptic-curve coordinates. In particular, every
computably enumerable subset of $\N$ is represented by a suitable
fixed value of $e$ in this one polynomial.
\end{corollary}
\begin{proof}
Fix an effective enumeration of machines. The Davis--Putnam--Robinson--Matiyasevich theorem, \cite[Basic Theorem, p.~3]{matiyasevich1971},
applied to the universal halting relation gives a fixed
$D\in\Z[e,n,z_1,\ldots,z_s]$, with $s\ge1$, such that, for $e,n\in\N$,
\[
 \text{machine $e$ halts on $n$}
 \quad\Longleftrightarrow\quad
 \exists\mathbf z\in\Z^s\ D(e,n,\mathbf z)=0.
\]
To use integer witnesses in this statement, replace each nonnegative
witness by a sum of four integer squares. In the positive-integer
convention of the cited source, apply the theorem to the predicate
that machine $a-1$ halts on $b-1$, substitute $a=e+1$, $b=n+1$,
and use one plus that sum for each witness; see also
\cite[p.~2]{matiyasevich1971}. A dummy witness ensures $s\ge1$ if needed.
For each $z_i$ take a separate block $\mathbf w_i$ of $r$ rational
variables, where $r$ is the witness count of $P_*$. Define
\[
 U(e,n,\mathbf z,\mathbf w_1,\ldots,\mathbf w_s)
   =D(e,n,\mathbf z)^2+\sum_{i=1}^s P_*(z_i,\mathbf w_i)^2.
\]
A rational zero makes every summand zero. Theorem~\ref{thm:target}
then forces every $z_i$ to be an integer and yields the displayed
halting condition. Conversely, integer witnesses for $D$ have separate
rational witness blocks for $P_*$, giving a zero of $U$. All blocks
are private and the polynomial is fixed independently of $e,n$.
No variable, degree or coefficient bound proved for $P_*$ is asserted
for this new polynomial $U$.
\end{proof}

OpenAI already obtained Turing degree $0'$ for this set in
\cite[Corollary 7.1, pp.~56--57]{openaihtp}; that paper explicitly
does not assert many-one completeness. The stronger direction below
is $0'\le_m\operatorname{HTP}(\Q)$, obtained from the integer
definition through the universal polynomial. The reverse many-one
reduction is the standard enumeration of rational solutions.

\begin{corollary}[The many-one degree of $\operatorname{HTP}(\Q)$]
\label{cor:htp-degree}
\[
 \operatorname{HTP}(\Q)\equiv_m 0',
 \qquad 0'=\{e\in\N:\text{machine $e$ halts on input $e$}\}.
\]
Here $\operatorname{HTP}(\Q)$ is the set of codes of integer-coefficient
polynomials, of arbitrary finite arity, having a rational zero, under a
fixed effective coding.
\end{corollary}
\begin{proof}
Given $e$, specialize the fixed polynomial of
Corollary~\ref{cor:rational-universal} at its ordinary integer parameters
to obtain $U(e,e,\mathbf x)$. Computing its coefficient list and code is
effective, and
\[
 e\in0'\quad\Longleftrightarrow\quad
 \ulcorner U(e,e,\mathbf x)\urcorner\in\operatorname{HTP}(\Q).
\]
This proves $0'\le_m\operatorname{HTP}(\Q)$.

Conversely, from the code of a polynomial $f$ construct a machine that
ignores its input and enumerates all rational tuples of the appropriate
arity. It evaluates $f$ using exact rational arithmetic and halts as soon
as it finds a zero. Its index $j(f)$ is computable from the finite code
of $f$, and
\[
 f\in\operatorname{HTP}(\Q)\quad\Longleftrightarrow\quad j(f)\in0',
\]
where the left side uses the code of $f$. Thus
$\operatorname{HTP}(\Q)\le_m0'$. Invalid polynomial codes, if present in
the coding convention, are sent to a fixed nonhalting index. Using
rational-coefficient polynomial codes gives the same result, since
clearing their finitely many denominators is effective and preserves
rational zeros. The two reductions prove the assertion.
\end{proof}

More generally, an integer Diophantine problem translates to a rational
one by adjoining $P_*$ for every integer witness, with private witness
blocks. Natural witnesses can first be expressed as sums of four integer
squares. Bounds for the resulting polynomial require their own composition
calculation; the bounds for $P_*$ do not automatically transfer.

For example, a machine that enumerates ZFC proofs and halts upon finding
a contradiction has a rational Diophantine halting condition by
Corollary~\ref{cor:rational-universal}. This equivalence does not assert
that a contradiction exists, and no additional explicit polynomial is
specified here.

% Begin degree-eight construction
\section{The base residual compiler}
\label{sec:degree8-original}

This section records the base residual compiler used as a construction
stage for the quartic in Section~\ref{sec:quartics}. Its degree-eight
sum-of-squares encoding fixes the arithmetic interfaces and conventions;
it is not a third final polynomial specification. The construction uses
a joint comparison of four elliptic indices, a Lucas equation, and the
five-contact height estimate proved below. The archived counts describe
this base compiler before its reciprocal revision and quadratic lift.

\begin{proposition}[Finite specification of the base compiler]\label{prop:o8-finite}
The prescription below specifies a polynomial
\[
 F_8\in\Z[X,V_1,\ldots,V_{15671}]
\]
of degree exactly eight, with exactly $129497$ nonzero monomials and
maximum absolute coefficient
\[
 685783949153203534233601<2^{80}.
\]
These statements are finite identities about the displayed construction
and its integer coefficient table.
\end{proposition}

\begin{theorem}[Arithmetic interpretation]\label{thm:o8-conditional}
For every
$x\in\Q$,
\[
 x\in\Z\quad\Longleftrightarrow\quad
 \exists(v_1,\ldots,v_{15671})\in\Q^{15671},\qquad
 F_8(x,v_1,\ldots,v_{15671})=0.
\]
\end{theorem}

The required height estimate is proved below; Theorem~\ref{thm:height}
uses a different contact set. The seven-variable pole lemma below is proved from
Theorem~\ref{hyp:RC}; it is not inferred from Theorem~\ref{hyp:PC},
which has different variables and a restricted completeness domain.
The finite coefficient calculation does not establish the height bound.
The supplied unpublished Freitag archive is the provenance of this
construction and of the original coefficient table. It is not cited here
as an arithmetic theorem. The published inputs already used in
Sections~\ref{sec:model} and \ref{sec:mult} remain
available, together with the quaternion trace theorem specified below.

\subsection{The height input and the exact pole predicate}

Put
\[
 S_* =\{p:p\le76176\},\qquad
 \mathcal D=\{0,1,2401,-343/32\},
\]
\[
 H_*=2(4608^{165888}+6^{100})
          \left(\prod_{p\le76176}p\right)^{24}+128.
\]
For $a\in\Q\setminus\mathcal D$, let $M_5(a)$ be the product of
the primes outside $S_*$ at which the point $(a:1)$ has positive odd
contact order with one of $\infty,0,1,2401,-343/32$. Explicitly, choose
local homogeneous coordinates $(u:v)$ and subtract
$\min(v_p(u),v_p(v))$ from the valuations of $v$ and of
$u-\delta v$ for $\delta\in\mathcal D$.

\begin{theorem}[The five-contact height estimate]\label{hyp:o8-height}
For every $a\in\Q\setminus\mathcal D$,
\[
 h(a)\le H_*M_5(a)^{24}.
\]
The displayed constant satisfies $\log_2H_*<4783227$.
\end{theorem}
\begin{proof}
Write $s=a$, $M(s)=M_5(a)$, $H=H_*$,
$B_*=\prod_{p\le76176}p$, and
$\mathcal C=\{\infty,0,1,2401,-343/32\}$.
We prove the geometric comparison, the non-CM estimate through the
modular level, and the CM estimate over the same fiber field.
All heights use the original stable Faltings normalization.
\paragraph{Recovering the contact line from published quotients}
Let $B/\mathbb Q$ be the indefinite quaternion algebra of discriminant $210$,
$\mathcal O$ a maximal order, $X$ its maximal-level Shimura curve, and $W$
its Atkin--Lehner group. The standard genus formula gives
\[
 g(X)=1+\frac{\varphi(210)}{12}=5,
\]
because the elliptic-point terms vanish: $5$ splits in $\mathbb Q(i)$ and
$7$ splits in $\mathbb Q(\sqrt{-3})$, so neither quadratic field embeds in
$B$. The group $W$ is $(\mathbb Z/2\mathbb Z)^4$ with involutions $w_m$
indexed by square classes of divisors of $210$.

The published models of \cite[Theorem 4.2 and Table 2, pp.~938--939]{GonzalezRotger2006} have the following coefficient rows
$(a,b,c,d,e)$ for $y^2=ax^4+bx^3+cx^2+dx+e$:
\[
\begin{array}{c|rrrrr}
 m&a&b&c&d&e\\\hline
30&-43&-686&-2915&-1372&-172\\
42&-43&600&-986&-600&-43\\
70&-43&256&-130&-768&-387\\
105&-43&2&505&12&-1548
\end{array}
\]
The cited theorem identifies each with $X/\langle w_m\rangle$.
The genus formula and the Atkin--Lehner group law used here also appear
in the same published paper, Proposition 5.2 and the paragraph following
Proposition 5.3, pp.~940--941.

Here is a deduction of all the required branch information, rather than
an additional assumption about it. Each of these four genus-one quotients
has eight fixed points upstairs, by Riemann--Hurwitz. A stabilizer for the
$W$ action on a characteristic-zero smooth curve is cyclic, and therefore
has order at most two. Each nonempty fixed set of an involution is a union
of $W$-orbits of size eight. Riemann--Hurwitz for $X\to X/W$ now gives
\[
 8=16(2g(X/W)-2)+\sum_{1\ne w\in W}\#\operatorname{Fix}(w).
\]
The four known fixed sets already contribute $32$. Thus $g(X/W)=0$,
the total is $40$, and there is exactly one further orbit of eight fixed
points. Inertia generators around the branch points have product one.
The four classes $30,42,70,105$ are independent in $W$, and their product
is the class $210$. Hence there are exactly five branch points, each of
index two, with distinct labels $30,42,70,105,210$. Since the involutions
are defined over $\mathbb Q$ and each label occurs only once, each branch
point is rational. Three of them trivialize the genus-zero quotient.

Normalize $b_{210}=\infty$, $b_{105}=0$, $b_{70}=1$, and write
$b_{42}=\lambda$, $b_{30}=\mu$. On a genus-one quotient the residual
group $(\mathbb Z/2)^3$ has translation subgroup of order four, necessarily
the full two-torsion group after choosing an origin. Its quotient has the
same Jacobian $j$-invariant: it is multiplication by two over the algebraic
closure. The remaining double cover of $X/W$ branches at the four points
other than $b_m$. Therefore the four quartic invariants determine the four
Legendre invariants of these four-point subsets.

For a quartic put
\[
 I=12ae-3bd+c^2,\quad
 J_q=72ace+9bcd-27ad^2-27b^2e-2c^3,\quad
 j=\frac{6912I^3}{4I^3-J_q^2}.
\]
The exact computation gives the Legendre parameters
$1/2401,32/375,1/225,5/1029$, in row order. Enumerating the six cross-ratio
values of each of the first two parameters and testing the last two
invariants leaves exactly
\[
 \lambda=2401,\qquad\mu=-343/32.
\]
The script \texttt{check\_height.py} reproduces this rational computation.
Every nonzero determinant of the five primitive contact forms has prime
support in $\{2,3,5,7\}$.

\paragraph{Integral geometry and the degree 576 field}
Choose $\mu_0\in\mathcal O$ with $\mu_0^2=-210$. The embedding criterion
allows this choice in a maximal order; maximal orders are conjugate here.
Use the maximal moduli stack $\mathscr X$ and the entire principal-level
three scheme $Y$, including its geometric components. Yuan's integral
moduli theorem \cite[Theorem 1.1 and Sections 2.1--2.2]{Yuan2024} makes $\mathscr X$ regular, proper and semistable over
$\mathbb Z$, with a universal abelian surface, and smooth away from $210$.
Its generic stabilizer is precisely $\{\pm1\}$. Principal level three is
fine and gives a universal full three-torsion marking.

Work first over $R_0=\mathbb Z[1/210]$. The effective level-change group
has order $36$, prime to every residue characteristic of $R_0$. Its generic
action is free. For a finite tame group acting on a smooth relative curve,
the fixed locus of an element is smooth locally along its fixed points;
an isolated special fixed point consequently lifts. Thus no new fixed
point occurs here, and $Y_{R_0}\to X_{R_0}$ is finite etale. The $W$ action
extends to the smooth stable model of $X$. Its stabilizers remain cyclic,
of order at most two. Distinct involution fixed loci cannot intersect;
each fixed locus is etale over the base. Consequently $X_{R_0}/W$ is a
smooth genus-zero curve with five disjoint branch sections of index two.
The three rational sections already chosen identify it with
$\mathbb P^1_{R_0}$. Their values are the five elements
of $\mathcal C$.

Set $K=\widehat{\mathcal O}^{\times}$, $U=U(3)$,
$Z=\mathbb A_f^\times$, and $N=N_{B(\mathbb A_f)^\times}
(\widehat{\mathcal O})$. Normalization preserves $3\widehat{\mathcal O}$,
so $U$ is normal in $N$. The canonical adelic action is defined over the
reflex field $\mathbb Q$ \cite[2.2.5 and 2.3.1]{Deligne1979}; this refers to the entire canonical model, not
an individually selected complex component. Every central finite idele
is a rational scalar times a unit idele. At three its scalar residue is
$1$ or $-1$, which a rational sign removes. Thus $Z\subset\mathbb Q^\times U$
acts trivially on $Y$.

At three the quaternion algebra is ramified. Its maximal order modulo
three has $81$ elements, radical of size $9$, and hence $72$ units.
It follows that
\[
 1\longrightarrow (\mathcal O/3\mathcal O)^\times/\{\pm1\}
 \longrightarrow\Gamma=N/(ZU)\longrightarrow W\longrightarrow1,
 \qquad |\Gamma|=36\cdot16=576.
\]
Above an unbranched point the actions of $W$ and of the effective
level-change kernel are free. The fiber therefore is a torsor of size
$576$ for this constant group. The Galois permutations commute with its
regular action, so belong to its regular centralizer, also of order
$576$. The normal splitting field $L$ consequently satisfies
$[L:\mathbb Q]\mid576$: the Galois group is a subgroup of that
regular centralizer. In particular $[L:\mathbb Q]\le576$.

All fiber surfaces and their three-torsion markings are defined over $L$.
For two such surfaces, a homomorphism annihilating the three-torsion
factors through $[3]$. Hence reduction embeds
\[
 \operatorname{Hom}(A,B)/3\operatorname{Hom}(A,B)
 \hookrightarrow\operatorname{Hom}(A[3],B[3]).
\]
The finite Galois image on the integral Hom lattice is trivial modulo
three. The principal congruence subgroup modulo three in an integral
general linear group is torsion-free, so this image is trivial. Thus
every geometric homomorphism between fiber surfaces is defined over $L$.

At a prime outside $210$, tame local coordinates for the quotient are
$w=z^2$ up to a unit at a branch section, and etale elsewhere. A base
point with even contact gives no ramification after an unramified
extension. The universal abelian scheme extends at the resulting
integral point of $Y_{R_0}$. Therefore
\[
 \operatorname{rad}(D_LN_{\rm bad})\mid B_*M(s).
\tag{R}\label{eq:h5-R}
\]
This assertion includes contact with infinity. It is not a statement
obtained merely from the numerical prime support of pairwise determinants.

The normalizer correspondences give isogenies between the conjugate
surfaces. In the non-CM case the geometric endomorphism algebra is $B$;
Skolem--Noether adjusts a conjugate isogeny to commute with its chosen
quadratic subfield. All such maps lie over $L$ by the Hom argument.
These facts establish the usual GL$_2$-type and compatibility hypotheses;
they do \emph{not} establish the separate bounded descent to a
GL$_2$-type variety over $\mathbb Q$.

\paragraph{The metric comparison, with its exact additive budget}
Put $t=32s$, $(a_1,a_2,a_3,a_4)=(0,32,76832,-343)$ and
$P(t)=\prod_i(t-a_i)$. The sections
\[
 q_0=(dt)^2/P(t),\qquad q_1=tq_0
\]
are holomorphic bicanonical sections on $X$, with no common zero. At a
finite branch point use $t-a_i=z^2$; at infinity use $t=z^{-2}$.
At infinity $q_0$ vanishes to order two and $q_1$ does not vanish.

For each prime let $\delta_p=\max_{i<j}v_p(a_i-a_j)$, so
\[
 J=\prod_p p^{\delta_p}=2^{10}3^2 5^3 7^4=2765952000.
\]
Both $J^2q_i$ extend across the maximal regular stack. Indeed, at the
generic point of a vertical component semistability makes $p$ a
uniformizer, and the relative differential satisfies $v(du)\ge v(u)$.
Writing $r_i=v(t-a_i)$ gives $v(dt)\ge\max r_i$. If $v(t)<0$, all $r_i$
equal $v(t)$ and both sections have nonnegative valuation. Otherwise all
$r_i\ge0$ and at most one exceeds $\delta_p$. In either case
$2\max r_i-\sum r_i\ge-2\delta_p$. This proves the required vertical
integrality; regularity handles codimension two.

For the archimedean estimate write $r_i=|t-a_i|$, $U_0=1+\sum_i r_i$ and
\[
 ds_\rho^2=\frac{U_0}{\prod_i r_i}|dt|^2,\qquad
 C_{\rm met}=16(12\cdot76832+9)^3=12540153541944538512.
\]
Its pullback is smooth and positive at every branch point and infinity.
The curvature formula, $\Delta U_0=\sum_i r_i^{-1}$ and
$|\nabla U_0|^2\le16$ give
\[
 -K_\rho\ge\frac{\sum_i\prod_{j\ne i}r_j}{2U_0^3}.
\]
When $|t|\le2\cdot76832+2$, at most one $r_i<1/2$, and
$U_0\le12\cdot76832+9$. The displayed bound is at least
$C_{\rm met}^{-1}$. Outside this disk, $|t|/2\le r_i\le3|t|/2$
and $U_0\le7|t|$, giving the stronger bound $1/1372$.
The maximum principle against the curvature $-1$ metric implies
$ds_\rho^2/C_{\rm met}\le ds_{\rm hyp}^2$. Yuan's Petersson norm is
$\|d\tau\|=2\operatorname{Im}\tau$. Consequently
\[
 \|q_0\|_{\rm Pet}\le4C_{\rm met}/U_0,\qquad
 \|q_1\|_{\rm Pet}\le4C_{\rm met}|t|/U_0\le4C_{\rm met}.
\]

Yuan's isometry \cite[Theorem 1.1 and metric definitions, pp.~2--3]{Yuan2024} identifies the square of the Hodge bundle with
$210\omega_{\mathscr X}^{\otimes2}$ with the induced metric. The Hodge
height in his normalization equals the original stable Faltings height
plus $\log\pi$ in dimension two. Thus
\[
 h_{\bar\omega^{\otimes2}}(y)=2h_F(A_y)+2\log\pi+\log210.
\]
The sign of $\log210$ is positive: the sublattice $210\omega^{\otimes2}$
has height $h_{\bar\omega^{\otimes2}}-\log210$. Evaluate on an extension
where the point extends to the proper maximal stack; its universal
abelian scheme computes stable height. The integral sections and the
supremum estimate yield
\[
 h(s)\le2h_F(A_y)+2\log\pi+\log210+2\log J
                       +\log(4C_{\rm met})+\log32.
\]
The integer comparison
$210J^2(4C_{\rm met})32\cdot16<2^{160}$, together with $\pi<4$ and
$\log2<3/4$, bounds the additive term by $120<128$. Therefore
\[
 h(s)\le2\max(0,h_F(A_y))+128.
\tag{G}\label{eq:h5-G}
\]


% Proposed repair of the original H5 arithmetic step. Independent review pending.
% Uses the existing baseline bibliography keys, plus GaudronRemond2023 and
% the published 2021 edition of Voight, chapter 43.
\paragraph{A level bound for the non-CM fibers}
\label{sec:h5-level-repair}

This argument allows the scalar character order and coefficient field to
vary. It does not assert bounded same-field descent to a rational
GL$_2$-type variety. Let $s$ be a rational point outside the five contacts,
let $L/\mathbb Q$ be the normal fiber field constructed above, and let
$A/L$ be a non-CM fiber surface. Thus $[L:\mathbb Q]\mid576$, all
geometric homomorphisms between conjugate fibers are defined over $L$,
and $\operatorname{End}^0(A_{\overline{\mathbb Q}})=B$, the quaternion
division algebra of discriminant $210$. Put
\[
 R=B_*M(s),\qquad T=\{p:p\mid R\}.
\]
The geometric argument gives that $L$ is unramified outside $T$ and
$A$ has good reduction outside the primes above $210$. All conjugate
QM structures lie in the same normalizer orbit.

\paragraph{Good reduction after full torsion.}
First a QM surface with the indicated division action has potentially
good reduction at every finite place. Indeed, pass to an extension with
semistable reduction, using \cite[Theorem 3.4, p.~406]{silverbergzarhin}.
Every endomorphism extends uniquely to the N\'eron model by its mapping
property. On the identity component of the geometric special fiber it
preserves the maximal torus $\mathcal T$: this torus is the maximal
connected affine subgroup, and every homomorphism from a torus to an
abelian variety is zero. If its rank $t$ is positive, restriction gives
a unital homomorphism
\[
 B\longrightarrow\operatorname{End}^0(\mathcal T)=M_t(\mathbb Q),
 \qquad 1\le t\le2.
\]
The kernel is a two-sided ideal of the simple algebra $B$, and the map
is unital, so it is injective. For $t=1$ dimensions prohibit this;
for $t=2$ it would be an isomorphism from a division algebra to
$M_2(\mathbb Q)$, also impossible. Hence $t=0$. The semistable identity
component is abelian, which is good reduction. These N\'eron-model and
reduction conventions are those of \cite[\S3, pp.~405--406]{silverbergzarhin}.

Set $E=L(A[11])$. We next justify that $E/\mathbb Q$ is Galois.
For each $\sigma$, the normalizer relation supplies a compatible
isogeny ${}^{\sigma}A\to A$ of degree prime to $11$. Here is the
prime-to-$11$ detail. In the analytic description
$A_z=\mathbb C^2/(\mathcal O(z,1)^t)$ of the published 2021 edition
of Voight \cite[equations (43.1.2)--(43.1.3)]{voight2021}, a positive normalizer element
$b=\left(\begin{smallmatrix}a&b'\\c&d\end{smallmatrix}\right)$
gives the scalar compatible map $(cz+d)I:A_{bz}\to A_z$.
When $\mathcal O b\subset\mathcal O$, its degree is
$[\mathcal O:\mathcal O b]=\operatorname{Nrd}(b)^2$.
At $11$ the algebra is split and the maximal-order normalizer is
$\mathbb Q_{11}^*\operatorname{GL}_2(\mathbb Z_{11})$.
Rescale $b$ rationally to make it an $11$-adic unit, then clear the
remaining denominators by an integer prime to $11$. This produces
the claimed isogeny. Changes of level marking use isomorphisms.
All these geometric homomorphisms are $L$-defined by construction.
Consequently the $11$-torsion representations of $A$ and its conjugates
over $L$ have the same kernel. Thus every conjugate of $E$ equals $E$.

The faithful action on $A[11]$ now gives
\[
 [E:\mathbb Q]\le576|\operatorname{GL}_4(\mathbb F_{11})|
 <2^{10}11^{16}<2^{66}.
\tag{LR1}
\]
For residue characteristic different from $11$, full rational
$11$-torsion gives semistable reduction over $E$ by
\cite[Theorem 3.5, p.~406]{silverbergzarhin}. The preceding torus
argument makes this good reduction. In particular it applies at
each of $2,3,5,7$.

\paragraph{The scalar lift with controlled support.}
Choose $B$-compatible isogenies $\mu_\sigma:{}^\sigma A\to A$ over $L$,
with $\mu_1=1$, and inflate this finite choice to $G_{\mathbb Q}$.
Their discrepancies $c(\sigma,\tau)$ commute with $B$, so lie in
$\mathbb Q^*$. Put $a_\sigma=(\deg\mu_\sigma)^{-1/4}>0$ and
$\nu_\sigma=a_\sigma\mu_\sigma$, with trivial Galois action on the
added algebraic coefficients. Since a scalar $c$ has degree $c^4$,
the discrepancies of the $\nu_\sigma$ form a sign cocycle
$\varepsilon\in Z^2(G_{\mathbb Q},\{\pm1\})$, inflated from
$\operatorname{Gal}(L/\mathbb Q)$.

The sign class is unramified, hence zero, locally outside $T$.
At each odd prime where its Brauer invariant is $1/2$, choose an odd
Dirichlet character modulo $p$; at $2$ use the odd character modulo
$4$ if required. Their product $\chi$ is tame at every odd prime,
has local conductor at most $4$ at $2$, and has order
\[
 m\le2\prod_{\substack{p\in T\\p\ne2}}(p-1)\le2R.
\]
For constant coefficient groups, the boundary in
$1\to\{\pm1\}\to\mu_\infty\xrightarrow{2}\mu_\infty\to1$
is the obstruction to a character square root. Local reciprocity
identifies this obstruction with the value at $-1$: the cyclic
torsion factor has precisely this obstruction, while the valuation
and principal-unit factors admit character square roots.
Thus the boundary of $\chi$ equals $[\varepsilon]$ at every finite
place. Brauer reciprocity and the global-to-local injection give
equality globally. A pointwise square root, corrected by a sign
cochain, gives
\[
 b:G_{\mathbb Q}\longrightarrow\mu_{2m},\qquad
 b^2=\chi,\qquad\delta b=\varepsilon.
\tag{LR2}
\]
Only this \emph{sign} cochain is asserted to take root-of-unity values;
the full scalar splitting also involves the $a_\sigma$.

On $G_L$, $b$ is a conjugation-invariant character. Outside $T$ its
inertia characters are quadratic. Only finitely many are nontrivial;
twist by the quadratic character of the product of those odd primes
to kill them. At each odd prime there is a unique nontrivial quadratic
inertia character. This changes no square or coboundary in (LR2),
and introduces possible ramification only at $2\in T$.
The kernel field $L'$ of $b|G_L$ is normal over $\mathbb Q$ and
\[
 [L':\mathbb Q]\le2m[L:\mathbb Q]\le2304R.
\tag{LR3}
\]

\paragraph{The compatible modular representations.}
Choose a number field containing all normalization scalars, the values
of $b$, and a splitting field of $B$. For each coefficient prime $\ell$
and an embedding into $\overline{\mathbb Q}_\ell$, set
\[
 \rho_\ell(\sigma)=b(\sigma)^{-1}a_\sigma\mu_\sigma\circ\sigma.
\]
The cocycle identity makes this a representation commuting with $B$;
after splitting $B$, matrix units give $\rho_\ell=r_\ell\oplus r_\ell$
with $\dim r_\ell=2$. On $G_{L'}$, $\rho_\ell$ is the Tate action of
$A$. Faltings' semisimplicity and endomorphism theorems and $\operatorname{End}^0(A)=B$
show that the multiplicity representation is absolutely irreducible.
Proper-smooth comparison and finite descent, as in
Lemma~\ref{lem:de-rham-descent}, make $r_\ell$ de Rham with Tate
Hodge--Tate weights $\{0,-1\}$; Berger's monodromy theorem makes it
potentially semistable.

It is odd. On Betti homology complex conjugation followed by the
holomorphic isogeny $\mu_c$ exchanges the two Hodge summands, each
a simple $B\otimes\mathbb C$-module. Its action on the multiplicity
space is therefore nonscalar. After the scalar correction its square
is one, so its eigenvalues are $1,-1$. The representation is unramified
outside $T\cup\{\ell\}$, by (LR2) and the reduction and field properties.
Pan's regular two-dimensional modularity theorem applies at $\ell=5$.
We use its published formulation in \cite[Section 7, p.~624]{calegari-icm2022},
with the definition of a geometric representation on pp.~611--612,
and the detailed author proof in \cite[Theorem 1.0.4]{pan2022}.
It requires no additional residual-image restriction at this prime and
identifies $r_5$ with a cuspidal modular representation up to a character
twist. The published formulation is used as stated; no byte or text
identity between Pan's author version and the journal version is asserted.

The twist can be taken finite and the weight is two, without invoking
a separate local-algebraicity theorem for characters. The twisting
character is a de Rham subrepresentation of the relevant Hom
representation. Comparing its integer weight with $\{0,-1\}$ and
the modular weights $\{0,1-k\}$ forces $k=2$ and twist weight zero.
On $G_{L'}$, the identity
$(\det r_5)^2=\det V_5(A)=\chi_{\rm cyc}^2$ shows that
$\det r_5/\chi_{\rm cyc}$ has finite image, since $G_{L'}$ has finite
index in $G_{\mathbb Q}$. The modular determinant
divided by $\chi_{\rm cyc}$ also has finite image. The square of the
twisting character consequently has finite image, so the character
itself has finite image. Absorb it in the modular form and take the
primitive weight-two newform $f$ of level $N$.

The conductor at $5$ is not omitted. For each prime
$p\notin T\cup\{3,5\}$, reduction of $\mu_\sigma$ composed with
relative $p$-Frobenius is a quasi-endomorphism of the good reduction
of $A$ at a chosen place. After clearing denominators its characteristic
polynomial is independent of the Tate prime. Half its trace, multiplied
by $b(\sigma)^{-1}a_\sigma$, is therefore one algebraic number giving
the trace of every $r_\ell$. Enlarge the coefficient field to contain
the coefficients of $f$. The equality with the trace of $f$ at five
is an equality of algebraic numbers; transporting it to an embedding
at three and applying Chebotarev and semisimplicity identifies $r_3$
with the three-adic realization of $f$. Thus use $r_3$ to compute the
conductor at $5$ and $r_5$ at every other prime. The modular level is
the conductor of these realizations by
\cite[\S\S0.3--0.7, Theorem A, pp.~409--411]{carayol}
and newform theory \cite[Theorems 1 and 4]{casselman}.
Here the totally real field is $\mathbb Q$, of odd degree one, so
Carayol's extra finite-place hypothesis for even degree does not apply.
His unramified normalization twist and contragredient preserve the
conductor; his theorem allows arbitrary nebentypus.
Li's Fricke relation and Shimura's
Jacobian construction, as in the baseline modularity argument, also
place $r_5$ in the Tate module of $J_1(N)$.

\paragraph{The primitive level equals the local conductor.}
Let $f\in S_2(\Gamma_0(N),\chi)$ be the primitive newform obtained
after absorbing the finite character twist, and let $\Pi_f$ be its
unitary cuspidal automorphic representation. The character $\chi$
is arbitrary; its conductor need only divide $N$.
For every prime $p$, write $a(\Pi_{f,p})$ for the local conductor
exponent. Then
\[
 a(\Pi_{f,p})=v_p(N).
\tag{CL1}
\]
Here is the exact published newform input. In
\cite[Theorem~1, p.~302]{casselman}, the local conductor ideal is
the largest ideal $\mathfrak c$ for which the appropriate central
character eigenspace under $\Gamma_0(\mathfrak c)$ is nonzero;
at that ideal the eigenspace is one-dimensional. The remark on
pp.~306--307 identifies this ideal with the exponent in the local
epsilon factor. Theorem~4 on p.~313 gives the global level spaces:
they vanish unless $\mathfrak c(\Pi)\mid M$, and at
$M=\mathfrak c(\Pi)\prod_p p^{b_p}$ their dimension is
$\prod_p(b_p+1)$. The paragraph immediately following that theorem
identifies all primitive newforms with the minimal level
$M=\mathfrak c(\Pi)$. This proves (CL1).

This statement includes nebentypus. Casselman's classical setup on
p.~309 allows every character $\epsilon$ modulo $M$ and uses
$f|_k\gamma=\epsilon(a)^{-1}f$ for
$\gamma=\left(\begin{smallmatrix}a&b\\c&d\end{smallmatrix}\right)
\in\Gamma_0(M)$. Since $ad\equiv1\pmod M$, this is exactly
$\epsilon(d)f$. Thus the usual nebentypus convention is already
covered. No squarefree-level hypothesis occurs.

For a coefficient characteristic $\ell\ne p$, Carayol's published
Theorem~A gives the full local Weil--Deligne compatibility, not only
an equality at unramified primes
\cite[\S\S0.3--0.7, pp.~409--411]{carayol}.
Its base field here is $\mathbb Q$, its real weight is two, and its
extra finite discrete-series-place condition applies only to even
degree totally real fields, so imposes no additional hypothesis
here. Theorem~A includes the dyadic cases; the weaker Theorem~B
is not the result being used.

Carayol normalizes reciprocity by sending geometric Frobenius to
a uniformizer. His Hecke correspondence in \S0.5 differs from
local Langlands by a contragredient and an unramified norm twist.
It is characterized there by equality of local $L$- and
epsilon-factors, including character twists. The conductor is
unchanged by either duality or an unramified twist. In particular,
the passage to the Tate-module convention used for $r_\ell$ changes
no conductor exponent. Equality of good Frobenius traces identifies
the global representation supplied by Carayol with the already
constructed irreducible modular representation, in this convention.
This uses Chebotarev and Brauer--Nesbitt; it does not use the
unpublished irreducibility observation in Carayol's remark after
Theorem~A. Since the matching representation is irreducible, its
matching semisimplification also forces the original representation
to be irreducible.

The epsilon-factor conductor in (CL1) therefore equals the Artin
conductor of $r_\ell|_{G_{\mathbb Q_p}}$:
\[
 a_p(r_\ell)=a(\Pi_{f,p})=v_p(N),\qquad p\ne\ell.
\tag{CL2}
\]
The Artin conductor here includes its tame term and Swan term;
equivalently one uses the full Weil--Deligne conductor, including
the monodromy contribution. One must not discard monodromy by
replacing the local representation with its semisimplification on
inertia. Carayol's compatibility supplies precisely the required
full local object. His epsilon-factor normalization and
Casselman's conductor-zero additive character give the same
exponent; changing an additive character only introduces the usual
common dimension-dependent shift, so no such change is needed.

Apply (CL2) with $\ell=5$ at every $p\ne5$ and with $\ell=3$ at
$p=5$. The companion proof identifies both $r_5$ and $r_3$ with
the realizations of this same $f$. Thus the asserted conductor
bounds control every prime divisor of the primitive level, even
when $p^2\mid N$. No compatibility theorem at $p=\ell$ is needed.

\paragraph{The conductor bound.}
Let $F=E\mathbb Q(\chi)$, where $\mathbb Q(\chi)$ means the kernel
field of the character, not its value field. On $G_F$, $b$ is a
conjugation-invariant quadratic character. Its kernel gives a Galois
extension $F'/\mathbb Q$ of degree at most two over $F$.
For $p=2,3,5,7$, let $e_E,e_{F'}$ denote local ramification indices
over $\mathbb Q_p$. From (LR1), $v_p(e_E)\le65$.
At odd $p$ both added extensions are tame, whereas at $2$ each
adds ramification index at most two. Consequently
\[
 v_p(e_{F'})\le65\ (p=3,5,7),\qquad v_2(e_{F'})\le67.
\tag{LR4}
\]
The possibly unbounded tame ramification index causes no problem.
For a finite Galois extension of $\mathbb Q_p$ with ramification
index $e$ and different exponent $\delta$, an Eisenstein polynomial
over the maximal unramified subfield gives
\[
 \delta\le e-1+ev_p(e).
\]
Indeed the derivative terms have distinct valuations modulo $e$,
and its leading derivative term has the stated valuation. In upper
numbering the different formula is
\[
 \delta/e=1-1/e+\int_0^\infty(1-|G^u|^{-1})\,du.
\]
Before the last positive break the integrand is at least $1/2$.
Thus (LR4) bounds this break for $F'/\mathbb Q_p$ by $136$.
These ramification and conductor conventions are those of
\cite[\S\S4.1--4.3, pp.~155--159]{serrelocal}.
Over $F'$ the scalar correction is trivial and $A$ has good reduction.
Above that break, therefore, the two-dimensional representation is
trivial on upper wild groups. Its Swan conductor is at most $272$,
and its tame conductor is at most two. Hence
\[
 a_p(r)\le274<300\qquad(p=2,3,5,7).
\]
For $p>7$, $L/\mathbb Q$ is tame because its degree divides $576$,
$A/L$ has good reduction, and $b$ is trivial on wild inertia:
on that pro-$p$ group its square is tame and its remaining image has
order at most two. Thus $a_p(r)\le2$ there; it is zero outside $T$.
Using the coefficient-prime comparison above proves
\[
 N\le210^{300}R^2.
\tag{LR5}
\]

\paragraph{The original constant is preserved.}
Apply the accompanying modular-Jacobian transfer calculation with
$C=210^{300}$ and (LR3)--(LR5). It gives
\[
 h_F(A)<2^{24}210^{3600}R^{24}
 <2^{28824}R^{24}<4608^{165888}R^{24}.
\tag{LR6}
\]
The last comparison uses $210<2^8$, $4608>2^{12}$, and
$28824<1990656=12\cdot165888$. Together with the separately proved
CM bound and geometric height comparison, (LR6) yields exactly
\[
 h(s)\le
 \bigl(2(4608^{165888}+6^{100})B_*^{24}+128\bigr)M(s)^{24}.
\]
No contact, coefficient constant, or radical exponent has been changed.
% UNAPPLIED CANDIDATE: global review and published Pan match are pending.

% Actual level-based repair: no uniform coefficient-field bound is used.
% The modular representation and level estimate are proved in the companion
% height-level-repair fragment. This file proves its height consequence.
\paragraph{Transfer from the level, allowing unbounded coefficient degree.}
Let $A/L$ be the non-CM quaternionic surface. Write $R\geq1$ for the
squarefree product of the prescribed rational bad primes, including $2$,
and assume $[L:\mathbb Q]\leq576$. The modularity construction provides
a weight-two newform of level $N$ and a field $L'$ over which the Tate
modules of $A$ and $J=J_1(N)$ share an irreducible constituent. In the
notation of the sign-splitting construction, the auxiliary Dirichlet
character has order
\[
 m\leq 2\prod_{\substack{p\mid R\\p\ne2}}(p-1)\leq2R.
\]
The sign-splitting cochain has values in $\mu_{2m}$. If one instead
uses the full scalar splitting map, it also includes the square root
of the degree map; that factor equals one on $G_L$. Thus in either
notation the restriction to $G_L$ is a character valued in
$\mu_{2m}$, and its kernel field satisfies
\[
 d'=[L':\mathbb Q]\leq2m[L:\mathbb Q]\leq2304R.
\tag{LT1}
\]
The quadratic twist used to remove extra ramification does not increase
this bound, since $2$ divides $2m$. There is no bound here on $m$
independent of $R$.

Faltings' Hom theorem, followed by scalar extension in the Hom space,
gives a nonzero $L'$-homomorphism $A\to J$. Absolute simplicity of $A$
and Poincar\'e reducibility therefore give an $L'$-isogeny
$A\times Z\to J$ for an abelian variety $Z/L'$. Put $g=\dim J$;
then $g\geq2$. The modular genus estimate
\cite[(5.9), p.~13]{vonKaenel2021} gives $g\leq N^2/24$.
The standard Belyi map has degree at most $N^2$ and
\cite[Theorem~1.1.1 and \S2.3]{javanpeykar} gives
\[
 h_F(J)\leq13000000N^{12}.
\tag{LT2}
\]
All heights in this argument use the original stable Faltings metric.

Suppose the companion local argument has proved $N\leq C R^2$,
where $C\geq1$. Then
\[
 g\leq C^2R^4/24,\qquad
 h_F(J)\leq13000000C^{12}R^{24}.
\]
Set $l=\log C$ and $r=\log R$. The published
Gaudron--R\'emond Theorem~1.9(1) \cite[pp.~7--8]{GaudronRemond2023}, combined with half-logarithmic isogeny-height variation,
gives
\[
 h_F(A)+h_F(Z)\leq h_F(J)+V,
 \quad
 V=2g^2\left(8g^2\log(7g)+\log d'
       +\log\max(1,\log d',h_F(J))\right).
\tag{LT3}
\]
Here the source variety for the isogeny estimate is $J_{L'}$.
The three logarithms in (LT3) satisfy respectively
\[
 \log(7g)\leq2l+4r,\quad
 \log d'<8+r,\quad
 \log\max(1,\log d',h_F(J))<17+12l+24r.
\]
For the last assertion, the bound in (LT2) also dominates $1$ and
$\log d'$; use $8+r\leq9R^{24}$.
It follows that
\[
 V\leq
 \frac{C^8R^{16}}{24^4}(32l+64r)
 +\frac{C^4R^8}{288}(25+12l+25r).
\tag{LT4}
\]
Using $l\leq C^4/4$, $r\leq R^8/8$, and also
$r\leq R^{16}/16$ in the second summand gives
\[
 \frac{V}{C^{12}R^{24}}
 \leq\frac{16}{24^4}+\frac{25+3+25/16}{288}<1.
\tag{LT5}
\]
The Bost lower bound in \cite[(2.2), p.~5]{vonKaenel2021}
gives $h_F(Z)\geq-\tfrac12\dim Z\log(2\pi^2)>-\tfrac32g$.
Consequently
\[
 h_F(A)<(13000000+1+1)C^{12}R^{24}
       <2^{24}C^{12}R^{24}.
\tag{LT6}
\]
This calculation needs no uniform bound on the degree of the newform's
coefficient field and no assertion that all its geometric factors are
isogenous to $A$.

For the companion conductor estimate one may take $C=210^{300}$.
Since $210<2^8$ and $4608>2^{12}$,
\[
 2^{24}C^{12}=2^{24}210^{3600}
 <2^{28824}<2^{1990656}<4608^{165888}.
\tag{LT7}
\]
Thus the non-CM height inequality is strictly within the original
coefficient budget, once the stated modularity and level estimate have
been established. In the five-contact application $R=B_*M(s)$, so
the radical exponent is exactly $24$, as required.

% Original 15671 architecture: independent CM branch, version 1.
% This lemma is conditional on the separately audited field/ramification lemma.
\paragraph{The CM height estimate over the same field}
Let $A/L$ be an abelian surface of CM type, let $B$ be a quaternion
division algebra over $\mathbb Q$ embedded in
$\operatorname{End}^0_{\overline L}(A)$, and suppose that all geometric
endomorphisms of $A$ are defined over $L$. Assume
\[
 d=[L:\mathbb Q]\le576,\qquad
 \operatorname{rad}(|D_L|)\le R,\qquad R\ge1.
\]
All heights below are stable Faltings heights in Faltings's original
normalization. We prove
\begin{equation}\label{eq:original-cm-bound}
 h_{\rm F}(A)<854R<6^{100}R^{24}.
\end{equation}
No assertion about the existence or degree of $L$ is included in this lemma.

First, $A_{\overline L}$ is isogenous to $E^2$ for a CM elliptic curve
$E$. Here is a dimension argument which makes the alternative cases
explicit. If $A$ were geometrically simple, its division endomorphism
algebra would act faithfully on the four-dimensional rational homology.
Its dimension therefore divides four. As it contains $B$, it would equal
$B$; but $B$ has no commutative semisimple subalgebra of dimension four,
contrary to CM type. If the two elliptic factors were not isogenous,
their endomorphism algebra would be a product of commutative fields and
could not contain $B$. Finally, if $E$ had no CM, the endomorphism algebra
would be $M_2(\mathbb Q)$, and the embedding of the four-dimensional
division algebra $B$ would be an impossible isomorphism.

The elliptic factors and the isogeny may be taken over $L$. Indeed,
the rational idempotents in $\operatorname{End}^0(A)$ are defined over
$L$, and the image of an integral multiple of such an idempotent is an
elliptic subvariety over $L$. Endomorphisms and homomorphisms of the
resulting factors extend to rational endomorphisms of $A$ by the
inclusions and rational projections; these are fixed by
$\operatorname{Gal}(\overline L/L)$. Thus the factor isogenies, and all
endomorphisms of $E$, are defined over $L$.

Put $F=\operatorname{End}^0(E)$, an imaginary quadratic field, and
$\mathcal O=\operatorname{End}(E)$. Differentiation gives an injective
unital map $F\longrightarrow\operatorname{End}_L(\operatorname{Lie}E)=L$:
it is nonzero, and every nonzero elliptic endomorphism in characteristic
zero is separable. Consequently $F\subset L$.

We next pass to the maximal order without extending $L$. Choose
$f\ge1$ with $f\mathcal O_F\subset\mathcal O$, choose an integral
basis $\alpha_1,\alpha_2$ of $\mathcal O_F$, and form the finite
subgroup scheme
\[
 K=\sum_{i=1}^2(f\alpha_i)(E[f])\subset E.
\]
Each $f\alpha_i$ is defined over $L$, so $K$ and the quotient
$E'=E/K$ are defined over $L$. Under a complex uniformization
$E(\mathbb C)=\mathbb C/\Lambda$, this subgroup is
$\mathcal O_F\Lambda/\Lambda$; thus $E'(\mathbb C)$ has lattice
$\mathcal O_F\Lambda$ and endomorphism ring $\mathcal O_F$.
The ring cannot be larger, since the rational endomorphism algebra
remains $F$. This is also a direct special case of the same-field
maximal-order theorem of R\'emond, but the construction just given is
the proof used here. It follows that $A'=(E')^2$ is $L$-isogenous to $A$.

Write $\Delta_F$ for the fundamental discriminant of $F$. The
discriminant tower formula implies that every rational prime dividing
$\Delta_F$ divides $D_L$. A quadratic fundamental discriminant has
odd prime exponents one and exponent at two at most three. Therefore
\[
 |\Delta_F|\le4\operatorname{rad}(|D_L|)\le4R.
\]
The multiplier four includes the possible factor $2^3/2$.

For completeness we give an elementary analytic bound on $h_{\rm F}(E')$.
CM elliptic curves have potentially good reduction. After a finite
extension with everywhere good reduction, the elliptic discriminant
formula of Silverman, reproduced as equation (6.2) in
\cite{cmheight}, gives
\[
 h_{\rm F}(E')=\frac1{[L':\mathbb Q]}\sum_{\sigma:L'\hookrightarrow\mathbb C}
 \left(-\log(2\pi)-\log|\eta(\tau_\sigma)|^2
                     -\tfrac12\log\operatorname{Im}\tau_\sigma\right).
\]
Here $L'$ is used only to evaluate the stable height; it is not the
field used in the isogeny estimate below. The discriminant term is zero
because reduction is good. We use the individual-embedding formula
(6.2), not the class-average assertion of Proposition 6.1, so no
hypothesis on Weierstrass coefficients lying in $\mathbb Q(j)$ is needed.

For each embedding choose a reduced quadratic parameter of discriminant
$\Delta_F$. If $y=\operatorname{Im}\tau$, then
\[
 \frac{\sqrt3}{2}\le y=\frac{\sqrt{|\Delta_F|}}{2a}
                   \le\frac{\sqrt{|\Delta_F|}}2,
 \qquad a\in\mathbb Z_{\ge1}.
\]
Put $r=|e^{2\pi i\tau}|<1/16$. The product defining $\eta$ gives
\[
 -\log|\eta(\tau)|^2
 \le\frac{\pi y}{6}+2\sum_{n\ge1}\frac{r^n}{1-r^n}
 \le\frac{\pi y}{6}+\frac{2r}{(1-r)^2}
 <\frac{\pi\sqrt{|\Delta_F|}}{12}+\frac{32}{225}.
\]
Also $-\tfrac12\log y\le\tfrac14\log(4/3)<1/12$.
Discarding the negative term $-\log(2\pi)$ and using $\pi<4$ shows
\[
 h_{\rm F}(E')<\sqrt{|\Delta_F|}+1\le2\sqrt R+1\le3R.
\]
The height is additive on products, hence $h_{\rm F}(A')<6R$.

Apply \cite[Theorem 1.9(1)]{GaudronRemond2023} over $L$, with source
$A'$ and target $A$ and with no extension of the base field.
For a surface, its quantity $\Xi(A')$ is
\[
 \Xi(A')=\left(14^{32}d\max(1,\log d,h_{\rm F}(A'))\right)^8.
\]
The general assertion supplies an isogeny $A'\to A$ of degree at most
$\Xi(A')^2$; no identification of CM type with that paper's
separate ``MM'' condition is made or needed. The isogeny-height
inequality then gives
\[
 h_{\rm F}(A)\le6R+
 8\log\left(14^{32}\,576\max(1,\log576,6R)\right).
\]
One can see the last inequality directly from invariant differentials:
after semistable base change, pullback on determinant differentials has
integral cokernel, while its archimedean squared norm is multiplied by
the isogeny degree. Thus the height difference is half the logarithm
of that degree minus a nonnegative normalized cokernel length.

Finally, $\log14<3$, $\log576<7$, and
$\max(1,\log576,6R)\le7R$. Since $\log7<2$ and $\log R\le R$,
\[
 h_{\rm F}(A)<6R+8(96+7+2+\log R)
             \le854R.
\]
As $6^4=1296>854$ and $R\ge1$, this proves
\eqref{eq:original-cm-bound}. In particular the original
$6^{100}R^{24}$ allowance absorbs this CM estimate with the field
degree still exactly bounded by 576.

\paragraph{The unconditional numerical certificate and scalar detector}
The specified integer $H$ satisfies
\[
 \log_2H<4783227.
\]
Indeed the exact sieve enumerates $7499$ primes in $S_*$, with largest
prime $76163$, and their product has bit length $109445$. Hence
$B_*<2^{109445}$. With $A_0=4608^{165888}$, we have
$6^{100}<A_0$ and $128<A_0$, whence $H<8A_0B_*^{24}$. Since $4608<2^{13}$,
\[
 \log_2H<3+165888\cdot13+24\cdot109445=4783227.
\]
This is a check of the number H, not a proof that H is a height constant.

For completeness, the scalar product used by the original polynomial is
\[
 r(\beta)=\beta(\beta-1)(\beta-2401)(32\beta+343)(2\beta-1).
\]
At a prime $p>76176$, if $v_p(\beta)<0$, its valuation is
$5v_p(\beta)$ and so detects odd contact with infinity. If $\beta$ is
$p$-integral and has odd contact with a finite member of $\mathcal C$,
only that factor has positive valuation; every other factor is a unit.
The extra root $1/2$ introduces only exceptional collision primes
(notably factors of $4801$ and $718$). Thus every prime of $M(\beta)$ has
odd valuation in $r(\beta)$. If a separate comparison proves
$M(\beta)^{48}\mid s\ne0$, the five-contact estimate just proved gives
$h(\beta)\le H\sqrt{|s|}$.

Combining the non-CM and CM estimates with \eqref{eq:h5-R} and
\eqref{eq:h5-G} gives
\[
 h(s)\le 2(4608^{165888}+6^{100})(B_*M(s))^{24}+128
       \le H_*M(s)^{24},
\]
because $M(s)\ge1$. The numerical certificate above proves the
remaining assertion about $H_*$.
\end{proof}


Define $\Psi(b)$ by seven rational witnesses $t,U,e,p,q,s,h$ and four
equations, with $B=17b$:
\begin{align}
 U(1+bt^2)&=b(1-bt^2),&
 p^2&=B(1-e)(1+3e),\label{eq:o8-pole}\\
 h&=e(B-U^2),&
 Bq^2&=(s^2-Bh)(s^2+3Bh).\notag
\end{align}

\begin{lemma}[Exact seven-variable pole statement]\label{hyp:o8-pole}
If $b\ne0$ and $\Psi(b)$, then $v_p(b)$ is not a negative odd integer
for every $p\notin\{2,3,17\}$. Conversely, $\Psi(b)$ holds whenever
$b>0$, the denominator of $b$ in lowest terms is a square, and $b$ is
a square in each of $\Q_2,\Q_3,\Q_{17}$.
\end{lemma}
\begin{proof}
These are Lemmas~\ref{lem:original-seven-parity-soundness} and
\ref{lem:original-seven-parity-completeness}. Their proof gives the
simultaneous two-twist construction explicitly, and applies the existing
Theorem~\ref{hyp:RC} and Lemma~\ref{lem:corank}. It introduces no new
simultaneous-realization assumption.
\end{proof}

The lemma includes the degenerate values of the seven witnesses
permitted by \eqref{eq:o8-pole}. No inverse condition on $e$, $h$, $s$,
or $U$ may be inserted into this particular polynomial. For $b\ne0$,
the first equation itself implies $1+bt^2\ne0$. The predicate is soluble
at $b=0$; the norm argument below accounts for that case where necessary.

% Additive proof for the original 15671-variable polynomial. No circuit changes.
\subsubsection{The seven-witness parity predicate and its finite correction}
\label{subsec:original-seven-parity}
Put $B=17b$. The predicate used here is exactly the existence of
$t,U,e,p,q,s,h\in\mathbb Q$ satisfying
\begin{align}
 U(1+bt^2)&=b(1-bt^2),\label{eq:seven-param}\\
 p^2&=B(1-e)(1+3e),\label{eq:seven-first}\\
 h&=e(B-U^2),\label{eq:seven-h}\\
 Bq^2&=(s^2-Bh)(s^2+3Bh).\label{eq:seven-second}
\end{align}
Denote it by $\Psi(b)$. No nonzero conditions are added to this formula.

\begin{lemma}[The exact one-sided assertion]
\label{lem:original-seven-parity-soundness}
If $b\ne0$ and $\Psi(b)$, then $v_\ell(b)$ is not a negative odd
integer for any prime $\ell\notin\{2,3,17\}$.
\end{lemma}
\begin{proof}
Suppose $k=v_\ell(b)<0$ is odd. Equation~\eqref{eq:seven-param}
forces $1+bt^2\ne0$: if it vanished, its right side would be $2b\ne0$.
If $t=0$, then $U=b$. Otherwise $v_\ell(bt^2)$ is odd and nonzero,
so $(1-bt^2)/(1+bt^2)$ is a unit. Thus $v_\ell(U)=k$ in all cases,
and $v_\ell(B-U^2)=2k$.

If $e=0$, equation~\eqref{eq:seven-first} gives $p^2=B$, impossible.
Write $j=v_\ell(e)$. If $j<0$, the right side of
\eqref{eq:seven-first} has valuation $k+2j$, which is odd.
If $j>0$, it has valuation $k$, again odd. Consequently $j=0$.
This includes the possibilities $p=0$, $e=1$, and $e=-1/3$.
Now $h\ne0$ and $v_\ell(Bh)=3k$ is odd. If $s\ne0$, its square
has even valuation, different from $3k$, so the two factors on the
right of~\eqref{eq:seven-second} have the same finite valuation.
Their product has even valuation and is nonzero. The same conclusion
holds when $s=0$, since the product is $-3(Bh)^2$.
The left side is either zero or has odd valuation, a contradiction.
\end{proof}

The omitted guards therefore cause no loss of the asserted soundness.
At $b=0$, the equations require $U=p=h=s=0$ and permit arbitrary
$t,e,q$. For $b\ne0$, the six-witness source formula is recovered
exactly by $e_{\rm old}=Be$, $p_{\rm old}=Bp$, $q_{\rm old}=Bq$,
with $s,t,U$ unchanged. Its first curve residual is $B^2$ times
\eqref{eq:seven-first}'s residual and its second is $B$ times
\eqref{eq:seven-second}'s residual after~\eqref{eq:seven-h}.
Division by $B$ reverses this transformation.

\begin{lemma}[Simultaneous realization from the existing restricted converse]
\label{lem:original-seven-simultaneous-realization}
Use Theorem~\ref{hyp:RC}, the restricted converse $\mathsf{RC}$ for
$E_l:y^2=x(x-l)(x+3l)$ with $l>0$ squarefree and $l\equiv11\pmod{12}$.
Let $\beta,\theta>0$ be squarefree, with disjoint supports avoiding
$2,3$, such that $\beta\equiv17\pmod{24}$,
$\theta\equiv1\pmod{24}$, and each is a local square at every prime
dividing the other. Then there is a positive squarefree $n$, prime
to $6\theta$, for which both $E_n$ and $E_{n\theta}$ have rational
points with nonzero abscissa in $\beta\mathbb Q^{\times2}$.
The case $\theta=1$ is allowed.
\end{lemma}
\begin{proof}
Here are the descent steps, so that no separate realization theorem
is being assumed. Write $B_0=\operatorname{supp}(\beta)$,
$J=\operatorname{supp}(\theta)$, and $d=|B_0|\ge1$.
For a support prime $q_i$, put $s_i=[-1/q_i]$, $\lambda_i=[3/q_i]$,
and define over $\mathbb F_2$
\[
 (Kv)_i=\sum_{j\ne i}[q_j/q_i](v_j+v_i).
\]
Add $d-1$ prime slots $I$ of type $(s,\lambda)=(0,0)$ and one slot
$q_0$ of type $(1,1)$. Let the $B_0$-by-$I$ symbol matrix have columns
$e_i+e_d$ for $1\le i<d$. Set all other symbols incident with $I$
to zero, prescribe $[q_0/q_b]$ as the corresponding row sum, and
obtain reverse symbols by reciprocity. The $q_0$--$J$ symbols remain
free. Mutual squareness gives zero row sums for both directed
$B_0$--$J$ blocks; hence $K\mathbf1_{B_0}=0$ on both supports.

Choose $a$ with ones at one fixed member of $B_0$ and at $q_0$.
The matrix $K^\sharp=K+\lambda a^t$ kills $\mathbf1$ on both sides.
After deleting $q_0$, its smaller-support matrix has the form
\[
 M_0=\begin{pmatrix}D&C\\C^t&0\end{pmatrix},
 \qquad \mathbf1^tD\mathbf1=1.
\]
The columns of $C$ span the even-sum hyperplane. If $M_0(x,y)^t=0$,
then $x=c\mathbf1$; multiplication of the first equation by
$\mathbf1^t$ gives $c=0$, and injectivity of $C$ gives $y=0$.
This also proves the case $d=1$, where $M_0=(1)$.
Adding $J$ preserves $M_0$. Each free $q_0$--$J$ bit toggles just
one new diagonal entry. The enlarged determinant is multilinear in
these bits, with full-product coefficient $\det M_0=1$, so some
Boolean choice makes it nonzero. CRT and Dirichlet's theorem realize
the prescribed symbols by distinct fresh primes, with types imposed
modulo $12$. Thus $n=\prod_{B_0\cup I\cup\{q_0\}}q$ and $n\theta$
are squarefree and congruent to $11$ modulo $12$.

The projection $T=1-\mathbf1_{B_0}a^t$ satisfies $K=K^\sharp T$.
Nonsingularity of the deleted matrix gives
$\ker K^\sharp=\langle\mathbf1\rangle$, and therefore
\[
 \ker K=\langle\mathbf1,\mathbf1_{B_0}\rangle,
 \qquad \lambda\notin\operatorname{im}K
\]
on both supports. For the latter assertion, $Kx=\lambda$ would
make $Tx-\mathbf1_{B_0}$ constant, contradicted by applying $a^t$.

For clarity, the local descent needed here is the manuscript's full
two-descent in coordinates $(x+3l,x-l)$. The two torsion classes are
$(l,1)$ and $(3,-l)$. At an odd good prime, the local image consists
of the four pairs with even valuations; at a prime dividing $l$,
the two torsion classes span it by their independent valuation
vectors. At $3$ it is the set of all first-coordinate classes with
square second coordinate. At $2$ both coordinate valuations are
even: the only problematic case $x-l=2u$, $u$ odd, would give
$(y/2)^2\equiv l\pmod8$, impossible for $l\equiv3\pmod4$.
When both differences are divisible by $4$, their difference $4l$
and the cubic equation give even valuations for both. At infinity
the first coordinate is positive, and the positive-second-coordinate
subgroup has index two. Its possible representatives consequently are
$(3^\delta\prod q_i^{u_i},\prod q_i^{v_i})$ with necessary equations
\[
 Ku+\operatorname{diag}(\lambda)v=\delta\lambda,
 \quad (K+\operatorname{diag}(s))v=0,
 \quad (\lambda+s)^tv=0.
\]
Reciprocity gives $K+K^t=ss^t+\operatorname{diag}(s)$.
Hence $\ker(K+\operatorname{diag}(s))=\ker K^t\cap s^\perp$
is a line. The functional $\lambda^t$ is nonzero on this line,
because $\lambda\notin\operatorname{im}K$ and
$\lambda^t\mathbf1=0$, whereas $s^t\mathbf1=1$.
The displayed equations force $v=0$, $\delta=0$, and $u\in\ker K$.
There are at most eight Selmer classes.
The class $(\beta,1)$ is actually local everywhere: it is the identity
at $2$ and infinity, the first torsion class at $3$, satisfies the
torsion comparison at support primes by $K\mathbf1_{B_0}=0$, and
is unramified at the other odd primes. It is independent of the two
torsion classes since $B_0$ is nonempty and $q_0\notin B_0$.
Thus each two-Selmer group has dimension exactly three.

Lemma~\ref{lem:corank} gives full two-primary Selmer
corank one. Explicitly, the Kummer sequence gives
$r+\dim\Sha[2]=1$; writing
$\Sha[2^\infty]=(\mathbb Q_2/\mathbb Z_2)^t\oplus F$, the
alternating Cassels--Tate pairing on the finite quotient makes
$\dim F[2]$ even, so $F=0$ and $r+t=1$.
Apply $\mathsf{RC}$ to $n$ and $n\theta$, and the manuscript's
published forward analytic-rank-one theorem. Both curves have rank
one and finite $\Sha$, hence $\Sha[2]=0$. Kummer exactness now
represents $(\beta,1)$ by actual rational points. These are not
two-torsion, so their abscissas are nonzero of class $\beta$.
When $\theta=1$ the same point can be used twice.
\end{proof}

\begin{lemma}[The completeness scope actually used]
\label{lem:original-seven-parity-completeness}
Assume only the existing $\mathsf{RC}$ and the published descent and
forward-rank facts used above. If $b>0$ has square reduced denominator
and is a square in $\mathbb Q_2,\mathbb Q_3,\mathbb Q_{17}$, then
$\Psi(b)$ holds.
\end{lemma}
\begin{proof}
Choose $t\in\mathbb Q$ close to a solution of $bt^2=1$ at infinity
and at $2,3,17$. In~\eqref{eq:seven-param} this makes $U$ arbitrarily
small at those places. We can and do require
$V=B(B-U^2)>0$ and $V$ locally square at all three primes.
Explicit sufficient conditions are $U^2<B$ in $\mathbb R$,
$v_2(U^2/B)\ge3$, and $v_3(U^2/B),v_{17}(U^2/B)\ge1$.
The denominator tends locally to $2\ne0$ and is positive globally.

Let $\beta,\theta$ be the positive squarefree representatives of
$B,V$. The number $B=17b$ is square at $2$, a nonsquare unit class
at $3$, and has odd valuation at $17$. Thus
$\beta\equiv17\pmod{24}$ and $\beta>1$;
$\theta\equiv1\pmod{24}$ and $17\nmid\theta$.
At any prime $\ell\mid\beta$ outside $2,3,17$, the valuation of
$b$ is odd and therefore positive, since its reduced denominator is
square. The same quotient calculation as in soundness gives
$v_\ell(U)=v_\ell(b)$; hence $1-U^2/B$ is a square unit.
This proves that $V$, and hence $\theta$, is square at every prime
of $\beta$, including $17$ by construction. The supports are disjoint.
Conversely, if $\ell\mid\theta$, then $\ell\notin\{2,3,17\}$
and $v_\ell(B)$ is even. Odd valuation of $B(B-U^2)$ forces
$v_\ell(B)=2v_\ell(U)$ and strict cancellation, so
$B/U^2\in1+\ell\mathbb Z_\ell$ is square. This proves the other
mutual-square condition. If $U=0$, then $\theta=1$ and this last
argument is unnecessary.

Apply the preceding lemma. Write $V=\theta c^2$ and rescale the
second point by $(x,y)\mapsto(c^2x,c^3y)$ to put it on $E_{nV}$.
If the first point has abscissa $Ba^2$, divide both curve parameters
by $a^2$ and the point coordinates by $a^2,a^3$. The resulting
points have abscissas $B,Bs^2$ with $s\ne0$, on $E_{e_0}$ and
$E_{e_0V}$ respectively. Write their ordinates as $y_1,y_2$.
Set
\[
 e=e_0/B,\qquad p=y_1/B,\qquad
 q=y_2/(B^2s),\qquad h=e(B-U^2).
\]
Direct substitution verifies all four equations. Every division
made in choosing these witnesses is by a nonzero number; none is
an additional equation or guard in the original polynomial.
\end{proof}

\paragraph{Finite correction and the unguarded norm root.}
At $2$ use valuation parity followed by unit coordinates $[-1],[5]$;
at $3$ and $17$ use valuation parity followed by the nonsquare-unit bit.
The columns for $2,3,5,7,11,13,17$ are exactly
\[
\begin{pmatrix}
1&0&0&0&0&0&0\\0&1&0&1&1&0&0\\0&1&1&0&1&1&0\\
0&1&0&0&0&0&0\\1&0&1&0&1&0&1\\0&0&0&0&0&0&1\\
0&1&1&1&1&0&0
\end{pmatrix}.
\]
The accompanying exact computation verifies rank seven and all 128
distinct subset signatures. Rational solutions of $e^2-e=0$ are
exactly $0,1$. Therefore the original eight bits and two products
\[
 A=(1+e_2)(1+2e_3)(1+12e_{13})(1+16e_{17}),\quad
 D=(1+4e_5)(1+6e_7)(1+10e_{11})(1-2\eta)
\]
give exactly $AD=(-1)^\eta\prod p^{e_p}$, always nonzero.
For nonzero integral $r$, choose $\delta=\operatorname{sign}(r)$
and the unique subset product $\mu$ having the same local signature
as $\delta r$. Taking $AD=\delta\mu$ makes $ADr$ positive and
locally square at $2,3,17$.

Choose $J_0$ sufficiently deep at all places over primes dividing
$\mu r$, then $S$ in a sufficiently divisible subgroup satisfying
the anchor. Since $f=z^2$ times a formal unit, the coefficients of
$f(J_0),f(S),z(S)/z(J_0)$ can attain any prescribed finite lower
valuation bounds. This remains true at ramified and bad primes by
working in sufficiently deep formal subgroups. Thus at $\ell\mid r$
all normalized leaves are integral, and at $\ell\mid\mu$ with
$\ell\nmid r$ all nonconstant leaves can be made divisible by $\ell$.
The norm tree $N_r(a,b)=a^2+rb^2$ has its constant leaf $1$ on
the all-left path; its root $W$ is a unit at these primes.
In the actual application $r=\beta(\beta-1)(\beta-2401)
(32\beta+343)(2\beta-1)$ with $\beta=X^2+2$ and $X\in\mathbb Z$.
It is nonzero and $|r|>1$ since $X^2\ne2399$ and $\beta\ge2$.
Consequently $W\ne0$, and $b=ADrW^2$ is positive, locally square
at $2,3,17$, and has square reduced denominator. Completeness applies.

For soundness no global guard $W\ne0$ is needed. At any prime where
$v_\ell(r)$ is odd, the two nonzero summands of $a^2+rb^2$ have
different valuation parity. Such a norm vanishes only if both inputs
vanish. Recursing to the constant leaf $1$ proves $W\ne0$ at every
prime relevant to an odd-contact argument. Outside the fixed
exceptional set, $AD$ is a unit and
$v_\ell(ADrW^2)=v_\ell(r)+2v_\ell(W)$.
The exact original norm argument therefore applies unchanged. If
$r$ has no relevant odd valuation, there is no local conclusion to
prove. No guard or witness has been inserted, and the original
polynomial and coefficient bound are unchanged.

\subsection{Elliptic charts and the anchor}

Use the normalized curve and point
\[
 K=\Q(\sqrt2),\quad d=3-\sqrt2,\quad D=d^2=11-6\sqrt2,
 \quad E:y^2=x^3-Dx,
\]
\[
 P_*=(57/8-7\sqrt2/2,-165/8+417\sqrt2/32),
 \qquad G=2E(K)=\Z P_*.
\]
The identification of $G$ follows from Proposition~\ref{prop:generator}
by the unit change of variables $\lambda=3+2\sqrt2$.
Write $f=1/x$ and $z=-x/y$. A nonzero point of $G$ is represented by
eight field variables $u,v,w,z,f,h,x,y$ satisfying
\begin{equation}\label{eq:o8-chart}
 f=u^2,\quad v^2+df=1,\quad w^2+v^2=2,\quad h=vw,
 \quad zh+u=0,\quad xf=1,\quad yz+x=0.
\end{equation}
Each field equality means its two rational coefficient equalities.
\begin{lemma}[The sparse point chart]\label{lem:o8-chart}
The equations \eqref{eq:o8-chart} parametrize precisely the nonzero points
of $G$, and their reciprocal functions are defined.
\end{lemma}
\begin{proof}
It gives $x=u^{-2}$, $y=vw/u^3$ and the three square root differences
$x=(1/u)^2$, $x-d=(v/u)^2$, $x+d=(w/u)^2$.
Here is a direct verification that these squares give precisely the
required doubles. Write $e_1=0,e_2=d,e_3=-d$ and choose
$r_i^2=x-e_i$ with $r_1r_2r_3=-y$. Put
\[
 B=r_1r_2+r_1r_3+r_2r_3,\qquad
 X'=x+B,\qquad
 Y'=-(r_1+r_2)(r_1+r_3)(r_2+r_3).
\]
Since $X'-e_i=(r_i+r_j)(r_i+r_k)$, the point $(X',Y')$ lies on
$E$. None of the pair sums is zero, because the $e_i$ are distinct.
Its tangent slope is $-(r_1+r_2+r_3)$, and the duplication equations
give $[2](X',Y')=(x,y)$.
Conversely, if $(x,y)=[2](a,b)$ is nonzero, then $b\ne0$ and direct
duplication gives
\[
 x-e_i=\left(
 \frac{(a-e_i)^2-(e_i-e_j)(e_i-e_k)}{2b}\right)^2.
\]
Choose the signs of these square roots to have product $y$, and recover
$u,v,w,h,z,f$ from \eqref{eq:o8-chart}. The group $G$ is torsion-free, so its nonzero points have
$x,y\ne0$. Thus the sparse chart and the membership equations in
Section~\ref{sec:model} have the same nonzero image.
The inverse relation $xf=1$ excludes zero denominators. In particular
$u,h,z,y$ are nonzero.
\end{proof}

Subtraction uses the equations
\[
 \ell x-y=j,\quad \ell a+b=j,\quad
 x'=\ell^2-x-a,\quad \ell x'+y'=j,\quad
 f'x'=1,\quad z'y'+x'=0
\]
for $(x,y)-(a,b)$, with an inverse for the rational norm of $x-a$.
Doubling uses the usual tangent formula with its nonzero denominators.

The baseline duplication estimates give
\[
 4h(x(Q))-\log392\le h(x(2Q))
       \le4h(x(Q))+\log676.
\]
Since $h(x(P_*))=\log41$, passage to the canonical height gives
$2\widehat h(P_*)\ge\log41-\frac13\log392>1$, and hence
\begin{equation}\label{eq:o8-indexheight}
 h(x(nP_*))\ge n^2-6\qquad(n\ne0).
\end{equation}
The deliberately weaker constant six suffices throughout this section.

Let $\mathcal T(a)$ be the seven-variable formula
\begin{equation}\label{eq:o8-T}
 u_0^2-2u_1^2-3u_2^2+6u_3^2=1,\qquad
 (a-2u_0)^2-8v_1^2-12v_2^2+24v_3^2=4.
\end{equation}
Then $\mathcal T(\Q)=\Z_{(2)}\cap\Z_{(3)}$. Indeed the quaternion
algebra $(2,3)$ is ramified exactly at $2,3$ and split at infinity.
Koenigsmann's trace-sum theorem, Proposition~6 of
\cite{koenigsmann-defining}, identifies the sum of two norm-one trace
sets with the intersection of these two local rings. Eliminating the
scalar coordinate of the second norm-one element gives exactly
\eqref{eq:o8-T}, with no translation of $a$.

Introduce twenty-five rational variables
\begin{equation}\label{eq:o8-chain}
 c_0=3,\qquad c_{i+1}=c_i^2\quad(0\le i<24),\qquad C=c_{24}.
\end{equation}
Thus $C=3^{2^{24}}$, but this large integer is not a coefficient in
$F_8$. For $S=sP_*\ne O$ and $f(S)=f_0+f_1\sqrt2$, impose
$Ca_i=f_i$ and $\mathcal T(a_i)$, for $i=0,1$.
At the prime above $3$, $P_*$ reduces to the nonidentity point
$(\sqrt2,0)$ of order two, and $v_3(z(2P_*))=1$.
The formal group therefore gives
$v_3(f(sP_*))=2+2v_3(s)$ whenever this valuation is positive.
The anchor forces
\begin{equation}\label{eq:o8-anchor-bound}
 |s|\ge3^{2^{23}-1}>(20H_*)^2.
\end{equation}
For the last inequality use Theorem~\ref{hyp:o8-height},
$\log_2 3>3/2$, and
$3(2^{23}-1)=25165821>4(5+4783227)=19132928$.
The allowed indices contain a nonzero subgroup: sufficiently divisible
indices make $f(sP_*)/C$ integral at both $2$ and $3$. The anchor can
therefore be imposed together with any finite list of formal divisibilities.

% Original 15,671-witness anchor. No polynomial or constant is changed.
\subsubsection{The literal seven-witness semilocal predicate}\label{subsec:o8-anchor-proof}

For $t\in\mathbb Q$, let $\mathcal T(t)$ mean that there are
$u_0,u_1,u_2,u_3,v_1,v_2,v_3\in\mathbb Q$ with
\[
 u_0^2-2u_1^2-3u_2^2+6u_3^2=1,
 \qquad
 (t-2u_0)^2-8v_1^2-12v_2^2+24v_3^2=4.
\]
Then
\[
 \mathcal T(t)\quad\Longleftrightarrow\quad
 t\in\mathbb Z_{(2)}\cap\mathbb Z_{(3)}.
\]
Indeed, use the quaternion convention $i^2=2$, $j^2=3$ and $ij=-ji$.
Its reduced norm is $x_0^2-2x_1^2-3x_2^2+6x_3^2$ and its reduced trace
is $2x_0$. Put $v_0=(t-2u_0)/2$. The two displayed equations say
exactly that $t$ is a sum of two reduced traces of norm-one elements.
No eighth variable is needed, since $v_0$ is recovered by this rational
expression with fixed nonzero denominator~$2$.

In the notation of \cite[Definition~4 and Proposition~6]{koenigsmann-defining},
this is $T_{2,3}=S_{2,3}+S_{2,3}$. That proposition identifies
$T_{a,b}$ with the intersection of the local integer rings at the
ramified places, with a separate real interval only when infinity is
ramified. Here the Hilbert symbols are $-1$ at $2$ and $3$, and $1$
elsewhere. At $3$ this follows from $(2/3)=-1$; at every other odd
prime both parameters are units. At $2$, writing $2=2^1\cdot1$ and
$3=2^0\cdot3$, the usual exponent is
$0+1\cdot(3^2-1)/8+0=1$. The algebra splits over $\mathbb R$ because
both parameters are positive. Thus the ramification set is precisely
$\{2,3\}$, which proves the claimed equivalence in both directions.
This uses Koenigsmann's sum of two trace sets, not a predicate with
an additional finite offset.

\subsubsection{The original anchor and its exact index subgroup}\label{subsec:o8-anchor-subgroup}

Write $K=\mathbb Q(\rho)$, $\rho^2=2$, and use the literal normalized
curve and generator
\[
 E:y^2=x^3-(11-6\rho)x,
 \qquad
 P=\left(\frac{57}{8}-\frac72\rho,
          -\frac{165}{8}+\frac{417}{32}\rho\right).
\]
As in the original construction, use Proposition~\ref{prop:generator}, transported by the unit change
of variables, to obtain $2E(K)=\mathbb ZP$. For a nonzero $S=nP$, put
$f(S)=1/x(S)=f_0+f_1\rho$. Retain the exact constant chain
\[
 c_0=3,\qquad c_{j+1}=c_j^2\quad(0\le j<24),\qquad C=c_{24}.
\]
It forces $C=3^{2^{24}}$, with $25$ shared rational variables and
$25$ quadratic-or-linear equations. The anchor introduces two
variables $a_i$, imposes $Ca_i=f_i$, and imposes $\mathcal T(a_i)$
for each $i\in\{0,1\}$. It has $16$ private witnesses and six
quadratic equations, exactly as in the saved module.

Set $N=2^{24}$ and $m=N/2=2^{23}$. We claim that, for nonzero $n$,
\[
 \text{the anchor holds on }nP
 \quad\Longleftrightarrow\quad
 n\in 2\cdot3^{m-1}\mathbb Z.
\]
First consider the place above~$3$. It is unramified quadratic, with
residue field $\mathbb F_9$ and normalized valuation $v(3)=1$.
For rational $a,b$ one has
$v(a+b\rho)=\min(v_3(a),v_3(b))$: after dividing by the smaller
power of $3$, linear independence of $1,\bar\rho$ over $\mathbb F_3$
precludes cancellation.
The curve has good reduction there. Direct calculation gives
\[
 \bar P=(\bar\rho,0),\qquad
 z(2P)=
 \frac{3396356159088-4135204317756\rho}{3287811625151},
 \qquad z=-x/y.
\]
The reduction of $P$ has order two. The denominator in this expression
is a $3$-unit and both numerator coefficients have valuation one.
Consequently $nP$ is in the formal kernel exactly when $n$ is even,
and $v(z(2P))=1$.

For completeness, the valuation of every such multiple can be computed
without an unproved uniformity assertion. The integral formal group
in the parameter $z$ has inverse $-z$, so its multiplication series
are odd. Their linear coefficient is the corresponding integer.
Thus, for a $3$-unit integer $u$, $[u]z=uz+O(z^3)$ preserves positive
valuation, whereas $[3]z=3z+O(z^3)$ increases valuation by exactly one
at depth at least one. The higher terms have valuation at least
$3v(z)>v(z)+1$. Apply this repeatedly to $[n/2](2P)$, allowing
negative integers by negation. It gives
\[
 v(z(nP))=1+v_3(n),\qquad
 v(f(nP))=2+2v_3(n)\quad(2\mid n).
\]
The second equality follows from
$f=z^2(1-(11-6\rho)f^2)$ on the formal branch.
The integral formal-group construction and its linear multiplication
coefficient are given in \cite[IV\S1, Example~IV.2.2.3 and
Proposition~IV.2.3(a)]{silverman}. No minimal-model assertion at
the bad prime~$2$ is used below.

At the place above~$2$, normalize $v(2)=1$, so $v(\rho)=1/2$.
The two summands of a rational pair have valuations of opposite
integer parity after multiplication by~$2$, and hence cannot cancel.
The displayed coordinates give
\[
 v(x(P))=-3,\qquad v(y(P))=-9/2,\qquad v(z(P))=3/2.
\]
Thus $P$ lies on the formal branch at infinity, even though reduction
at~$2$ is bad. Every integral multiple lies in this formal subgroup.
One can see the integral branch directly by setting $w=-1/y$:
the equation becomes $w=z^3-(11-6\rho)zw^2$, which has a unique
solution $w\in z^3\mathcal O[[z]]$; its group law is the integral
formal group just cited. In particular $f(nP)$ is integral at~$2$.
Since
\[
 v(a+b\rho)=\min(v_2(a),v_2(b)+1/2),
\]
its rational coefficients are both $2$-integral. As $C$ is a $2$-unit,
the two anchor conditions at~$2$ impose no further restriction on $n$.

At~$3$, the two anchor conditions are equivalent to
$v(f(nP))\ge N$. Positive valuation of $f$ implies negative valuation
of $x$, hence membership in the formal kernel, so $n$ is even.
The preceding exact formula makes the remaining condition
$v_3(n)\ge m-1$. Conversely these two conditions give both required
coordinate divisibilities at~$3$, and the condition at~$2$ already
holds. The proved characterization of $\mathcal T$ supplies its
rational witnesses, proving the subgroup claim.

In particular, anchor indices contain a nonzero finite-index subgroup
and have positive representatives. They can be combined with any
finite collection of requirements that indices be divisible by fixed
integers: take a common multiple with $2\cdot3^{m-1}$. More generally,
any finitely many formal-neighborhood conditions whose preimages are
finite-index subgroups can be imposed in the same way. This statement
does not claim compatibility with arbitrary nonzero residue conditions.

The old and new reciprocal coordinates are equivalent for this anchor.
With $\lambda=3+2\rho$, one has $f_{\rm old}=\lambda^2f$ and
\[
 \begin{pmatrix}f_{{\rm old},0}\\f_{{\rm old},1}\end{pmatrix}
 =\begin{pmatrix}17&24\\12&17\end{pmatrix}
 \begin{pmatrix}f_0\\f_1\end{pmatrix}.
\]
This matrix has determinant one and integral inverse, so it preserves
coordinatewise membership in
$C(\mathbb Z_{(2)}\cap\mathbb Z_{(3)})^2$ in both directions.

\subsubsection{The unchanged numerical cutoff}\label{subsec:o8-anchor-cutoff}
In this calculation $H=H_*$ is the constant in Theorem~\ref{hyp:o8-height}.

The anchor gives $|n|\ge3^{2^{23}-1}$. Conditional only on the separately
audited original height bound $\log_2H<4783227$, it follows that
$|n|>(20H)^2$. Indeed $\log_2 3>3/2$ and $\log_2 20<5$, and
\[
 3(2^{23}-1)=25165821
 >19132928=4(5+4783227).
\]
Thus the original $24$ squarings suffice whenever that exact height
input is established. This calculation does not replace the required
five-contact height theorem by a six-contact theorem or by a larger
constant. No part of this proof changes the saved polynomial.

\subsection{The exact norm tests}

For a nonzero rational contact $r$ and rational tests $X_i$ of positive
orders $L_i$, apply the binary operation
$N_r(a,b)=a^2+rb^2$ to the ordered leaves
$1,X_1/r^{L_1},\ldots,X_k/r^{L_k}$ in the following fixed order.
First pair $1$ with $X_1/r$, where $L_1=1$. Group the remaining
tests by increasing order, retaining their original order within each
group. Pair successive tests within each group and carry any last
unpaired test. Then pair successive resulting nodes at each level and
carry any last node until only a root $W$ remains. The polynomial
equations for the bottom pairs are
\[
 rw-r-X_1^2=0,\qquad
 r^{2L}w-X_i^2-rX_j^2=0;
\]
an unpaired normalized leaf satisfies $r^Lw-X_i=0$.
Higher nodes satisfy $w-a^2-rb^2=0$. Powers of $r$ are computed
by binary squaring. Add an inverse for $r$, a variable for $W^2$,
and the predicate $\Psi(rW^2)$.

At a prime with $R=v_p(r)$ odd, the two summands in every norm node
have distinct valuation parity. There is no cancellation. The constant
leaf is on the all-left path, so $v_p(W)\le0$ if $R<0$.
Consequently $v_p(rW^2)>0$ first implies $R>0$.
It also gives $v_p(W)>-R/2$. Descending any norm node, a bound
$v_p(N_r(a,b))>-cR$ with $c<1$ gives
$v_p(a)>-cR/2$ and $v_p(b)>-(c+1)R/2$.
Every normalized leaf therefore has valuation greater than $-R$.
Thus $v_p(X_i)>(L_i-1)R\ge L_i-1$, or $v_p(X_i)\ge L_i$.
The root cannot vanish at such a prime. Indeed $N_r(a,b)=0$ forces
$a=b=0$ over $\Q_p$, and descent would make its constant leaf zero.
Thus Lemma~\ref{hyp:o8-pole} applies even though the polynomial
has no separately imposed inverse for $W$.

For completeness at a positive prime $r=p$, tests of orders $L_i$
can be made so deep that every nonconstant normalized leaf is
$p$-integral and the root $W$ is a $p$-adic unit. Then $rW^2$ has
square reduced denominator. If $r$ is a local square at $2,3,17$,
the completeness half of Lemma~\ref{hyp:o8-pole} applies.

\subsection{Forty-one contacts and one common comparison point}

Set
\[
 R=47027=31\cdot37\cdot41,\quad W_0=25878772920,
 \quad c=18484837801,\quad V_0=32W_0.
\]
For nine rational targets $a_i$, put $A_i=Ra_i+3$ and
$w_i=A_iv-u$. Use the forty-one nonzero contacts
\begin{equation}\label{eq:o8-contacts}
 v/V_0,\qquad u-\delta v\ (\delta\in\mathcal D),\qquad
 -(w_i-\delta v)\ (\delta\in\mathcal D,\ 1\le i\le9).
\end{equation}
For integer $a_i$, the prime-pattern argument of
Section~\ref{sec:mult} applies on
$v=32W_0(c+W_0y)$, $u=c+W_0x$ and gives positive ordinary prime
values for every distinct contact, all congruent to $c$ modulo $W_0$.
Repeated targets use the same prime values. These are local squares
at $2,3,17$. The precise application uses at most forty-one distinct
affine forms and their successor exclusions. At $31,37,41$ there are
at most eight slopes since every $A_i$ reduces to $3$; at $p\ge43$
there are at most forty. Choosing two slopes outside that set proves
local admissibility. The small-prime conditions are built into $c$.
The full uniform fiber sieve and same-prime order argument are given
in Subsection~\ref{subsec:o8-pattern-proof} below. No conditions relating orders at
different selected primes are needed.

An odd normalized contact of $u/v$ or $w_i/v$ outside $S_*$ forces
some raw contact in \eqref{eq:o8-contacts} to have odd valuation.
To see this, let $m$ be the minimum of the two coordinate valuations.
For a finite contact of odd depth, if $m$ is even its distinguished
linear form has odd valuation; if $m$ is odd another finite contact
form has valuation $m$. At an odd contact with infinity the same
argument uses the denominator if $m$ is even and a finite form if
$m$ is odd. All differences of distinct finite contact points are units
outside $S_*$. The signs and $V_0$ do not affect parity there.

For nonzero points $Q_i=n_iP_*$, $1\le i\le4$, take $Q_0=P_*$,
one common $J=jP_*\ne O$, and four rational proxies $b_i$.
At each raw contact introduce $X_0=hP_*$ and $X_i=m_iP_*$ and
the subtraction charts $Y_i=X_i-4Q_i$, $0\le i\le4$.
The points $4Q_i$ are computed by two doublings. Put
\[
 L(Z)=\sum_{j=0}^{15}
       \frac{(-1)^j\binom{2j}{j}D^j}{4j+1}Z^{4j+1}.
\]
This polynomial has degree sixty-one; the first omitted logarithm term
has degree sixty-five. Test order one for
\[
 f(X_i)\ (0\le i\le4),\quad f(J^\sigma),\quad
 f(Y_i^\sigma)\ (0\le i\le4),
\]
and order sixty-four for
\[
 L(z(X_i))/L(z(X_0))-b_i\ (1\le i\le4),\qquad
 z(Y_i^\sigma)/z(J^\sigma)\ (0\le i\le4).
\]
The inverses in these formulas are imposed explicitly. For each field
test $a+b\sqrt2$ of order $L_i$, use the rational tests
$(\theta^2-2)b$ and $a+\theta b$ of that order, together with
$\theta^2-2$ of order one. There is one common rational $\theta$
for all forty-one comparisons. This gives twenty-three rational tests
of order one and eighteen of order sixty-four.

At an odd raw contact the norm tests select a place above $p$, and
the conjugate parameter ratios imply
$h\equiv4$ and $m_i\equiv4n_i\pmod {p^{64}}$.
In particular $h$ is a unit. The logarithm expansion and its degree
sixty-five remainder give
\begin{equation}\label{eq:o8-proxies}
 b_i\in\Z_p,\qquad b_i\equiv n_i\pmod {p^{64}}.
\end{equation}
This uses the same formal-group estimates as Section~\ref{sec:mult}.

For completeness choose a consistent orientation at each distinct
selected prime. Let $d_{1,p},d_{2,p}$ be the two reduction orders.
The ordinary-prime conditions give $\gcd(d_{1,p},d_{2,p})=4$ and
$p\nmid d_{1,p}d_{2,p}$. Choose
$j=\operatorname{lcm}_p d_{2,p}$. For each occurrence at $p$ choose
\[
 h\equiv0\pmod {d_{1,p}},\qquad
 h\equiv4\pmod {d_{2,p}},\qquad
 h\equiv4\pmod {p^{64+v_p(j)}}.
\]
These congruences are compatible. Avoid $h=0,4,-4$, and set
$X_i=hn_iP_*$. They give all the required formal conditions and
parameter ratios relative to the single point $J$. Increase the last
precision if necessary and choose $\theta$ last by the Chinese
remainder theorem. Repeated occurrences of a prime use the same
orientation. Every subtraction denominator is nonzero, and the
formal-logarithm unit calculation at $p$ proves $L(z(X_0))\ne0$.

% Independent arithmetic justification for the unchanged original 15671-witness circuit.
\subsubsection{The original forty-one-prime pattern}\label{subsec:o8-pattern-proof}
Let
\[
 \mathcal D=\{0,1,2401,-343/32\},\quad R=47027=31\cdot37\cdot41,
 \quad W=25878772920=4\prod_{q\le29}q,
\]
\[
 c=18484837801,\qquad V_0=32W=828120733440.
\]
For integers $a_1,\ldots,a_k$, $k\le9$, put $A_i=Ra_i+3$,
$w_i=A_iv-u$. There are $5+4k$ displayed positions
\begin{equation}\label{eq:old41-list}
 v/V_0,\qquad u-\delta v,\qquad
 u-(A_i-\delta)v\quad(\delta\in\mathcal D,\ 1\le i\le k).
\end{equation}
For every lower bound these positions can simultaneously be positive
primes beyond that bound, congruent to $c$ modulo $W$, with the
following property. For the curve
$E:y^2=x^3-(3-\sqrt2)^2x$, both reductions above each prime are good,
and their cardinalities $d_1,d_2$ satisfy
\[
 \gcd(d_1,d_2)=4,\qquad p\nmid d_1d_2.
\]
This is a same-prime statement and asserts no cross-prime gcd condition.

\begin{proof}
Use $u=c+Wx$, $v=V_0(c+Wy)$. Put
\[
 \mathcal S=\mathcal D\cup\bigcup_i(A_i-\mathcal D),\quad
 K_s=V_0s\in W\mathbb Z,
 \quad\Psi_0=c+Wy,\quad\Psi_s=c+Wx-K_s(c+Wy).
\]
Remove duplicate targets. The remaining slopes are all distinct.
Indeed, the nonzero integral differences in $\mathcal D-\mathcal D$
are $\pm1,\pm2400,\pm2401$, none divisible by $31$.
The integral sums in $\mathcal D+\mathcal D$ are
$0,1,2,2401,2402,4802$, with residues $0,1,2,14,15,28$ modulo $31$,
none equal to $A_i\equiv3$. These facts exclude intersections between
secondary families and between primary and secondary families.
The $m=1+|\mathcal S|\le41$ forms have pairwise independent linear
parts. Repeated targets may be restored afterwards with repeated values.

At $q\mid W$, the forms equal $c$; direct reduction gives
$c\equiv1\pmod8$, $c\equiv3\pmod7$, and $c\equiv1\pmod q$ at
all other odd primes $q\le29$. Thus the forms and their successors
are units at every odd $q\mid W$. For $q>29$, the affine map from
$(x,y)$ to $(u,v)$ is invertible. At $q=31,37,41$, all $A_i=3$
modulo $q$, so $|\mathcal S\bmod q|\le8$. For $q\ge43$ this
cardinality is at most $40$. In either case choose distinct residues
$b,e$ outside the slopes. Put $\gamma=V_0^{-1}$ modulo $q$ and,
reversing $b,e$ if necessary, require $e-b\ne-\gamma$. Then
\[
 v=(e-b)^{-1},\quad u=bv
\]
satisfies
\[
 u-sv=v(b-s)\ne0,\quad u-sv+1=v(e-s)\ne0,
 \quad\gamma v\ne0,\quad\gamma v+1\ne0.
\]
This proves all the value and successor local conditions. No condition
on a successor is imposed at $2$.

Add an unused integer coordinate $z$, and choose a fixed rectangular
box $\Omega$ inside the cone $v>0$, $u/v>\max(0,|A_1|,\ldots,|A_k|)+2413$ in the
homogeneous leading coordinates. All forms on $X\Omega$ lie between
$\alpha X$ and $\beta X$ for positive fixed constants. Fix $Q>82$
containing every prime divisor of $W$ and every nonzero $K_s-K_t$.
Impose one good residue at every prime up to $Q$. For a cutoff $L>Q$,
exclude $\Psi_\nu=0,-1$ at every additional prime up to $L$.
The surviving proportion is
\[
 \gg_{\mathbf a}\prod_{Q<q\le L}(1-2m/q)
 \gg_{\mathbf a}(\log L)^{-2m},
\]
by \cite[Theorem~1, modulus~2]{williams-mertens}.

For fixed $L$, apply \cite[Main Theorem, p.~1763]{greentao} to every
surviving lattice residue in the convex box. The singleton partition
in \cite[Definition~1.5 and Lemma~1.6]{greentao} bounds complexity by
$m-2\le39$. The inputs $\mathrm{MN}(39)$ and $\mathrm{GI}(39)$ follow
from \cite[Theorem~1.1]{moebius} and \cite[Theorem~1.3]{gtz}.
At imposed primes the local factor is $(1-1/q)^{-m}$; at the other
primes it is $1+O(m^2/q^2)$. The singular product therefore has a
positive lower bound independent of $L$. All coefficient and modulus
bounds are fixed before $X$ tends to infinity. Summing the finitely
many asymptotics and removing proper prime powers gives
\begin{equation}\label{eq:old41-lower}
 \#\{\text{prime patterns}\}\ge
 c_{\mathbf a}\frac{X^3}{(\log L)^{2m}(\log X)^m}.
\end{equation}
The constant $c_{\mathbf a}>0$ is independent of $L$; the threshold
for $X$ may depend on it. Proper powers contribute at most
$O_{\mathbf a}(X^{5/2}(\log X)^{m+1})$ before removing von Mangoldt
weights, since each fixed value has $O_{\mathbf a}(X^2)$ fibers.

We give the uniform upper estimates needed to impose the order condition.
For forbidden sets of cardinality $\rho(q)$ in a two-dimensional box
of side $N$, put $g(q)=\rho(q)/q^2$, $h(q)=g(q)/(1-g(q))$, and extend
multiplicatively on squarefree support. CRT gives
\[
 A_d=N^2g(d)+R_d,\qquad |R_d|\ll Nd+d^2.
\]
With $G(Z)=\sum_{d\le Z}h(d)$, take
\[
 \lambda_d=\frac{\mu(d)}{G(Z)\prod_{q\mid d}(1-g(q))}
 \sum_{\substack{e\le Z/d\\(e,d)=1}}h(e)\quad(d\le Z).
\]
Then $\lambda_1=1$, $|\lambda_d|\le d$, and the Selberg quadratic
form equals $G(Z)^{-1}$. To verify the last assertion, use
$y_r=\sum_{r\mid d}\lambda_dg(d)=\mu(r)h(r)/G(Z)$ and
$1/g(n)=\sum_{r\mid n}1/h(r)$. The number of sifted points is thus
\begin{equation}\label{eq:old41-upper}
 \ll N^2/G(Z)+NZ^6+Z^8.
\end{equation}
For $\rho(q)=2q-1$, supported on $q\equiv1\pmod4$ with at most one
prime omitted, $G(Z)\gg_Q\log Z$, uniformly in that omission.
For $\rho(q)=(m-1)q$ at all $q>Q$,
$G(Z)\gg_{Q,m}(\log Z)^{m-1}$. These follow from
\cite[Theorem~1, moduli~4 and~2]{williams-mertens}. Truncation is
justified by using $z=Z^{1/512}$ and the probability distribution
$h(d)/\prod_{q\le z}(1+h(q))$: its expected $\log d$ is at most
$40\sum_{q\le z}(\log q)/q<\tfrac12\log Z$ for large $Z$.
Thus at least half the mass lies below $Z$. Removing one norm-sieve
prime changes its product by a factor at least $1/2$.

For a selected prime, $(7/p)=-1$ and $p\equiv1\pmod8$. The two
reductions of $3-\sqrt2$ have product $7$ and opposite square classes,
so their traces are opposite. Their discriminants $64d^6$ are units.
The automorphism $(x,y)\mapsto(-x,iy)$ is rational over $\mathbb F_p$.
By the Tate-module injection \cite[III.7.4]{silverman}, its rational
centralizer is $\mathbb Q(i)$: at the auxiliary prime $3$ the
centralizer has dimension two because $T^2+1$ is irreducible.
Frobenius is integral by \cite[V.2.3.1(b)]{silverman}; hence it is
$a+bi\in\mathbb Z[i]$. The trace and determinant formulas
\cite[III.8.6]{silverman} give
\[
 a^2+b^2=p,\qquad d_1=p+1-2a,\qquad d_2=p+1+2a.
\]
Full rational two-torsion gives $4\mid d_1,d_2$, and
$v_2(d_1+d_2)=2$, so their gcd has $2$-part exactly $4$.
Also $0<d_i<2p$ and $d_i\ne p$, so $p\nmid d_1d_2$.
An odd common divisor $\ell$ divides $a$ and $p+1$, and forces
$b^2\equiv-1\pmod\ell$. The successor conditions exclude $\ell\le L$.

The number of bad prime values of any fixed form is
\begin{equation}\label{eq:old41-bad}
 \ll_{\mathbf a}X/(L\log X)+X^{7/8}+\sqrt X\log X.
\end{equation}
Indeed, for $L<\ell\le X^{1/8}$, write
$a=\ell r$, $b=b_0+\ell s$, with at most two $b_0$ modulo $\ell$.
The box has side $N\asymp\sqrt X/\ell\gg X^{3/8}$.
Use the norm case of \eqref{eq:old41-upper} with $Z=X^{1/100}$;
its error is uniformly $o(N^2/\log X)$. Sum
$O(X/(\ell^2\log X))$. For $\ell>X^{1/8}$, $a\ne0$ gives
$O(\sqrt X/\ell)$ choices with no additive constant, and the
$b$ count is $O(\sqrt X/\ell+1)$. Summation proves the remaining terms.

On a fiber $\Psi_0=p$, use $(x,z)$; the other forms are
$c+Wx-K_s p$. On a fiber $\Psi_{s_0}=p$, use $(y,z)$ and
\[
 x=K_{s_0}y+\frac{p-c(1-K_{s_0})}{W},\qquad
 \Psi_s=p+(K_{s_0}-K_s)\Psi_0.
\]
The quotient is integral for an actual fiber. For
$Q<q\le X^{1/100}<p$, the remaining zero lines are distinct and
parallel, with total cardinality $(m-1)q$. The second case of
\eqref{eq:old41-upper} bounds each fiber by
$O_{\mathbf a}(X^2/(\log X)^{m-1})$, independently of $p,L$.
Multiply by \eqref{eq:old41-bad} and sum over $m$ positions.
The bad patterns are at most
\[
 C_{\mathbf a}\frac{X^3}{L(\log X)^m}
 +O_{\mathbf a}\left(
 X^{23/8}/(\log X)^{m-1}+X^{5/2}/(\log X)^{m-2}\right).
\]
Choose $L$ so large that $C_{\mathbf a}/L$ is less than half
$c_{\mathbf a}/(\log L)^{2m}$, then choose $X$ large.
A good pattern survives, proving the assertion.
\end{proof}

\subsubsection{A shared auxiliary point with private shifts}\label{subsec:o8-common-J-proof}
Let $\mathcal P$ be the finite set of distinct selected primes, all
larger than $76176$. At each $p$ fix an orientation of its two places,
and let $d_{p,1},d_{p,2}$ be the corresponding group cardinalities.
Set
\[
 j=\operatorname{lcm}_{p\in\mathcal P}d_{p,2}>0,
 \qquad J=[j]P_*.
\]
For each $p$, independently choose an integer $h_p$ satisfying
\begin{equation}\label{eq:old41-private-crt}
 h_p\equiv0\pmod{d_{p,1}},\qquad
 h_p\equiv4\pmod{d_{p,2}},\qquad
 h_p\equiv4\pmod{p^{64+v_p(j)}}.
\end{equation}
The first two congruences are compatible because their modulus gcd
is $4$. The last is compatible because $p\nmid d_{p,1}d_{p,2}$.
The resulting arithmetic progression contains an $h_p$ outside
$\{0,4,-4\}$. Put
\[
 X_0=[h_p]P_*,\qquad X_i=[h_p n_i]P_*,\qquad
 Y_i=[(h_p-4)n_i]P_*,\quad n_0=1.
\]
These are private for each contact occurrence, whereas $J$ is common.
Use the same orientation at repeated occurrences of $p$.

At the primary place all $X_i$ are formal; at the opposite place
$J$ and all $Y_i$ are formal. The formal logarithm is a homomorphism
and an isometry on $p\mathbb Z_p$ by
\cite[IV.6.4 and its proof; VII.2.1--2.2]{silverman}. The hypotheses hold
because $p$ is odd, split, and of good reduction. For any integers
$k,j$ whose multiples of $P_*$ lie in the formal group,
\[
 \frac{\log([k]P_*)}{\log([j]P_*)}=\frac{k}{j}.
\]
To see this without defining a logarithm of $P_*$ itself, apply the
homomorphism law to $[j]([k]P_*)=[k]([j]P_*)$. Nonzero formal
multiples have nonzero logarithm. Therefore
\[
 v_p\left(\frac{z(Y_i)}{z(J)}\right)
 =v_p\left(\frac{(h_p-4)n_i}{j}\right)\ge64,
 \qquad
 \frac{\log(X_i)}{\log(X_0)}=n_i.
\]
This proves the required local comparison conditions. Although other
$d_{q,2}$ may be divisible by $p$, their effect is exactly the term
$v_p(j)$ in \eqref{eq:old41-private-crt}; no cross-prime gcd assertion
is needed.

For the original truncation
\[
 L(z)=z\sum_{r=0}^{15}\frac{(-1)^r\binom{2r}{r}}{4r+1}(Dz^4)^r,
\]
the full logarithm has the same series extended over all $r\ge0$.
This follows by integrating
$(1+4Dz^4)^{-1/2}$, the invariant differential in the parameter
$z=-x/y$. If $v_p(z)\ge1$ and $p>76176$, every omitted term has
valuation at least $65v_p(z)$: its integer numerator can only increase
valuation, and for $n=4r+1\ge65$,
$n v_p(z)-v_p(n)\ge65v_p(z)$.
For $n<p$ this is immediate; for $n\ge p$ use $n-65\ge v_p(n)$.
Thus $L(z)/\log(z)\in1+p^{64}\mathbb Z_p$. The required primary
truncated-log ratios are congruent to the integers $n_i$ modulo $p^{64}$.

Conversely, in an accepted comparison at any $p>76176$, its conjugate
parameter-ratio tests imply
\[
 m_i-4n_i\equiv0\pmod{p^{64}},\qquad h-4\equiv0\pmod{p^{64}},
\]
since $j$ is a nonzero integer. Hence $h$ is a $p$-unit. The same
logarithm identity and truncation estimate give
$L(X_i)/L(X_0)\equiv m_i/h\equiv n_i\pmod{p^{64}}$ and show
that this ratio is integral. The rational proxy test then gives the
claimed integrality and congruence for every $b_i$.

\subsubsection{The common root approximation and parity boundary}\label{subsec:o8-root-proof}
For each distinct selected $p$, retain the chosen square root
$\rho_p$ of $2$. After constructing the private points, each field
test $A+B\sqrt2$ of order $k$ has $v_p(A+B\rho_p)\ge k$.
Choose a common integer $\theta$ sufficiently close to $\rho_p$ at
every selected prime that
$v_p(\theta-\rho_p)\ge k-v_p(B)$ for every nonzero coefficient $B$,
and at least one. There are finitely many requirements; CRT supplies
such an integer. Since $\theta+\rho_p$ is a unit,
\[
 v_p((\theta^2-2)B)\ge k,\qquad v_p(A+\theta B)\ge k.
\]
Repeated primes require the same root and impose compatible precision
bounds. Conversely, the root test and these two rational tests imply
the field test at the selected root. This verifies the circuit's one
common $\theta$, not a separate root variable per contact.

Every raw contact parameter in the complete assignment is a positive
prime congruent to $c$ modulo $W$. It is therefore a square in
$\mathbb Q_2,\mathbb Q_3,\mathbb Q_{17}$. The binary norm batch has
unit root at its contact prime, because normalized leaves are integral
and its leftmost constant leaf is $1$. Its root is positive and nonzero
in the real embedding. The argument $r\mathcal W^2$ consequently is
positive, has square reduced denominator, and is locally square at
$2,3,17$. These are the exact hypotheses of the original specialized
parity-completeness input. This verifies its consumption but does not
constitute a new proof of that input, which is supplied here by
Lemma~\ref{lem:original-seven-parity-completeness}. In particular no integrality of
the comparison points at these three small primes is required here.

\subsection{Scalar height and the defining relation}

Put $\beta=X^2+2$ and
\[
 r(\beta)=\beta(\beta-1)(\beta-2401)(32\beta+343)(2\beta-1).
\]
This is nonzero for every rational $X$, since $X^2=2399$ has no
rational solution. Introduce $J_0\ne O$ and apply a scalar norm
batch to the two coefficients of each of
$f(J_0),f(S),z(S)/z(J_0)$, of orders $1,1,1,1,48,48$.
In this batch replace the pole argument by $\delta\mu rW^2$, where
$\delta\in\{1,-1\}$ and $\mu$ is a squarefree product of
$2,3,5,7,11,13,17$. Precisely, eight Boolean variables produce
\[
 A=(1+e_2)(1+2e_3)(1+12e_{13})(1+16e_{17}),\quad
 B=(1+4e_5)(1+6e_7)(1+10e_{11})(1-2\eta),
\]
and the pole argument is $ABrW^2$. These factors have no valuation
outside $S_*$. Every odd five-contact of $\beta$ gives an odd
valuation of $r(\beta)$ there. At infinity its valuation is five
times that of $\beta$; at a finite contact all other factors are units.
Thus the scalar tests force $M_5(\beta)^{48}\mid s$, and
Theorem~\ref{hyp:o8-height} gives
\begin{equation}\label{eq:o8-scalarheight}
 h(\beta)\le H_*\sqrt{|s|}.
\end{equation}

For integral $X$, the scalar $r$ is a nonzero integer. The square
classes of $2,3,5,7,11,13,17$ form a basis of
$\Q_2^*/\Q_2^{*2}\times\Q_3^*/\Q_3^{*2}
\times\Q_{17}^*/\Q_{17}^{*2}$.
Choose $\delta$ and $\mu$ so that $\delta\mu r$ is positive and
a square at these three places. Choose $J_0$ in the required formal
subgroups and then choose $S$ sufficiently divisible relative to $J_0$
at all primes dividing $r\mu$, together with the anchor conditions.
All normalized nonconstant leaves are then integral at primes dividing
$r$, and vanish at primes dividing $\mu$ but not $r$. The root is
a unit at these primes. Consequently $\delta\mu rW^2$ has square
reduced denominator, and Lemma~\ref{hyp:o8-pole} supplies the
scalar pole witnesses.

Take $A=nP_*,S=sP_*,T=tP_*,Z=zP_*\ne O$, and use
$(\beta,\xi,\tau,\zeta)$ as the four proxies. Impose
\begin{equation}\label{eq:o8-Lucas}
 \tau^2-\xi(\xi-2)\tau\zeta+(\xi-2)^2\zeta^2=1.
\end{equation}
The nine rational targets in \eqref{eq:o8-contacts} are a nonzero
common denominator $e$ and the eight coefficients of
$ex(A),ex(S),ex(T),ex(Z)$. Require its forty-one comparison batches,
the scalar batch, and the anchor just described.

If every associated five-contact radical has its $K$th power dividing
a nonzero integer $N$, Theorem~\ref{hyp:o8-height} and
$Ra_i+3=u/v+w_i/v$ give
\[
 h(a_i)<2H_*|N|^{24/K}+14,
 \qquad h(x(Q))<7H_*|N|^{24/K}
\]
for each of the four points. The second bound follows by applying the
first to $e$ and both coefficients of $ex(Q)$ and absorbing the
constant smaller than fifty into $H_*$.

\begin{lemma}[Lucas growth]\label{lem:o8-lucas}
If $s,t,z\in\Z$, $|s|=a\ge7$, $z\ne0$, and
\[
 t^2-s(s-2)tz+(s-2)^2z^2=1,
\]
then $|t|>\exp(a-3)$. Conversely each such $s$ admits positive
integers $t,z$ satisfying this equation.
\end{lemma}
\begin{proof}
Put $w=\operatorname{sgn}(s)(s-2)z$ and
$\rho=(a+\sqrt{a^2-4})/2$. The equation becomes
$t^2-atw+w^2=1$. Every positive norm-one unit of $\Z[\rho]$ is a
power of $\rho$: multiply it by a power into $[1,\rho)$, write the
result as $v=A+B\rho$, and use
$B=(v-v^{-1})/(\rho-\rho^{-1})\in[0,1)\cap\Z$.
Thus $v=1$. If $U_0=0$, $U_1=1$ and
$U_{k+1}=aU_k-U_{k-1}$, the identities for $\rho^{\pm k}$ give
$|w|=U_k$ and $|t|\in\{U_{k-1},U_{k+1}\}$ for some $k\ge1$.
Modulo $a-2$ one has $U_k\equiv k$, and modulo $a+2$ one has
$U_k\equiv(-1)^{k-1}k$. The divisibility of $w$ by $|s-2|$ therefore
gives $k\ge a-2$ when $s>0$ and $k\ge a+2$ when $s<0$.
Induction gives $U_k\ge(a-1)^{k-1}$, whence
$|t|\ge(a-1)^{a-4}>e^{a-3}$ for $a\ge7$.
For existence take $k=a-2$, $z=U_k/(a-2)$ if $s>0$, and
$k=a+2$, $z=U_k/(a+2)$ if $s<0$, in both cases with $t=U_{k+1}$.
\end{proof}

\begin{proof}[Proof of Theorem~\ref{thm:o8-conditional}]
Suppose the relation holds, and let
$B=\max(|n|,|s|,|t|,|z|)$ and
$\Delta=t^2-s(s-2)tz+(s-2)^2z^2-1$.
If $\Delta\ne0$, the congruences \eqref{eq:o8-proxies} imply that
every coordinate contact radical to the sixty-fourth power divides
$\Delta$. Since $|\Delta|<5B^4$, the index-height and shared-height
bounds give
\[
 B^2-6<7H_*|\Delta|^{3/8}<14H_*B^{3/2}.
\]
Hence $B<(14H_*+6)^2\le(20H_*)^2$, contradicting the anchor.
Thus $\Delta=0$. Lemma~\ref{lem:o8-lucas} gives
$|t|>\exp(|s|-3)$ because $z\ne0$.

If $\beta\ne n$, put $N=|\operatorname{num}(\beta-n)|$ in lowest
terms. The same congruences give
\[
 B^2<7H_*N^{3/8}+6,\qquad N\le2e^{h(\beta)}B,
\]
whence $\log B<\frac8{13}\log(16H_*)+\frac3{13}h(\beta)$.
Writing $a=|s|>(20H_*)^2$, \eqref{eq:o8-scalarheight} implies
$h(\beta)<a/20$ and $\log(16H_*)<a/20$. Therefore
$\log B<11a/260<a-3$, another contradiction.
We obtain $\beta=n\in\Z$, and $X\in\Q$, $X^2=\beta-2\in\Z$
imply $X\in\Z$ by reducing its denominator.

Conversely let $X\in\Z$ and set $A=\beta P_*$. Choose the scalar
height data and then an anchored $S$ with all required divisibilities.
Lemma~\ref{lem:o8-lucas} gives positive integers $t,z$ satisfying the
index equation; set $T=tP_*$ and $Z=zP_*$. Choose an integer $e$
clearing their abscissa coefficients. The ordinary prime pattern
provides \eqref{eq:o8-contacts}. Set the proxies equal to their integer
indices, choose $J$ and the private points as above, and finally choose
the common $\theta$. The norm and pole completeness statements supply
all remaining witnesses. No subsequent choice changes the target points.
\end{proof}

\subsection{The base compiler prescription and finite verification}

Every equation above is implemented with named arithmetic intermediates
so that each rational residual has degree at most four. A field product
uses $(a+b\sqrt2)(c+d\sqrt2)=(ac+2bd)+(ad+bc)\sqrt2$.
Scalar inverses are equations $uv=1$. Field inverses use the two
coefficient equations of the field product being one; subtraction has
the separate scalar inverse of its abscissa-difference norm.
The logarithm uses Horner evaluation in $(Dz^4)^2$ with block length
two. The two doublings of the fixed point $P_*$ also use arithmetic
intermediates; inserting its large precomputed fourth multiple as a
coefficient would give a different polynomial. The chart is exactly
\eqref{eq:o8-chart}, including the intermediate $h=vw$.

For each nonzero rational residual, multiply by the least common multiple
of its coefficient denominators, then divide all coefficients by their
positive greatest common divisor. Denote the resulting integral rows
by $R_1,\ldots,R_{15107}$. In the variable order specified by
\texttt{degree8-original/build.py} and its residual table, define
\begin{equation}\label{eq:o8-polynomial}
 F_8=\sum_{i=1}^{15107}R_i^2.
\end{equation}
This equation is an exact finite prescription for every coefficient.
The residual table, its full collected expansion, and the construction
certificate accompany the manuscript. No huge constant is treated as
an additional free parameter. In particular, all twenty-five variables
in \eqref{eq:o8-chain} belong to the stated witness count.

The recursive module counts are
\[
\begin{array}{c|rr}
 \text{module}&\text{new variables}&\text{residuals}\\\hline
 \text{pole}&7&4\\
 \text{anchor}&16&6\\
 \text{scalar batch, including pole}&34&31\\
 \text{forty-one-test batch, including pole}&58&55\\
 \text{one joint comparison, including batch}&373&360\\
 \text{root, including all its calls}&15646&15082\\
 \text{shared anchor chain}&25&25
\end{array}
\]
The root has $303$ direct variables and $285$ direct residuals; its
calls are one anchor, one scalar batch, and forty-one comparisons.
Thus $303+16+34+41\cdot373=15646$ and
$285+6+31+41\cdot360=15082$. Expansion of \eqref{eq:o8-polynomial}
gives the finite assertions of Proposition~\ref{prop:o8-finite}.
Over $\Q$, a sum of squares vanishes exactly when every summand
vanishes, so the residual interpretation proves the theorem.
This computation checks the polynomial and its compilation. The proof
of the height estimate is the arithmetic argument given above.

% End degree-eight construction
% Begin explicit quartic constructions
\section{An explicit quartic defining the integers}
\label{sec:quartics}
This is the second final polynomial of the manuscript. The original huge
polynomial remains unchanged. The base residual compiler in
Section~\ref{sec:degree8-original} supplies construction machinery and
inherited arithmetic inputs; it is not a third final polynomial.

\begin{theorem}[An explicit quartic definition]\label{thm:quartic}
With precisely the inherited arithmetic hypotheses and qualifications of
Section~\ref{sec:degree8-original} and Appendix~\ref{app:sources}, the
coefficient file below defines one integer polynomial $\Phi$ of total
degree $4$, with $72\,925$ nonzero monomials in $4\,973$ variables.
For every $q\in\Q$,
\[
 q\in\Z\quad\Longleftrightarrow\quad
 \exists(w_1,\ldots,w_{4972})\in\Q^{4972},\qquad
 \Phi(q,w_1,\ldots,w_{4972})=0.
\]
There are $4\,972$ existential witnesses besides the original input.
The coefficient bound is
\[
 \|\Phi\|_\infty=685783949153203534233890<2^{80}.
\]
\end{theorem}

\subsection{Literal specification and coefficient file}
The complete coefficient table is \nolinkurl{quartics/quartic-4973/polynomial.expanded.txt}.
After one comment line, each line gives a nonzero signed integer coefficient
followed by a nondecreasing list of variable indices. Repeated indices give
their multiplicities as exponents; a coefficient alone is the constant.
Thus \texttt{-6 0 2 2} denotes $-6X_0X_2^2$. The original input is $X_0$.
Like terms are collected, and the rows are in lexicographic index-list order.
The occurrence-frequency relabeling is recorded in
\nolinkurl{quartics/quartic-4973/variable-map.json} and fixes index zero.
The gzip companion has exactly the same uncompressed bytes.
The coefficient table has $1\,622\,053$ bytes, including
its comment and line endings, and SHA-256
\begin{center}\footnotesize\ttfamily
094f6a7ed17e7cc8d5d17a2f809127bbaffd0fe343cb2a32d18ac97e61279876
\end{center}

This specifies every coefficient without an auxiliary choice of witnesses.

\subsection{Arithmetic source and retained foundations}
\label{subsec:quartic-source}
The selected source uses the following retained height, contact and
reciprocal-logarithm proofs before the new source equivalence is established.
An element of $K$ is represented by two rational coefficients, with
\[
 (a_0,a_1)(b_0,b_1)=(a_0b_0+2a_1b_1,a_0b_1+a_1b_0),\qquad
 a^{-1}=\frac{(a_0,-a_1)}{a_0^2-2a_1^2}\quad(a\ne0).
\]
The inverse denominator is nonzero by the irrationality of $\sqrt2$.

% BLOCK HEIGHT19.
\begin{theorem}[A sharper five-contact height estimate]\label{thm:quartic-height19}
Keep $S_*$, $\mathcal D$ and $M_5$ from Theorem~\ref{hyp:o8-height}, and put
\[
 B_*=\prod_{p\le76176}p,\qquad
 A_{19}=2^{109}210^{1800},\qquad H_{19}=2A_{19}B_*^{19}+128.
\]
For every $a\in\mathbb Q\setminus\mathcal D$,
\[
 h(a)\le H_{19}M_5(a)^{19},
 \qquad \log_2H_{19}<2093966.
\]
Here $h$ is the absolute logarithmic Weil height.
\end{theorem}
\begin{proof}
We use the geometric construction and compatible representations in the
proof of Theorem~\ref{hyp:o8-height}, not merely its final exponent-$24$
inequality. Write $R=B_*M_5(a)$ and $C=210^{300}$.
The geometric comparison there is
\[
 h(a)\le2\max(0,h_F(A))+128.
\]
All $h_F$ in this proof use the original stable Faltings metric.
The CM estimate in that proof is $h_F(A)<854R$.

In the non-CM case the same construction gives an absolutely simple
quaternionic surface $A/L$, a normal field $L'/\mathbb Q$, and the final
primitive weight-two newform $f$, after the finite twist has been
absorbed, with
\[
 [L:\mathbb Q]\le576,\quad d'=[L':\mathbb Q]\le2304R,\quad
 N\le CR^2,\quad \operatorname{rad}(N)\mid R.
\]
Its two-dimensional compatible representation $r_\ell$ satisfies
\[
 V_\ell(A)|_{G_{L'}}\simeq
       (r_\ell\oplus r_\ell)|_{G_{L'}},\qquad
 (\det r_\ell)^2|_{G_{L'}}=\chi_{\rm cyc}^2|_{G_{L'}}.
\]
These are (LR3)--(LR5) and the determinant calculation in the existing
modularity argument. There is no bounded-degree hypothesis on the
coefficient field of $f$.

Let $\epsilon$ be the nebentypus of this final form. In the Tate-module
convention, with arithmetic Frobenius, $\det r_\ell=\epsilon\chi_{\rm cyc}$;
see \cite[\S12.5, (12.5.1)--(12.5.2) and Remark 12.5.1,
pp.~119--120]{diamond-im}. Equivalently, at a good prime its
characteristic polynomial is $X^2-a_pX+p\epsilon(p)$.
A consistent contragredient convention replaces the finite character by
its inverse, with the same kernel and order. The character $\epsilon^2$
is trivial on $G_{L'}$, so its image order divides $d'$. Consequently
\[
 t=\operatorname{ord}(\epsilon)\le2d'\le4608R.
\]
This bounds the character order, not the coefficient-field degree.

Put $H=\ker\epsilon\subset(\mathbb Z/N\mathbb Z)^*$ and
\[
 \Gamma_H=\left\{
 \begin{pmatrix}a&b\\c&d\end{pmatrix}\in\Gamma_0(N):
 d\bmod N\in H\right\}.
\]
Since $\epsilon(-1)=1$, this group contains $-I$.
Let $X_H=X_1(N)/H$ for the diamond action and $J_H=\operatorname{Jac}(X_H)$.
The curve, diamond action and quotient are defined over $\mathbb Q$.
The rational model here parametrizes a point of exact order $N$, up to
$H$, and is the model specified by \cite[\S2.1, (2.1.8)]{mascot-moduli};
it introduces no cyclotomic field-degree factor.

Write $q:J_1(N)\to A_f$ for the newform quotient. Its pulled-back
differentials are spanned by all Galois conjugates of $f$; see
\cite[Theorem~1(i),(iv), p.~526]{shimura-factors} and
\cite[\S3.1]{baker-et-al-finiteness}. Every $h\in H$ fixes those differentials,
since its eigenvalue on each conjugate is a conjugate of $\epsilon(h)=1$;
the diamond convention is \cite[\S3.2.1]{baker-et-al-finiteness}.
Homomorphisms of abelian varieties in characteristic zero are determined
by their differentials. Hence $q\circ\langle h\rangle=q$.
For
\[
 e_H=\frac1{|H|}\sum_{h\in H}\langle h\rangle
       \in\operatorname{End}_{\mathbb Q}^0(J_1(N))
\]
we therefore have $qe_H=q$.

Let $H_{\rm eff}$ be the image of $H$ in
$\operatorname{Aut}(X_1(N))$, and let $\pi:X_1(N)\to X_H$ be the quotient.
On Jacobians,
\[
 \pi^*\pi_*=\sum_{h\in H_{\rm eff}}h,\qquad
 \pi_*\pi^*=[|H_{\rm eff}|].
\]
The averages over $H$ and $H_{\rm eff}$ agree. In particular the trivial
action of $-1$ is not counted twice in the first identity.
The image of $e_H$ is thus $\mathbb Q$-isogenous to $J_H$, and its
restriction of $q$ is surjective. Poincare reducibility makes $A_f$
a $\mathbb Q$-isogeny factor of $J_H$.

The accepted realization of $r_\ell$ in $A_f$ puts it in $V_\ell(J_H)$.
Over $L'$, Faltings' Hom theorem and extension of scalars give a nonzero
homomorphism $A\to J_H$. Absolute simplicity and Poincare reducibility
give an $L'$-isogeny $A\times Z\to J_H$. In particular
$g=\dim J_H\ge2$. The Jacobian has its canonical principal polarization;
no bound on a polarization degree of $A$ or $Z$ is being assumed.

The map $X_H\to X(1)$ has degree
\[
 d_H=[\operatorname{SL}_2(\mathbb Z):\Gamma_H]
     =Nt\prod_{p\mid N}(1+1/p).
\]
There is no factor of two here, because $-I\in\Gamma_H$.
It is branched only over the cusp and the two elliptic points, so
$\deg_B(X_H)\le d_H$. The congruence-curve genus bound gives
$g\le1+d_H/12\le d_H$, using $g\ge2$.
These degree and genus conventions, and the resulting height inequality,
are recorded in \cite[\S4, (4.1)--(4.2), p.~8]{vonKaenel2021}.
For every prime $p\ge5$, $(1+1/p)^6\le p$, while the possible factors
at two and three have product at most two. It follows that
\[
 d_H\le9216CR^{19/6}<T,\qquad
 T=2^{14}CR^{19/6},\qquad g<T.
\]
By \cite[Theorem~1.1.1 and \S2.3]{javanpeykar},
\[
 h_F(J_H)\le13000000\,d_H^5g<2^{24}T^6.
\]
This is the generic Belyi-curve estimate; no squarefree-level restriction
or trivial-nebentypus assumption enters it.

Apply \cite[Theorem~1.9(1), pp.~7--8]{GaudronRemond2023} with source
$(J_H)_{L'}$ and target $A\times Z$. Combining its degree bound with
half-logarithmic isogeny-height variation gives
\[
 h_F(A)+h_F(Z)\le h_F(J_H)+V,\quad
 V=2g^2\left(8g^2\log(7g)+\log d'
     +\log\max(1,\log d',h_F(J_H))\right).
\]
The Bost lower bound in \cite[(2.2), p.~5]{vonKaenel2021} gives
$h_F(Z)>-3g/2$. Put $u=\log T$. Then
\[
 V\le16T^4(2+u)+2T^2(25+7u)\le T^6.
\]
Indeed $\log d'<8+\log R\le8+u$, and the last logarithm in $V$ is
less than $17+6u$. For $T\ge16$ the inequalities
$2+\log T\le T^2/32$ and $25+7\log T\le T^4/4$ follow at $16$
and remain true by differentiation. Thus
\[
 h_F(A)<(2^{24}+1)T^6+\tfrac32T
       <2^{25}T^6=A_{19}R^{19}.
\]
The earlier CM estimate is also less than this bound.
The geometric comparison and $M_5(a)\ge1$ prove the assertion.
Finally $B_*<2^{109445}$, $210<2^8$ and $A_{19}B_*^{19}>64$ give
\[
 H_{19}<4A_{19}B_*^{19}
       <2^{\,2+109+1800\cdot8+19\cdot109445}
       =2^{2093966}.
\]
\end{proof}

% BLOCK ROOT37: no denominator-clearing variable is added to this source.
\begin{lemma}[Thirty-seven contacts for rational coefficients]
\label{lem:quartic-rational37}
Keep $K=\Q(\sqrt2)$, $d=3-\sqrt2$, $D=d^2$, and
$E:y^2=x^3-Dx$ from Section~\ref{sec:degree8-original}.
Let $A,S,T,Z$ be nontorsion points of $E(K)$.
Order the eight rational coefficients of
$x(A),x(S),x(T),x(Z)$ as $a_1,\ldots,a_8$. Put
\[
 R_{\rm c}=47027,\quad W_0=25878772920,\quad c=18484837801,
 \quad V_0=32W_0,
\]
\[
 A_i=R_{\rm c}a_i+3,\quad w_i=A_iv-u.
\]
Use the following thirty-seven nonzero raw contacts:
\begin{equation}\label{eq:quartic-contacts37}
 r_0=v/V_0,\qquad u-\delta v\quad(\delta\in\mathcal D),
 \qquad -(w_i-\delta v)\quad(1\le i\le8,\ \delta\in\mathcal D).
\end{equation}
Here and below the order within $\mathcal D$ is
$(0,1,2401,-343/32)$. The common five positions are numbered $0,\ldots,4$;
the other positions are numbered $5,\ldots,36$ in the displayed order.

Every normalized odd contact of $u/v$ or $w_i/v$ outside $S_*$ is
detected by an odd valuation at a raw position belonging to that slope.
If the radicals for both slopes $u/v,w_i/v$ have their $e$th powers
dividing an integer $N\ne0$, then, with $H=H_{19}$,
\[
 h(a_i)<2H|N|^{19/e}+14.
\]
Consequently the two coefficients of any one of the four abscissas give
\begin{equation}\label{eq:quartic-direct-height}
 h(x(Q))<4H|N|^{19/e}+30<7H|N|^{19/e}
 \quad(Q\in\{A,S,T,Z\})
\end{equation}
whenever all its required positions have those divisibilities.

Conversely, for any fixed eight rational coefficients, the inherited
ordinary-prime construction supplies these raw contacts with every finite
contact a positive selected prime and $r_0$ a positive rational square
times a selected prime. Repeated positions may have equal values.
These contacts have the positivity and local squareclasses used by Lemmas~\ref{lem:quartic-conic-completeness} and \ref{lem:quartic-linear-pole}.
\end{lemma}
\begin{proof}
Nonzero raw contacts imply $v\ne0$ and exclude the four finite branch
values for every slope. At a prime outside $S_*$ all differences of
distinct members of $\mathcal D$ are units. Write $m$ for the minimum
of the valuations of the two coordinates of one slope. If an odd finite
normalized contact has $m$ even, its distinguished raw linear form has
odd valuation. If $m$ is odd, another finite form has valuation $m$.
For an odd contact at infinity use the denominator in the first case
and a finite form in the second. The unit $V_0$ and the signs do not
change parity. This proof uses rational coordinates and also covers
negative valuations; it does not presuppose that a contact is prime.

Apply Theorem~\ref{thm:quartic-height19} to the two slopes. The identity
$R_{\rm c}a_i+3=u/v+w_i/v$ and the elementary logarithmic height inequalities
give the first bound, since $\log(6R_{\rm c})<14$. For
$x(Q)=a+b\sqrt2$ use
$h(x(Q))\le h(a)+h(b)+h(\sqrt2)+\log2$.
This proves \eqref{eq:quartic-direct-height}, since $H>128$ and $|N|\ge1$.

For the converse choose a positive integer $D_0$ with $D_0a_i\in\Z$
for every $i$, after fixing the eight coefficients. This integer is a
choice in the existence proof, not a variable of the source. Set
\[
 T_0=V_0D_0^2,\qquad p_0=c+W_0y,\qquad
 u=c+W_0x,\qquad v=T_0p_0.
\]
For each distinct slope
$b\in\mathcal D\cup\bigcup_i(A_i-\mathcal D)$ put
$K_b=T_0b\in W_0\Z$ and $\Psi_b=u-K_bp_0$.
Together with $p_0$, these are at most thirty-seven integer affine forms,
each congruent to $c$ modulo $W_0$. Distinct slopes give independent
linear parts. Exact duplicate slopes are removed only while applying
the prime theorem and restored at all their original positions afterward.

The value and successor admissibility in
Section~\ref{subsec:o8-pattern-proof} remains valid. At primes dividing
$W_0$ use the same residue $c$. If a prime $\ell>29$ divides $D_0$,
all $K_b$ vanish modulo $\ell$, and $u=p_0=1$ avoids zero and minus one.
If $\ell>29$ does not divide $D_0$, the change from $(u,p_0)$ to $(u,v)$
is invertible. For $\ell=31,37,41$ there are at most eight slope residues,
since $R_{\rm c}a_i\equiv0$; for $\ell\ge43$ there are at most thirty-six.
Choose distinct missing residues $b_0,e_0$ and order them so that
$e_0-b_0\ne-T_0^{-1}$. Then
\[
 v=(e_0-b_0)^{-1},\qquad u=b_0v
\]
make every $\Psi_b$ different from zero and minus one, and make
$p_0=v/T_0$ different from zero and minus one. Interchanging the two
missing residues permits the required ordering, because $\ell$ is odd.

Choose a fixed positive cone
$u/v>\max_i|A_i|+2413$, $v>0$. Enlarge the fixed initial cutoff past
$82$, every prime dividing $D_0$, and the nonzero determinants of the
distinct forms. The inherited prime-pattern proof then applies with its
same unused third coordinate and with complexity at most $35$.
For its fiber upper estimates, fixing $\Psi_{b_0}=p$ gives the same
one-variable forms
$\Psi_b=p+(K_{b_0}-K_b)p_0$.
All coefficients, the cone, the denominator $D_0$ and the cutoff are
fixed before the box grows. Thus the proof uses precisely the inherited
fixed-system prime and successor estimates, not a uniform estimate in a
growing denominator. It gives the positive ordinary primes and the
same-prime elliptic-order conditions used there.

Now $r_0=D_0^2p_0$. At a prime $\ell\mid D_0$ outside $S_*$, both $u$
and $p_0$ are units, $v_\ell(v)=2v_\ell(D_0)$, and
$v_\ell(A_i)\ge-v_\ell(D_0)$. Hence $w_i=A_iv-u$ is a unit.
Both slopes have negative valuation $-2v_\ell(D_0)$, and their normalized
infinite depths are even. No new normalized odd contact is created at
such a denominator prime. At $p_0$, the odd valuation of $r_0$ is one.

For each occurrence choose its private comparison witnesses sufficiently
deep that its norm root is a unit at its assigned prime. The remaining
denominator valuations in the pole argument occur in the square of that
root. The raw contacts are positive, their reduced denominators are
squares, and they are squares in $\Q_2,\Q_3,\Q_{17}$ by the retained
residue $c$ and the square factor $D_0^2$.
Lemmas~\ref{lem:quartic-conic-completeness} and
\ref{lem:quartic-linear-pole} supply the selected ordinary-pole witnesses. Equal prime occurrences use the
same orientation and may use separate private points. No condition on
reduction orders at different primes has been added.
\end{proof}

% BLOCK RECIPROCAL: the released complete proofs with the stated uniform range.
\begin{lemma}[Reciprocal logarithm and its error]
\label{lem:quartic-reciprocal-log}
\label{lem:quartic-uniform-reciprocal}
Put $f=1/x$, $z=-x/y$, and define
\[
 a_0=1,\qquad a_1=-\frac25,\qquad
 a_j=\frac{4j-5}{4j+1}a_{j-1}\quad(j\ge2),\qquad
 G(w)=\sum_{j\ge0}a_jw^j.
\]
On the formal branch at the identity,
\[
 \log_E(z)=zG(Df^2).
\]
For every $J\ge0$ put $G_J(w)=\sum_{j=0}^{J}a_jw^j$.
For every integer $10\le e\le64$, set $m=\lceil e/4\rceil$ and
\[
 G_{m-1}(w)=\sum_{j=0}^{m-1}a_jw^j,\qquad
 L_e(z,f)=zG_{m-1}(Df^2).
\]
Let $p>76176$ be split and choose either place above it. If the point is
nonidentity and lies in the formal group at that place, so
$r=v_p(z)\ge1$ and $v_p(f)>0$, then
\[
 v_p\bigl(\log_E(z)-L_e(z,f)\bigr)\ge(4m+1)r,
 \qquad \frac{L_e(z,f)}{\log_E(z)}\in1+p^e\Z_p.
\]
\end{lemma}
\begin{proof}
The curve equation gives
\[
 f=z^2(1-Df^2),\qquad w=Df^2,\qquad Dz^4=\frac{w}{(1-w)^2}.
\]
Coefficient comparison in the defining recurrence gives
\[
 4w(1-w)G'(w)+(1+w)G(w)=1-w.
\]
Indeed the coefficient of $w$ is $5a_1+a_0=-1$, and each later
coefficient gives $(4j+1)a_j=(4j-5)a_{j-1}$. Differentiating the
curve identity yields $z\,dw/dz=4w(1-w)/(1+w)$. Consequently
\[
 \frac{d}{dz}\{zG(w)\}=\frac{1-w}{1+w}
                   =(1+4Dz^4)^{-1/2}.
\]
The square root has constant term one. This is the normalized invariant
differential for $z$, and $zG(w)$ has zero constant term, which proves
the formal identity.

At the stated place, $D$ is a unit and the formal branch gives
$v_p(f)=2r$, hence $v_p(w)=4r$. The denominator of $a_j$, for $j\ge1$,
divides $5\prod_{k=2}^j(4k+1)$ and therefore divides $(4j+1)!$.
Thus
\[
 v_p(a_j)\ge-v_p((4j+1)!)\ge-\frac{4j+1}{p-1}.
\]
For $j\ge m$ with $4j+1<p$, the coefficient is $p$-integral and
$v_p(a_jw^j)\ge4mr$. Otherwise $j\ge19044$, so $3j\ge4m$ and
$(4j+1)/(p-1)\le j$. Therefore
\[
 v_p(a_jw^j)\ge4jr-j\ge4mr.
\]
These lower bounds tend to infinity. They control the entire infinite
tail, including terms whose denominators are divisible by $p$.
Multiplication by $z$ proves the first estimate. Formal logarithm isometry
gives $v_p(\log_E(z))=r$, and proves the relative estimate.
\end{proof}

\begin{lemma}[Global nonvanishing]
\label{lem:quartic-reciprocal-nonzero}
For every $f\in K$ and every $J\ge2$,
\[
                 G_J(Df^2)\ne0.
\]
In particular $L_e(z,f)\ne0$ whenever $z\ne0$ and $10\le e\le64$.
\end{lemma}
\begin{proof}
Let $\nu_2$ be the rational exponent-of-two valuation and put
\[
 v(a+b\sqrt2)=\tfrac12\nu_2(a^2-2b^2)
             =\min\{\nu_2(a),\nu_2(b)+\tfrac12\}.
\]
The two terms in the norm have valuations of opposite parity when finite.
Norm multiplicativity and the displayed minimum prove the valuation
identities. Its finite values lie in $\tfrac12\Z$, and $v(D)=0$.
Every $a_j$ with $j\ge1$ has valuation one, because $a_1=-2/5$ and
all subsequent recurrence factors have odd numerator and denominator.
If $f=0$, the assertion is immediate. If $t=v(f)\ge0$, the constant
term of $G_J(Df^2)$ is uniquely least in valuation. If $t<0$, then
$t\le-1/2$ and the nonconstant valuations $1+2jt$ strictly decrease;
the last is at most $1-J<0$, below the constant term as well. In either
case a uniquely least term proves nonvanishing. No good-reduction
assumption at two is used.
\end{proof}

\subsection{Conventions and inherited hypotheses}
\label{sec:v4-inherited}

Retain the field and curve of Section~\ref{sec:degree8-original}:
\[
 K=\Q(\sqrt2),\qquad d=3-\sqrt2,\qquad D=d^2,
 \qquad E:\ y^2=x^3-Dx.
\]
For an affine point with nonzero abscissa write $f(P)=1/x(P)$ and
$p(P)=y(P)/x(P)$. The rational function $z=-x/y$ extends with value zero
at $O$ and $T_0=(0,0)$; its other poles are $(d,0)$ and $(-d,0)$.
The letter $J$ denotes a common comparison point. We write $\jmath$ for
$1/z(J)$, to distinguish that field element from a point or an integer index.
The implementation conjugates the complete alternate-comparison identity:
its boundary value is $1/\overline{z(J)}$, and its primary and target
coefficients in that identity are conjugated as well.

The following arguments retain all inherited qualifications of
Section~\ref{sec:degree8-original} and Appendix~\ref{app:sources}. In particular, the restricted converse, rank and original
arithmetic assertions retain their stated proof status. The prime-pattern
input retains its corrected author-version and erratum conventions,
including the stated version of arXiv:1009.3998v5. The Pan interface retains
the inspected author theorem and the cited Calegari formulation, without
claiming comparison with inaccessible journal text. The Yuan input retains
its author-attested translation qualification. The stable Faltings metric,
arithmetic-Frobenius, diamond-quotient and final-twist conventions are
unchanged. None of the finite source or polynomial checks strengthens these
assumptions.

In particular we use, in exactly their previous scopes, the generator and
torsion statements $2E(K)=\Z P_*$ and $E(K)_{\mathrm{tors}}=E[2](K)$;
the formal-group logarithm and valuation statements; the fixed-system
prime-pattern theorem with arbitrary prescribed finite exclusions; the
height, scalar, anchor and Lucas lemmas; and the restricted-converse and
simultaneous-realization inputs used for the pole predicate. The one
additional explicit trace application below is identified in
Lemma~\ref{lem:quartic-anchor-four}.

Put
\[
 B_* =\prod_{p\le76176}p,\qquad
 H=2(2^{109}210^{1800})B_*^{19}+128,
 \qquad \log_2 H<2093966.
\]
This is $H_{19}$ of Theorem~\ref{thm:quartic-height19}. The
source parameters remain high order $40$, beta order $12$, scalar orders
$(1,1,1,1,32,32)$ and anchor constant $C=3^{2^{25}}$.

We also retain Lemma~\ref{lem:quartic-reciprocal-log}. Explicitly,
write
\[
 \sqrt{1-w}\sum_{j\ge0}\frac{\binom{2j}{j}}{4^j(4j+1)}w^j
       =\sum_{j\ge0}a_jw^j,\qquad
 G_m(w)=\sum_{j=0}^m a_jw^j.
\]
The polynomials used here are $G_9$ and $G_2$. They are nonzero at every
argument $Df^2$ with $f\in K$, by the previously proved dyadic argument.
For a nonidentity formal-group point at either place above a split prime
$p>76176$, the approximation $-G_9(Df^2)/p(P)$ has relative order at least
$40$, and $-G_2(Df^2)/p(P)$ has relative order at least $12$. Actual formal
membership is required: small $z$ alone is not the premise of this lemma.

\subsubsection{The retained scalar, anchor and beta interfaces}
The scalar input is still $\beta=X^2+2$. The scalar batch tests the two
rational coefficients each of $f(J_0),f(S),z(S)/z(J_0)$, with orders
$(1,1,1,1,32,32)$ in that order. Its original pole predicate is retained.
The anchor constant is computed by
\[
 c_0=3,\qquad c_{i+1}=c_i^2\quad(0\le i<25),\qquad C=c_{25}.
\]
Thus the integer $C=3^{2^{25}}$ is a witness value, not a polynomial
coefficient. The four-witness anchor below consumes this same terminal
value. The height targets are the eight rational abscissa coefficients
of Lemma~\ref{lem:quartic-rational37}.

For every beta-bearing comparison, the common and beta quotient equations
are explicitly
\[
 \Lambda p_H=G_H,\qquad B_A=C_Ap_A,\qquad B_A\Lambda=G_A.
\]
The global nonvanishing in Lemma~\ref{lem:quartic-reciprocal-nonzero}
forces $\Lambda,p_H,B_A,p_A\ne0$. Hence $C_A$ is the ratio of the beta
and common approximate logarithms, with their two minus signs cancelled.
The beta residual is $C_A-2\beta$. In full comparisons $G_H$ has order
$40$ and $G_A$ order $12$; in beta-only comparisons both have order $12$.
The Lucas numerators use the normalized equations below, whose permitted
zero-ordinate branches must not be excluded by this beta argument.

\subsection{Selective comparisons and the local elliptic chart}

\begin{lemma}[Three-mode coverage]\label{lem:quartic-three-mode-coverage}
Number the unchanged raw contacts from $0$ to $36$. Retain full comparisons
at $0,\ldots,4$ and $21,\ldots,28$, beta-only comparisons at
$5,\ldots,12$, and Lucas-only comparisons at $13,\ldots,20$ and
$29,\ldots,36$. The contacts carrying Lucas information cover the three
points $S,T,Z$, and those carrying beta information cover $A,T$.
\end{lemma}
\begin{proof}
The first five positions are shared. Each of $A,S,T,Z$ has its original
block of eight finite positions, in that order. Thus the Lucas set contains
the five shared positions and all three finite blocks of $S,T,Z$; the beta
set contains the five shared positions and the two finite blocks of $A,T$.
The unchanged two-slope rational-height argument applies separately to
these sets. No position is deleted when two slopes, raw values, or assigned
primes coincide. The two maxima used below are consequently
$B_L=\max(|s|,|t|,|z|)$ and $B_{AT}=\max(|n|,|t|)$.
\end{proof}

\begin{lemma}[Doubled formal comparisons]
\label{lem:quartic-doubled-private-comparison}
At a split good prime above the cutoff, if $v(z(R))>0$, then either $R$ or
$R+T_0$ is in the formal group. Whenever the relevant rational function is
defined,
\[
 v(\LogE([2]R))=v(z(R)),
\]
with infinite valuation for $R=O$ or $T_0$.
\end{lemma}
\begin{proof}
For an affine nontorsion point,
$z(R)^2=x(R)/(x(R)^2-D)$. If $v(x)=0$, the right side has nonpositive
valuation. If $v(x)<0$, the point is formal. If $v(x)>0$, translation by
$T_0$ gives $x'=-D/x$ and $z'=-z$, so the translated point is formal with
the same parameter valuation. Doubling kills $T_0$ and is an isometry on
the formal group at an odd good prime. The asserted equality follows from
the retained logarithm isometry. The two other nonzero two-torsion points
are poles of $z$ and are not in the domain.
\end{proof}

Private primary points may therefore lie in $E(K)$ rather than $2E(K)$.
If $[2]P=[2\alpha(P)]P_*$, their coefficients $\alpha(P)$ may be
half-integral. They are integral at each consumed odd prime. One compares
logarithms only after writing an actual integer-multiple identity between
formal points; torsion is not discarded by equating these half-indices.
The primary $f$ tests still put the primary points themselves in the
formal group. An alternate test of
$z(P-Q)/z(J)$, together with the opposite-place $f(J)$ test, implies
$v(z(P-Q))>0$, and the lemma supplies the needed doubled logarithm.

\begin{lemma}[Joint primary and difference graph]
\label{lem:quartic-joint-chart}
Let $Q=(a,aq)$ be nontorsion and $\jmath\ne0$. Set
\[
 F_Q=1/a,\quad c=a+DF_Q,\quad
 r=\jmath/q,\quad s_Q=c/q,\quad t_Q=1/q^2,\quad u=rs_Q.
\]
The three field equations
\begin{align}
 k&=(A-r)p,\\
 kp+up-2Drf-4Dt_QA&=0,\\
 2D(A+r)f+s_Qk+4Dt_QA+\jmath q&=0
\end{align}
give the exact graph $f=1/x(P)$, $p=y(P)/x(P)$,
$A=\jmath z(P-Q)$ for affine primary points $P$ with $x(P)\ne0$
and finite $z(P-Q)$. In particular this primary chart excludes $O$ and
$T_0$. The equations permit $A=0$ and do not
require inverses of $A,k$ or of a difference parameter.
\end{lemma}
\begin{proof}
Set $x=p^2+Df$. The target identity is $c^2-q^4=4D$.
Reversible linear combinations of the displayed equations give
\[
 A U=\jmath H_2,\qquad A V=\jmath H_1,
\]
where
\[
 U=q(x+Df)+cp,\quad V=q^2(x-Df)-4D,\quad
 H_1=q(x+Df)-cp,\quad H_2=x-Df-q^2.
\]
Direct expansion yields
\[
 H_2V-H_1U=-4Dq^2(fp^2+Df^2-1).
\]
Thus $xf=1$ and $P=(x,xp)$ is on $E$. If $U=V=0$, then for
$v\in\{d,-d\}$ one has
$p=-2v/q$ and $x=v(a+v)/(a-v)$, precisely $P=Q+(v,0)$.
At either point $H_2=(4D-q^4)/q^2\ne0$: equality $q^4=4D$ would imply
$N_{K/\Q}(q)^2=28$, impossible for a rational norm. These genuine poles
cannot satisfy the graph equations. Else $U$ or $V$ is nonzero and
$A/\jmath$ is uniquely determined. The chord identities
\[
 H U=N H_2,\qquad H V=N H_1,
\quad H=(p-q)(1+DfF_Q),\quad
 N=pq(1-DfF_Q)+2D(f+F_Q)
\]
identify it with $z(P-Q)$ on the chord domain and hence at every regular
point. At $P=Q$ and $Q+T_0$ the value is zero. Opposition is covered by
$V\ne0$. At $p=0$, the curve has $x=\pm d$, $U=2qx\ne0$, and the graph
still gives its unique regular value. Conversely substituting any point
and its regular difference value reverses these identities and recovers
all three equations with $k=(A-r)p$.
\end{proof}

The last $p=0$ cases are needed below. They are excluded at consumed places
by the positive primary-$f$ valuation, but need not be excluded globally
for every normalized numerator.

\subsection{Private conics, packed field tests and the ordinary pole}

Each contact $r$ has its own rational variables $\theta,s_c,S_c$ with
\[
 S_c=s_c^2,\qquad \theta^2-rS_c=2.
\]
In particular $r,s_c\ne0$. For a field value $F_i=a_i+b_i\sqrt2$ of
required order $e_i$, define
\[
 Y_i=a_i+\theta b_i,\qquad
 T_i=2a_iY_i-a_i^2+2b_i^2
     =Y_i^2-(\theta^2-2)b_i^2,\qquad
 r^{2e_i}W_i=T_i.
\]
The upper tree uses successive pairs $N_r(u,v)=u^2+rv^2$, starting with
the literal leaf $1$ followed by the $W_i$ in stable order, carrying an
unpaired final child unchanged. Denote its root by $W$.

\begin{lemma}[Local conic and tree deductions]
\label{lem:quartic-private-conic-packing}
Assume the ordinary pole gives no negative odd valuation to nonzero
$rW^2$ outside its fixed exceptional primes. At any $p>76176$ with odd
$R=v_p(r)$, the equations above imply $R>0$ and choose one root
$\rho\in\Q_p$ of $2$ such that
$v_p(a_i+\rho b_i)\ge e_i$ for every retained field.
\end{lemma}
\begin{proof}
The conic makes $\delta=v_p(\theta^2-2)=R+2v_p(s_c)$ odd.
It cannot be negative, since then $v_p(\theta^2)$ would be odd; it cannot
be zero. Thus $\delta>0$, $\theta$ is a unit, and one root $\rho$ satisfies
$v_p(\theta-\rho)=\delta$ while the other factor is a unit.
The summands in $N_r(u,v)$ have opposite valuation parity and cannot
cancel. The literal leaf proves $W\ne0$. If $R<0$, ascent along that
leaf gives $v_p(W)\le0$, contradicting the pole conclusion. Hence $R>0$.
The pole conclusion gives $v_p(W)>-R/2$. Descent through the tree uses
\[
 v_p(N_r(u,v))>-cR\ \Longrightarrow
 v_p(u)>-cR/2,\quad v_p(v)>-(c+1)R/2,
\]
with $0\le c<1$, and gives $v_p(W_i)>-R$. Consequently
$v_p(T_i)>(2e_i-1)R\ge2e_i-1$, so $v_p(T_i)\ge2e_i$.
The two summands of $T_i$ also have opposite parity, yielding
$v_p(Y_i)\ge e_i$ and $\delta+2v_p(b_i)\ge2e_i$.
It follows that $v_p((\theta-\rho)b_i)\ge e_i$ and hence the claimed
field order. Zero summands are treated with infinite valuation. No
integrality of $a_i,b_i$ and no global squareclass of $F_i$ was assumed.
\end{proof}

The proof uses $v_p(T_i)\ge2e_i$, not the generally unjustified stronger
bound $2e_iR$. It is independent of the elliptic meaning of the fields.

\begin{lemma}[Conic completeness]\label{lem:quartic-conic-completeness}
In the retained integral-input construction, each contact
$r=d_0^2p_0>0$ admits the needed private conic and arbitrarily deep
approximation to its chosen split branch, independently of the other
contact occurrences.
\end{lemma}
\begin{proof}
Here $p_0\equiv1\pmod8$ and $d_0\ne0$ is a $p_0$-unit. The ring
$\Z[\sqrt2]$ is norm Euclidean: subtracting nearest integers from two
coefficients leaves $|a|,|b|\le1/2$ and $|a^2-2b^2|\le1/2<1$.
The split prime ideal therefore gives integers $a,b$ with $a^2-2b^2=p_0$,
after multiplying by the norm-minus-one unit if needed. Since $b\ne0$,
$(\theta_0,s_0)=(a/b,1/(d_0b))$ is a rational conic point. The parametrization
\[
 \Delta=\frac{2(rs_0-\theta_0 t)}{t^2-r},\qquad
 s_c=s_0+\Delta,\quad \theta=\theta_0+t\Delta
\]
and rational approximation give simultaneous approximation at the finite
set of required places. There is no rational pole $t^2=r$.
Near the selected root, $\delta=1+2v_{p_0}(s_c)$ can be made arbitrarily
large. After all field coefficients are fixed, choose it so that
$\delta+2v_{p_0}(b_i)\ge2e_i$ for every $i$. The original field orders
then make all $W_i$ integral; the constant path makes $W$ a $p_0$-unit.
At the real place the tree is positive, even if some $W_i$ are negative.
The extra approximations used for the following pole may be imposed at
the same time. No conic variable is shared between occurrences.
\end{proof}

\begin{proposition}[Linear-unit ordinary pole]
\label{lem:quartic-linear-pole}
Set $g=3-2\theta$. Replace only the ordinary pole module by
\begin{gather*}
 \ell=gW,\quad z_\ell=\ell^2,\quad c_r=re,\qquad
 p_r^2-17r-34c_r+51ec_r=0,\\
 w=17e-c_rz_\ell,\quad j_s=s_r^2,\quad k_r=rq_r,\\
 j_s^2+34j_sw-867w^2-17q_rk_r=0.
\end{gather*}
It has the one-sided pole conclusion used in
Lemma~\ref{lem:quartic-private-conic-packing}. It has integral-input
completeness under the same inherited realization and rank hypotheses.
The scalar-height pole module is unchanged.
\end{proposition}
\begin{proof}
At an odd prime with odd $v_p(r)$, the conic argument makes $\theta$
integral, and
$(3-2\theta)(3+2\theta)=1-4rs_c^2$ is a unit. Thus $g$ is a unit.
Suppose $b=rW^2\ne0$ has negative odd valuation outside $2,3,17$.
The first conic $p_r^2=17r(1-e)(1+3e)$ forces $e$ to be a unit: positive,
negative or infinite valuation of $e$ would make the right side have odd
valuation. Since $v_p(\ell)=v_p(W)$, one has
$v_p(w)=v_p(b)<0$, an odd integer. In
\[
 (s_r^2-17w)(s_r^2+51w)=17rq_r^2
\]
the left side is nonzero with even valuation, since the valuations of
$s_r^2$ and $w$ differ in parity. The right side is zero or has odd
valuation. This contradiction proves the asserted exclusion, including
zero ordinate branches of the first conic. If $W=0$, the conclusion has
no nonzero pole input; explicit witnesses are
$\ell=z_\ell=p_r=q_r=k_r=0$, $e=1$, $c_r=r$, $w=17$, $s_r=17$, $j_s=289$.

For completeness retain the generalized fixed-$U$ realization argument
under its original hypotheses. For $b>0$ with square reduced denominator
and square in $\Q_2,\Q_3,\Q_{17}$, it applies to a prescribed $U_0$ when
\[
 U_0^2/(17b)<1\text{ at infinity},\quad
 v_2(U_0^2/(17b))\ge3,\quad
 v_3(U_0^2/(17b)),v_{17}(U_0^2/(17b))\ge1,
\]
and $v_p(U_0)=v_p(b)$ at every other odd-support prime of $b$.
Indeed, put $B=17b$ and $V=B(B-U_0^2)$, and let
$\beta_{\rm sf},\eta_{\rm sf}$ be their positive squarefree representatives.
The displayed conditions make $V$ positive and a square in
$\Q_2,\Q_3,\Q_{17}$. Consequently
\[
 \beta_{\rm sf}>1,\quad \beta_{\rm sf}\equiv17\pmod{24},\quad
 \eta_{\rm sf}\equiv1\pmod{24},\quad 17\nmid\eta_{\rm sf}.
\]
At an odd-support prime of $B$ outside $2,3,17$, the square reduced
denominator makes $v_p(b)$ positive, and the hypothesis on $U_0$ gives
$v_p(U_0^2/B)>0$. Thus $1-U_0^2/B$ is a square unit, so
$\eta_{\rm sf}$ is a square in $\Q_p$ and $p\nmid\eta_{\rm sf}$.
The same conclusion at $17$ follows from the displayed smallness condition.
Conversely, at a prime dividing $\eta_{\rm sf}$ the necessary strict
cancellation in $B-U_0^2$ gives $v_p(B)=2v_p(U_0)$ and makes
$B/U_0^2$ a square unit. Thus $\beta_{\rm sf}$ is a square in $\Q_p$.
In particular $\gcd(\beta_{\rm sf},\eta_{\rm sf})=1$, and both
mutual-square conditions hold. If $U_0=0$, then $V=B^2$ and
$\eta_{\rm sf}=1$, so the converse condition is empty. These are precisely
the hypotheses consumed by the inherited simultaneous-realization proof;
no new rank or converse statement is substituted.

Now set $b=rW^2$ and $U_0=rW^2g$. In the complete construction
$r=d_0^2p_0$ and $W$ is a $p_0$-unit, so $v_{p_0}(b)=1$ and all
other valuations of $b$ are even. Hence its reduced denominator is a
square. The retained local squares of $r$ at $2,3,17$ also make $b$ a
square at those places. After the fields are fixed, approximate
the private conic also near $(\theta,s_c)=(3/2,1/(2\sqrt r))$ at infinity,
$2,3,17$. These local points exist by the original contact construction.
The continuous function $(r/17)W(\theta)^2g(\theta)^2$ vanishes there,
so it satisfies the displayed smallness conditions while the selected
prime still makes $W$ a unit. The local-unit argument gives the required
odd-support valuations of $U_0$. The resulting old three-conic witnesses satisfy
\[
 p_o^2=17b(1-e)(1+3e),\quad h_o=e(17b-U_0^2),\quad
 (s_o^2-17bh_o)(s_o^2+51bh_o)=17bq_o^2.
\]
Here $r,W,b$ are nonzero. Set
\[
 p_r=p_o/W,\quad w=h_o/b,\quad
 s_r=s_o/b,\quad q_r=q_o/(r^2W^3).
\]
These substitutions give $w=e(17-r\ell^2)$ and the final conic in the
displayed normalized system; the other coordinates are its definitions.
Thus the witness conversion uses only stated nonzero quantities.
Equivalently the unnormalized form
has $L=r\ell$, $z_L=r^2z_\ell$, $w_L=r^2w$, $s_L=rs_r$,
$q_L=r^2q_r$, $j_L=r^2j_s$, $k_L=r^2k_r$, with $e,p_r,c_r$ unchanged.
All divisions in the reverse scaling are powers of the conic-proved
nonzero $r$. This also proves that normalization loses no zero branch.
\end{proof}

\begin{lemma}[Four-witness anchor]\label{lem:quartic-anchor-four}
For rational $a,b$, both belong to $\Z_{(2)}\cap\Z_{(3)}$ if and only if
$9+3(a^2+b^2)$ is a norm-one trace for the quaternion algebra $(2,3)$.
Consequently the original two-coefficient anchor is equivalent to
\[
 CH-9C^2-3f_0^2-3f_1^2=0,\qquad
 H^2-8U^2-12V^2+24W_a^2-4C^2=0.
\]
\end{lemma}
\begin{proof}
By Definition 4, Observation 5 and Facts (a), (b), (c), (e) in the proof of
Proposition 6 of \cite{koenigsmann-defining}, the algebra is ramified
exactly at $2,3$, the trace set is unrestricted at the other places, and it
is integral at its two ramified places
$2,3$, contains integral traces with residues $1$ at $2$ and $0$ at $3$,
and is determined by these local trace conditions. These are the explicit
trace facts used here, in addition to the older sum-of-traces criterion.
If $m=\min(v_p(a),v_p(b))<0$, then at $3$ the sum of squares has valuation
$2m$, while at $2$ it has valuation $2m$ for unequal valuations and
$2m+1$ for equal finite valuations. Thus $9+3(a^2+b^2)$ would be
nonintegral, proving one implication, with zeros covered by infinite
valuation. Conversely integral $a,b$ give a trace congruent to zero at
$3$. At $2$ it is odd or belongs to $4\Z_2$. For the latter case put
$x_0=h/2\in2\Z_2$, where $h$ is this trace. If $x_0\equiv2\pmod4$, then
$(x_0^2-1)/3\equiv1\pmod8$ is a square, giving
$x_0^2-3v^2=1$. If $x_0\equiv0\pmod4$, then $1-x_0^2/6$ is a square,
and $u=v=1$, $w^2=1-x_0^2/6$ gives
$x_0^2-2u^2-3v^2+6w^2=1$. The local trace criterion supplies rational
witnesses. Finally divide the two homogeneous equations by $C^2$,
which is allowed because the retained constant $C$ is nonzero. This is
an equality of existential projections, not a rational formula choosing
quaternion witnesses.
\end{proof}

\subsection{Local and normalized Lucas tests}

Write
\[
 \mathcal F(s,t,z)=t^2-s(s-2)tz+(s-2)^2z^2-1.
\]
\begin{lemma}[Local Lucas residual]\label{lem:quartic-local-lucas}
The primary and alternate logarithm comparisons, before any proxy equality
is used, give local integral quotients $C_s,C_t,C_z$ congruent to the
integer root indices modulo $p^{40}$. An order-$40$ test of
$\mathcal F(C_s,C_t,C_z)$ implies $p^{40}\mid\mathcal F(s,t,z)$.
\end{lemma}
\begin{proof}
The implication follows by substitution into the polynomial with integer
coefficients. One may evaluate the field residual using
\[
 V=(C_s-2)C_z,\quad D_L=C_sC_t-V,\quad
 E_L=C_t^2-VD_L-1.
\]
These are quadratic equations, and $E_L=\mathcal F(C_s,C_t,C_z)$.
They replace the three proxy comparisons and the global proxy equation.
They do not impose $E_L=0$ as a field equality. Removal of the proxy
coordinates requires their exclusive incidence in that equation and those
comparison arguments. Integer completeness and the reverse input
implication use the two height contradictions below; no conversion of
arbitrary old proxy witnesses is asserted.
\end{proof}

\begin{lemma}[Normalized local Lucas evaluation]
\label{lem:quartic-normalized-lucas}
Let $\Lambda p_H=G_H$ be the common negative-logarithm equation. For the
three Lucas numerators replace their quotient blocks and the preceding
evaluation by
\begin{gather*}
 B_i\Lambda=G_i\quad(i=s,t,z),\qquad P B_s=p_s,\\
 (U+V)p_t=P B_t,\qquad K_Lp_t=(1-2P)B_t,\\
 Vp_z=(1-2P)B_z,\qquad E_L=U^2-K_LV-P^2.
\end{gather*}
An order-$40$ test of $E_L$ has the same required integer consequence.
Every integer input still admits the changed witnesses.
\end{lemma}
\begin{proof}
The global nonvanishing of the retained $G_i$ implies
$\Lambda,p_H,B_s,B_t,B_z\ne0$. Also $p_t\ne0$: otherwise the two
equations with left factor $p_t$ would force both $P=0$ and $1-2P=0$.
At a consumed place all primary points are formal, so their $p_i$ are
nonzero. Put $C_i=B_i/p_i$. Then $P=1/C_s$ and direct substitution gives
\[
 E_L=P^2\mathcal F(C_s,C_t,C_z).
\]
Since $C_s$ is locally integral, the implication
$v(E_L)\ge40\Rightarrow v(\mathcal F(C_s,C_t,C_z))\ge40$ is valid even
when $C_s$ is not a unit. No local denominator has been discarded.

For completeness the integer $s\ne0$ and all root points are fixed before
the prime-pattern theorem is used. Its prescribed finite-exclusion version
allows every chosen prime to avoid the divisors of $s$. Then $C_s$ is a
unit at the consumed place. The old order-$40$ residual implies the new
one, and the displayed equations assign the new witnesses. This imposes
no compatibility between group orders at different primes.

Globally $p_s$ or $p_z$ may vanish at an occurrence not being consumed in
the local argument. Such points have $x=\pm d$ and are covered by the
regular extension in Lemma~\ref{lem:quartic-joint-chart}; they must not be
excluded by reusing an earlier blanket nonzero assertion. At a consumed
place $C_s=2$ gives $P=1/2$ and $V=K_L=0$, since both $p_t$ and $p_z$
are nonzero there. No division by $1-2P$ is needed. Globally the branch
$P=1/2$, $p_z=0$ forces $K_L=0$ but need not force $V=0$; the latter
value may be chosen in the completeness construction.
\end{proof}

The retained field lists can now be stated without proxy coordinates.
In a full occurrence there are the five primary $f$ values, the common
$\overline{f(J)}$, the beta residual, $E_L$, and five alternates.
In a beta-only occurrence there are two primary $f$ values, the common
field, the beta residual and two alternates. In a Lucas-only occurrence
there are four primary $f$ values, the common field, $E_L$ and four
alternates. These lists have lengths $13,6,10$. The formal fields all have
order one. The beta residual and beta alternate have order $12$; the
Lucas residual and Lucas alternates have order $40$. The denominator
alternate has order $e_0=40$ in full and Lucas-only occurrences, and
$e_0=12$ in beta-only occurrences. Its $G$ precision is the same $e_0$.

\begin{theorem}[The two height contradictions]
\label{thm:quartic-selective-arithmetic}
Retain the inherited scalar and contact-height conclusions, and the integer
root indices $n,s,t,z$ of the chosen carrier, with $s,z\ne0$. Let
$\beta=X^2+2$ and $\beta'=2\beta$. Suppose the
local tests above give $C_i\equiv n_i\pmod{p^{e_i}}$ at all consumed
contacts, with beta residual $C_A-\beta'$. Put $H_*=5H$. Suppose
\[
 n_Q^2-6<7H_*N^{19/e},\qquad
 h(\beta')<a/40,\qquad \log(40H_*)<a/40,
 \qquad a=|s|>(24H_*)^{10},
\]
where the first inequality is the contact-height bound whenever the
appropriate radical powers divide $N$. Then the changed arithmetic source
has the same integer-input conclusion as Version 3.
\end{theorem}
\begin{proof}
If $\Delta=\mathcal F(s,t,z)\ne0$, every relevant order-$40$ radical
power divides $|\Delta|$. By three-mode coverage, all three heights in
$B_L=\max(|s|,|t|,|z|)$ satisfy the assumed inequality. Since
$|\Delta|<5B_L^4$, the same elementary estimate as in Version 3 gives
\[
 B_L^{2-76/40}<24H_*.
\]
This contradicts $B_L\ge a>(24H_*)^{10}$. Therefore $\Delta=0$.
The displayed lower bound on $a$ gives $|s|\ge7$, and $z\ne0$ was
retained. The Lucas lemma therefore gives the stronger individual conclusion
$|t|>\exp(a-3)$, not merely a bound on $\max(|t|,|z|)$.

If $\beta'\ne n$, put $N=|\operatorname{num}(\beta'-n)|$.
The beta contacts cover $A,T$, so with $B=\max(|n|,|t|)$ one has
\[
 B^2<7H_*N^{19/12}+6,\qquad
 N\le2e^{h(\beta')}B.
\]
Thus $B^{2-19/12}<40H_*e^{(19/12)h(\beta')}$ and
\[
 \log B<\frac{12+19}{40(24-19)}a=\frac{31}{200}a<a-3,
\]
contrary to $B\ge|t|$. Hence $\beta'=n$ is an integer. $\beta'=2\beta=2(X^2+2)$;
writing $X=u/v$ in lowest terms gives $v^2\mid2$, hence again $v=1$.
\end{proof}


\subsection{General root carrier}
\label{sec:quartic-general-root-variant}

\begin{lemma}[Integer doubled root indices]
\label{lem:quartic-doubled-indices}
Let $R=(1+\sqrt2,2\sqrt2)$, so direct duplication gives $2R=P_*$.
Every $Q\in E(K)$ has a unique integer $n_Q$ with $2Q=n_QP_*$.
It is nonzero exactly when $Q$ is nontorsion. For such a point,
\[
 h(x(Q))>n_Q^2/4-4.
\]
\end{lemma}
\begin{proof}
The inherited generator assertion gives $2Q=n_QP_*$. The torsion
assertion makes its kernel exactly $E[2](K)$, so equivalently
$Q=n_QR+\tau$, $\tau\in E[2](K)$. Uniqueness follows from the infinite
order of $P_*$. The unchanged height bounds give
\[
 n_Q^2-6\le h(x(2Q))\le4h(x(Q))+\log676.
\]
Since $(6+\log676)/4<4$, the stated inequality follows.
\end{proof}

The root points $A,S,T,Z,J_0,J$ may now be any nontorsion points of $E(K)$.
Their indices always mean the integer doubled indices of the lemma. They
are distinct from the half-integer coefficients used for private points
in the earlier carrier convention.

\begin{lemma}[Small general-root charts]
\label{lem:quartic-general-root-charts}
For $A,S,T,Z$ retain $x,y$ and the actual first-doubling square $q=x^2$,
and impose $y^2=xq-Dx$. For $S$ also retain $xf=1$, $yz+x=0$.
Keep the nonzero tangent denominator, or the direct tangent
equation of Lemma~\ref{lem:quartic-tangent-fusion}. For each of $J_0,J$
use $f,z,k$ with
\[
 k=fz,\qquad f=z^2-Dk^2,
\]
with $z\ne0$ enforced by its retained inverse. These charts reconstruct exactly
the required nontorsion root points.
\end{lemma}
\begin{proof}
The first equations are the curve equation with a reused square. The
nonzero tangent denominator excludes all points with $y=0$, hence all
torsion. The two extra S equations preserve its reciprocal and parameter.
In the other chart, $z\ne0$ implies $f\ne0$: $f=0$ would force $k=0$
and then $z=0$. Set $x=1/f$, $p=-1/z$, $y=xp$. Dividing the second
equation by $z^2$ gives $fp^2=1-Df^2$, so $P=(x,y)$ lies on the curve
and $xy\ne0$. Conversely every nontorsion point gives these witnesses.
These are rational maps on the stated nonzero domains. The new equations
cannot justify deleting an unrelated guard without its incidence proof.
\end{proof}

For clarity, the inherited scalar argument still applies to the new
indices. At each odd contact prime of $\beta$ outside $S_*$, the
constant-left scalar norm tree excludes negative odd raw valuation and
gives all six assigned coefficient orders. The first four put $S,J_0$
in the formal group at every place above that prime. Writing
$2S=sP_*$ and $2J_0=j_0P_*$, with nonzero integer $j_0$, the last two
give $v_p(s/j_0)\ge32$, hence $v_p(s)\ge32$. The logarithm comparison
uses the exact identity $[2j_0]S=[2s]J_0$. It is valid also at inert
primes: both rational coefficients are tested and the extension is
unramified outside two. The unchanged scalar raw polynomial and its
small-prime correction detect every prime of $M_5(\beta)$, including
an odd infinite contact. Consequently
\[
 M_5(\beta)^{32}\mid s,\qquad
 h(\beta)\le H M_5(\beta)^{19}\le H|s|^{19/32}.
\]
The height theorem's domain condition holds. The value $\beta=X^2+2$
is none of $0,1,-343/32$, and $\beta=2401$ would require the nonsquare
integer $2399$ to be a rational square. Doubling both indices therefore
changes neither the scalar ratio nor its divisibility argument.

Take $H_*=5H$ and $\beta'=2\beta$. The unchanged contact-height bound
$h(x(Q))<7HN^{19/e}$ and the preceding lemma give
\[
 n_Q^2-6<28HN^{19/e}+7<35HN^{19/e}=7H_*N^{19/e}.
\]
The scalar ratio is still $s/j_0$, since doubling both indices cancels
the factor two. Its order-$32$ tests give
$h(\beta)\le H|s|^{19/32}$. The full anchor at both $2$ and $3$ is
retained. Its $3$-adic consequence, using $2S=sP_*$, is
\[
 s\in2\cdot3^{2^{24}-1}\Z\setminus\{0\},\qquad
 \log_2|s|>25165822.5.
\]
Here formal doubling at $3$ is an isometry and the retained calculation
is $v_3(f(sP_*))=2+2v_3(s)$. The $2$-adic part of the anchor is not
replaced by this consequence. Since $\log_2H_*<2093969$, one has
\[
 |s|>(24H_*)^{10},\qquad |s|>(40H_*)^{32/13},\qquad
 h(2\beta)<|s|/40,\quad\log(40H_*)<|s|/40.
\]
For example the first two required binary cutoffs are respectively
$10\cdot2093974$ and $(32/13)\cdot2093975$, both below the displayed
anchor lower bound. Thus Theorem~\ref{thm:quartic-selective-arithmetic}
applies with unchanged orders. Both beta-bearing comparators must now
test $C_A-2\beta$; the scalar definition $\beta=X^2+2$ and its raw scalar
polynomial remain unchanged.


\subsubsection{Doubled targets}
\label{sec:quartic-target-two-variant}

\begin{lemma}[Targets $2Q_i$ and $P_*$]
\label{lem:quartic-target-two}
The general-root construction may use targets $2Q_i=n_iP_*$ and fixed
denominator target $P_*=2R$, requiring only one doubling of each variable
root. It preserves the local integer congruences and integer completeness.
\end{lemma}
\begin{proof}
For soundness write $2X_0=hP_*$, $2X_i=m_iP_*$ and $2J=j_dP_*$.
The root $J$ is nontorsion, so $j_d\in\Z\setminus\{0\}$.
The doubled differences have indices $h-2$ and $m_i-2n_i$.
Lemma~\ref{lem:quartic-doubled-private-comparison} and the alternate
tests give
\[
 v_p((h-2)/j_d)\ge e_0,\qquad
 v_p((m_i-2n_i)/j_d)\ge e_i,
\]
where $e_0=40$ in full/Lucas mode and $e_0=12$ in beta-only mode.
Since $j_d$ is integral, $h\equiv2$ and $m_i\equiv2n_i$ at those orders.
In particular $h$ is a unit at the odd prime. The exact identity
$[2h]X_i=[2m_i]X_0$ permits logarithm comparison on the primary formal
group, giving ratio $m_i/h$. The retained relative errors, with
$e_i\le e_0$, give $C_i\equiv n_i\pmod{p^{e_i}}$.

For completeness let $N_{p,1},N_{p,2}$ be the two full finite group
cardinalities supplied by the inherited ordinary-prime pattern. It gives
$\gcd(N_{p,1},N_{p,2})=4$ and $p\nmid N_{p,1}N_{p,2}$. Each reduction
has exactly four rational two-torsion points, so its doubling image has
order $M_{p,i}=N_{p,i}/4$. Consequently $M_{p,1}$ and $M_{p,2}$ are
coprime and prime to $p$. Every reduction of a doubled global point is
killed by the corresponding $M_{p,i}$; no global reduction-surjectivity
or cyclicity claim is needed.
Put $j_c=\operatorname{lcm}_p M_{p,2}$ and $J=j_cP_*$, whose doubled
index is $j_d=2j_c$. Separately for each contact choose
\[
 k\equiv0\pmod{M_{p,1}},\quad
 k\equiv1\pmod{M_{p,2}},\quad
 k\equiv1\pmod{p^{40+v_p(j_c)}},\qquad k\notin\{0,1,-1\}.
\]
The three moduli are compatible. Set $X_0=kP_*$ and $X_i=kn_iP_*$.
The primary points are formal at the first place. Their differences from
the chosen targets, and $J$, are formal at the second. The alternate
logarithm ratios are $(k-1)/j_c$ and $(k-1)n_i/j_c$, while the primary
ratios are $n_i$. This supplies every required order. The finite exclusions
avoid zero points and the old open-chart exceptional cases. There is no
simultaneous CRT condition on different contact occurrences.
\end{proof}

The fixed target can be specified directly by
\[
 \begin{array}{c|cc}
 &1&\sqrt2\\\hline
 x(P_*)&57/8&-7/2\\
 q=y(P_*)/x(P_*)&-87/41&129/164\\
 F_Q=1/x(P_*)&456/1681&224/1681\\
 s_Q=(x(P_*)+DF_Q)/q&-55123/14637&19301/58548\\
 t_Q=1/q^2&7624/14161&39904/127449
 \end{array}
\]
Direct substitution verifies the curve and target identities. The source
stores $x,q,F_Q$ and the conjugates of $s_Q,t_Q$. Its remaining coefficients
are affine in $\jmath$:
\[
 r=\jmath\left(-\frac{232}{357}+\frac{86}{357}\sqrt2\right),\qquad
 u=\jmath\left(\frac{291673}{127449}-\frac{29440}{42483}\sqrt2\right),
\]
when $\jmath=1/\overline{z(J)}$. The ten constant bindings and four
affine bindings replace the fixed derivation, subject to full incidence
checks. Their isolated coefficient calculation does not bound the complete
sum of squares.



\subsection{Exact implementation refinements}


\begin{lemma}[Direct tangent equation]\label{lem:quartic-tangent-fusion}
For a curve point with $q=x^2$, replace the inverse and slope equations
$(2y)h=1$, $\lambda=(3q-D)h$ by
$2y\lambda=3q-D$. This preserves all retained rational coordinates.
\end{lemma}
\begin{proof}
The new equation and the curve force $y\ne0$. If $y=0$, then $x=0$ or
$x^2=D$, and $3q-D$ is respectively $-D$ or $2D$, a contradiction.
The unique inverse is restored by $h=(2y)^{-1}$. There is also the
polynomial restoration
\[
 h=\frac{2\lambda(3q-2D)-9xy}{4D^2},
\]
because $(3q-D)(3q-2D)-9x(x^3-Dx)=2D^2$.
For a retained second doubling its input is on the curve: with
$t=X-x$, $Y_0=-y-\lambda t$, the identity
\begin{align*}
 Y^2-X^3+DX={}&y^2-x^3+Dx+t(2y\lambda-3x^2+D)\\
 &+t^2(\lambda^2-X-2x)+(Y-Y_0)(Y+Y_0)
\end{align*}
proves curve preservation from the first doubling equations. Target-two
sources have only the first four variable-root doublings; no second set
of savings or equations is present in that branch.
\end{proof}





\begin{lemma}[Shared field-square evaluation]
\label{lem:quartic-shared-field-square}
For the common field $F=a+b\sqrt2=\overline{f(J)}$, bind once
$c=a^2+2b^2$, $d_s=2ab$. In each designated order-one packed field replace
\[
 Y=a+\theta b,\qquad PW=2aY-a^2+2b^2
\]
by $PW=c+\theta d_s$, with the same $P=r^2$. This is an exact projection
preserving $W$, every higher norm node and the pole witnesses.
\end{lemma}
\begin{proof}
Substitution gives $2a(a+\theta b)-a^2+2b^2=c+\theta d_s$.
The inverse map restores $Y=a+\theta b$, without division. All zero
cases are included. The two coefficients describe $F^2$, but their
packed use still enforces the original order-one condition on $F$;
they must not be relabelled as a new order-one test of $F^2$. The same
identity permits reuse of an already available primary-$f$ square in
its own packed numerator, after checking that square's defining rows.
\end{proof}



\begin{lemma}[First packed denominator and norm fusion]
\label{lem:quartic-first-norm-fusion}
For the first order-one primary field only, replace its virtual denominator
$W_0=T_0/r^2$ by $W_0=T_0/r$, and its first norm $N=1+rW_0^2$ by
\[
 T_0=2a_0Y_0-a_0^2+2b_0^2,\qquad rN-r-T_0^2=0.
\]
Keep every other denominator, the conic, all higher tree nodes and the
ordinary pole. The arithmetic input projection is unchanged.
\end{lemma}
\begin{proof}
The earlier tree proof still gives $v_p(W_0)>-R$, so $v_p(T_0)>0$.
More generally $v_p(T)\ge2e-1$ suffices for an order-$e$ field test:
opposite parity gives $v_p(Y)\ge e$ and
$\delta+2v_p(b)\ge2e-1$, hence
$\delta+v_p(b)\ge e+(\delta-1)/2\ge e$.
Thus the first field still has order one, and every other deduction is
unchanged. For completeness the old depth-two choice has $v_{p_0}(T_0)\ge2$
and $R=1$, so the new virtual leaf is integral and the root remains a unit.
For fixed nonzero $r$, all new tree values remain rational polynomials in
$\theta$, allowing the same exceptional-place pole approximation.

After that arithmetic change, the displayed fusion is exact: set
$T_0=rW_0$ forward and $W_0=T_0/r$ backward. The retained conic proves
$r\ne0$. These maps preserve $N$ and all later nodes, including $N=0$
or $T_0=0$ and negative rational $r$. The unguarded case $r=0$ would be
invalid and is not admitted. The old $r^2$ coordinate remains for other
leaves. The changed field is primary $f(H)$, not the separately shared
common field $\overline{f(J)}$.
\end{proof}


\subsection{Composition, completeness and finite realization}
\label{sec:quartic-composition}

For integral input first set $\beta=X^2+2$ and
take $A=\beta P_*$, whose doubled index is $2\beta$. Choose $J_0,S$ together in the old subgroup with the
required scalar and full anchor divisibilities. Their doubled-index ratio is unchanged. The Lucas existence lemma
then gives $t,z$ for the now fixed $s$, and take $T=tR$,
$Z=zR$. Fix the eight abscissa coefficients and denominator data before
using the prime-pattern theorem. Exclude the divisors of the fixed nonzero
$s$ as well as every earlier finite exceptional set.

Next choose one common comparison point $J$ using the stated least common
multiple, then choose the private points separately for each occurrence,
using precisely the target-two CRT rule.
Construct all joint-graph, reciprocal-$G$, normalized-ratio and
shared-square fields. Only then choose each private conic point and its
norm/pole witnesses by simultaneous approximation. Repeated primes retain
their common orientation but separate conic and private-point witnesses.
The constructions above give every required equation. In the reverse
direction the conic and pole decode field orders at every odd-contact
prime, without an ordinariness assumption. The graph and doubled logarithm
arguments give integer-index congruences, and the two height contradictions
give the integer input. The exact fusions preserve the specified
retained coordinates; arithmetic source replacements use the full
soundness and completeness argument rather than a claim about arbitrary
old private witness tuples.

Finally each selected source must be compiled afresh into rational
quadratic residual equations with the original input coordinate retained.
For real, hence rational, values the all-weights-one sum of their squares
vanishes exactly when every residual vanishes. Clearing nonzero rational
row denominators preserves those zero sets and gives an integral polynomial
of degree at most four.
This elementary last step does not prove the inherited arithmetic inputs.
The exact source, complete flattening map, allowed algebraic transformations,
and full coefficient collection must be checked for the single selected
endpoint. None of the local identities, allocation counts or isolated
coefficient bounds in this appendix is an assertion of that final check.
\begin{theorem}[The selected arithmetic source]
\label{thm:quartic-selected-source}
With precisely the inherited arithmetic hypotheses and source-version
qualifications stated above, the source consisting of the selected
equations has rational input projection $\Z$. Its input coordinate is the
original rational number $X$.
\end{theorem}
\begin{proof}
First suppose that rational witnesses satisfy the source. At every prime
above the cutoff with an odd raw-contact valuation, the private conic,
linear-unit pole and constant-left norm tree give the required field
orders, by Lemmas~\ref{lem:quartic-private-conic-packing} and
\ref{lem:quartic-first-norm-fusion}. This argument starts at arbitrary
odd-contact primes and does not assume that they are ordinary. The joint
graph supplies actual primary points and their regular difference
parameters. Positive primary-$f$ valuations put the primary points in the
formal group; the common conjugate-$f$ test and each alternate test allow
Lemma~\ref{lem:quartic-doubled-private-comparison} to be applied to the
doubled difference. Lemma~\ref{lem:quartic-target-two}, with denominator
order $12$ in beta-only mode and $40$ otherwise, gives the integer-index
congruences for each retained numerator. The beta quotient equations and
Lemma~\ref{lem:quartic-normalized-lucas} therefore give the precise local
tests used in Theorem~\ref{thm:quartic-selective-arithmetic}.

The retained scalar system gives its order-$32$ divisibility, while the
four-witness anchor is equivalent to the original two-place condition.
The general-root height and anchor estimates establish every cutoff of
that theorem with $H_*=5H$. The three-mode position sets cover exactly
the two maxima needed in its proof. It follows first that
$\mathcal F(s,t,z)=0$, then that $2\beta=n\in\Z$. Write $X=u/v$ in
lowest terms with $v>0$. Since $2X^2=2\beta-4$ is integral, $v^2$ divides
$2$; hence $v=1$ and $X\in\Z$.

Conversely fix any $X\in\Z$. Set $\beta=X^2+2$ and $A=\beta P_*$.
Choose the scalar pair $S,J_0$ in the retained subgroup using the
inherited scalar and full anchor construction, and choose $t,z$ by the
inherited Lucas existence lemma for the resulting nonzero integer doubled
index $s$. Set $T=tR$, $Z=zR$. All four roots and the two common roots
are nontorsion and admit the general-root charts. Fix their rational
abscissa coefficients and denominator data before applying
Lemma~\ref{lem:quartic-rational37}. Its fixed-system prime theorem permits
the finite exclusion of divisors of $s$, as required by the normalized
Lucas construction.

Choose one common $J$ by the least common multiple in
Lemma~\ref{lem:quartic-target-two}; at each occurrence choose its private
points using that lemma's three compatible congruences. No congruences
between different occurrences are imposed. The joint graph, exact
reciprocal evaluators, beta quotients and normalized Lucas equations then
have the stated witnesses and local orders. The explicit order of choices
in Section~\ref{sec:quartic-composition} now chooses each private conic
after its field coefficients are fixed. Conic approximation supplies the
integral packed leaves and a unit norm root; the fixed-parameter pole
realization supplies its remaining witnesses. This remains true with the
first denominator changed to $r$. The shared square and direct tangent
equations are exact identities with the reconstruction maps proved above.
Thus every equation of the selected source has a rational solution at
this input. This proves both implications without identifying arbitrary
old and new private witness tuples.
\end{proof}

The finite evaluations used here can themselves be written quadratically.
For precision $40$ bind four field pairs
\[
 w=Df^2,\qquad w_2=w^2,\qquad w_3=w_2w,\qquad U_6,
\]
with the fourth pair defined by the quadratic equation
\[
 U_6=\frac{a_3+a_4w+a_5w_2+
                  w_3(a_6+a_7w+a_8w_2+a_9w_3)}{a_6}.
\]
The quadratic output expression is
\[
 G_9(w)=a_0+a_1w+a_2w_2+a_6w_3U_6,
 \qquad a_6=-\frac{418}{16575}\ne0.
\]
For precision $12$ bind only $w=Df^2$ and use
$G_2(w)=1+a_1w+a_2w^2$. Substitute these expressions directly into
$\Lambda p_H=G_H$ and the beta or normalized-Lucas equations $B_i\Lambda=G_i$. Every product is quadratic in
the bound rational coordinates. Expansion gives each required $a_j$
coefficient exactly. The scale is one; no finite-error estimate is used
as a substitute for these polynomial identities.
\subsection{Exact source specification and audit boundary}
The source just proved has $4\,935$ rational coordinates,
including input zero, and $4\,513$ residual equations.
Its distinct name in \nolinkurl{quartics/quartic-4973/source-entry.json} is
\begin{center}\footnotesize\ttfamily
target2-tangent-first-norm-normalized-linear-pole-k40-q12-a25-h19-s32-4a82d9a8dc1b
\end{center}
The compressed source \nolinkurl{quartics/quartic-4973/source-residuals.jsonl.gz} has SHA-256
\begin{center}\footnotesize\ttfamily
94f8a1972cbfc0cf6493d418e628a18482307dec5c427db72e86d3eacb1a6f5d
\end{center}
Its uncompressed bytes have SHA-256
\begin{center}\footnotesize\ttfamily
4a82d9a8dc1b5c92eb40a49751c0a6aae073c8f765ca5e276a74802fd56b84be
\end{center}
The source contains $11$ modules. Its complete construction map is \nolinkurl{quartics/quartic-4973/construction-map.json}.
The independent source audit checks the full module equations, exact
ordered expression lists and contact positions, root and scalar interfaces,
all retained guards, the anchor terminal coordinate and complete flattening.
It is recorded in \nolinkurl{quartics/quartic-4973/source-audit.json}. The source-auditor SHA-256 is
\begin{center}\footnotesize\ttfamily
58bf1da4007247bc4558eacce6bdbbc83d2394ac2a47f128390f20f26567d50a
\end{center}
Those checks establish finite source identities. The source's rational
input projection is proved in Theorem~\ref{thm:quartic-selected-source},
with the inherited qualifications unchanged. A source digest is not an
arithmetic proof, and the historical base compiler remains construction
machinery rather than another final polynomial.

\subsection{Quadratic reductions preserving the rational input projection}
\label{subsec:quartic-reductions}

We describe the algebraic reductions applied to the new, fixed source.
All scalar coordinates range over $\Q$, and coordinate zero is the original
input $X$. Each field equation denotes its two rational coefficient
equations. These reductions prove equality of rational input projections
to that source; its arithmetic interpretation is
Theorem~\ref{thm:quartic-selected-source}.

\paragraph{Definitions, row operations, and interfaces.}
An equation $cv-F=0$, with $c\in\Q^*$ and $F$ independent of $v$,
permits substitution $v=F/c$ in every other equation and deletion of $v$.
Restriction is the forward solution map, and adjoining $v=F/c$ is its
inverse. Each use below retains degree at most two after substitution.
The initial quadratic lift uses the reverse operation: introduce targets
for products of two available coordinates until the source residuals are
quadratic. Evaluating these product definitions and forgetting their
targets give the two solution maps. A previously defined product can be
reused only with its retained defining equation.

An operation replacing $R_i$ by $aR_i+bR_j$, while retaining $R_j$,
has inverse $R_i=(R_i'-bR_j)/a$ when $a\ne0$ and $i\ne j$.
Thus every recorded linear row operation preserves the simultaneous zero
set. A shear replaces an old witness by $v=v'+cw$, with fixed rational
$c$ and another witness $w$ retained. Its forward and reverse coordinate
maps are $v'=v-cw$ and $v=v'+cw$. This is an invertible rational coordinate
change and preserves quadratic degree.

Every removed internal coordinate must have no occurrence outside the
replaced rows, unless the stated substitution is made in \emph{every}
such row. Retained coordinates keep their identities and values. All new
coordinates are fresh, and input coordinate zero is never a target of
deletion or redefinition. These conditions apply to the current system
after all preceding transformations.
\paragraph{Elimination modulo retained equations.}
For a recorded relation elimination, a current defining row has the form
$x-N(y)=0$ up to nonzero rational scale, with $N$ quadratic and independent
of the deleted witness $x$. For each other affected row $F_j(x,y)$ the
certificate gives an exact polynomial identity
\[
 F_j(N(y),y)=G_j(y)+\sum_k M_{jk}(y)H_k(y),
\]
up to the recorded nonzero row scales. Every $H_k$ is an unchanged,
$x$-free row, and every new row $G_j$ is at most quadratic. Forgetting
$x$ takes an old solution to a new one, since the $H_k$ vanish.
Conversely adjoin $x=N(y)$: the same identities recover every affected
$F_j=0$. Complete occurrence checks ensure that no other row uses $x$.
This is an identity over $\Q$, not a division by a possibly zero field
value. The retained support rows are never replaced in this operation.

\paragraph{The recorded contact-power circuits.}
The circuit replacements preserve the boundary values $r^2,r^{24},r^{80}$.
The replacement definitions are
\[
 r_2=r^2,\quad r_4=r_2^2,\quad r_8=r_4^2,\quad
 r_{16}=r_8^2,\quad r_{24}=r_{16}r_8,\quad
 r_{40}=r_{24}r_{16},\quad r_{80}=r_{40}^2.
\]
Their exact outputs agree with the old acyclic product definitions.
For either direction retain $r$ and the boundary outputs, and assign
each internal target its displayed power of $r$. These maps work also
for $r=0$ and use no inverse. Each deleted internal target has no use
outside the replaced definitions; every new internal target is fresh.
The certificate checks those conditions and the exact output identities
in the current system. Only circuit events with these retained powers
are included in this description.

\paragraph{A quadratic presentation of the parity block.}
This replacement concerns only its rational equations, with unchanged
input $b$. Its old private variables $t,U,e,p,q,s,h$ satisfy
\begin{align*}
 U(1+bt^2)&=b(1-bt^2),&
 p^2&=17b(1+2e-3e^2),\\
 h&=e(17b-U^2),&
 17bq^2&=(s^2-17bh)(s^2+51bh).
\end{align*}
Use instead twelve private rational variables
$u,v,a,z,e,p,c,w,s,q,j,k$ and nine quadratic equations
\begin{gather*}
 z=u^2,\quad a=bv,\quad z+va=b^2,\quad c=be,\\
 p^2-17b-34c+51ec=0,\qquad w-17bc+cz=0,\\
 j=s^2,\qquad k=bq,\qquad
 j^2+34jw-867w^2-17qk=0.
\end{gather*}
For the forward map put $\delta=1+bt^2$. It is nonzero: otherwise
$bt^2=-1$ and the first old equation gives $0=2b$, a contradiction.
Keep $e,p,s,q$ and set
\[
 u=U,\quad v=2bt/\delta,\quad a=bv,\quad z=U^2,\quad
 c=be,\quad w=bh,\quad j=s^2,\quad k=bq.
\]
The identity
$[b(1-bt^2)]^2+b(2bt)^2=b^2(1+bt^2)^2$ proves the first conic;
substitution proves the other equations.

Conversely $u^2+bv^2=b^2$. Set $h=e(17b-z)$ and keep $e,p,s,q$.
If $b=0$, take $U=t=0$. If $b\ne0$ and $u+b\ne0$, take
$U=u$ and $t=v/(u+b)$. In this case the first old residual, multiplied
by $(u+b)^2$, is
\[
 (u-b)(u+b)^2+b(u+b)v^2
       =(u+b)(u^2+bv^2-b^2)=0.
\]
If $b\ne0$ and $u=-b$, then $v=0$; take $U=b$ and $t=0$.
In every case $U^2=z$, and the remaining old equations follow from
$c=be$, $w=bh$, $j=s^2$, and $k=bq$. The possible sign change of $U$
is harmless because all old witnesses and deleted lifting coordinates
are private. These maps prove equality of existential projections;
they are not asserted to be inverse on the private witness tuples.
This rule is applied only at the certificate's height-parity blocks;
the ordinary comparison pole is the linear-unit source predicate already
proved above. The actual multiplicity is given in the following census.

\subsection{Finite realization and proof of the quartic theorem}
The fresh quadratic lift introduces only defining product equations.
Their evaluation and forgetting maps preserve rational solutions while
fixing the original input. The certificate then applies exactly the
operations below, under the current-row and complete-incidence hypotheses
just proved. The resulting file \nolinkurl{quartics/quartic-4973/residuals.json.gz}
contains $4\,552$ nonzero integer residuals of degree at most two.
Every weight in \nolinkurl{quartics/quartic-4973/weights.json} is one, and
\[
 \Phi(X_0,\ldots,X_{4972})=\sum_{i=1}^{4552}R_i^2.
\]
At rational values each square is nonnegative, so this sum vanishes if and
only if all residuals vanish. Nonzero rational row normalization preserves
each row's zero set before this integral sum is formed.

The certificate \nolinkurl{quartics/quartic-4973/certificate.jsonl.gz} records $551$ transformations:
\[
\begin{array}{lr}
\text{height-parity replacements}&1\\
\text{contact-power circuit replacements}&13\\
\text{invertible linear row operations}&266\\
\text{rational definition eliminations}&121\\
\text{invertible rational shears}&2\\
\text{eliminations modulo retained equations}&148\\
\end{array}
\]
Its SHA-256 is
\begin{center}\footnotesize\ttfamily
81ca768459a13b76dbdf4bfe81e94fbab9ab99f14a5731f6b1a1b03e28c1df2b
\end{center}
The independent verifier \nolinkurl{quartics/quartic-4973/verify_certificate.py} has SHA-256
\begin{center}\footnotesize\ttfamily
b7180b01029b3d9364d94914401a4f1f66b190762f10b99a38e9fd4f49437ffe
\end{center}
It checks the named source binding, fresh lift, every recorded
polynomial identity and reconstruction map, defining order, current-row
supports, every occurrence of each removed coordinate, final quadratic
degrees, and the bijective variable map fixing input zero. It then compares
the complete collected integer expansion, including every cross term.
The endpoint verification receipt is \nolinkurl{quartics/quartic-4973/evidence/final-verification.json}.

\begin{proof}[Proof of Theorem~\ref{thm:quartic}]
Theorem~\ref{thm:quartic-selected-source} proves both rational-input
implications for the selected arithmetic source. For any fixed rational
input, every defining lift and every recorded transformation has the
forward and reverse solution maps proved above. Their composition preserves
the existence of rational witnesses and keeps that input fixed. Hence the
final residual system has rational input projection exactly $\Z$.
The all-one sum of its squares vanishes precisely on its simultaneous
zero set, so the same equivalence holds for $\Phi$.

The literal coefficient table is the verified complete integral expansion
of that sum. Its nonzero degree-four term makes its degree exactly four;
the collected coefficients give the stated monomial count and maximum.
The checked variable map gives the stated ambient coordinates, one of
which is the input. Thus all numerical
claims and the input-projection equivalence concern this one polynomial.
\end{proof}

The finite checks do not replace the inherited rank, height, prime-pattern,
formal-group, norm or pole arguments. Their qualifications remain those
stated in the source proof and Appendix~\ref{app:sources}.

% Begin Version 5 additions
\section{Arithmetic interfaces for the smaller quartic}
\label{sec:v5-interfaces}
% Additive interfaces for the 1272-coordinate source.
% These use the original exponent-24 constant H_*, not H_19.
% They do not alter any existing source equation or polynomial artifact.

\begin{lemma}[The original scalar certificate at order thirty-two]
\label{v5:scalar32-bridge}
Let $K=\Q(\sqrt2)$, $d=3-\sqrt2$, $D=d^2$, and
$E:y^2=x^3-Dx$, with the point $P_*$ of
Section~\ref{sec:degree8-original}. Put
\[
 \beta=X^2+2,\qquad
 r=\beta(\beta-1)(\beta-2401)(32\beta+343)(2\beta-1),
 \qquad X\in\Q.
\]
For $h_0,a\in\Z\setminus\{0\}$ and points $J_0=h_0P_*$ and $S=aP_*$,
take the
two rational coefficients, in the basis $1,\sqrt2$, of
$f(J_0),f(S),z(S)/z(J_0)$ in that order. Denote them by
$X_1,\ldots,X_6$, and put
\[
 (m_1,\ldots,m_6)=(1,1,1,1,32,32),\qquad l_i=X_i/r^{m_i}.
\]
Use the same ordered norm tree as the original scalar construction:
\[
 \begin{split}
 N_r(u,v)&=u^2+rv^2,\\
 n_0&=N_r(1,l_1),\quad n_1=N_r(l_2,l_3),\quad
 n_3=N_r(l_5,l_6),\\
 W&=N_r\bigl(N_r(n_0,n_1),N_r(l_4,n_3)\bigr).
 \end{split}
\]
Let
\[
 c\in\left\{\pm\prod_{p\in I}p:
                   I\subseteq\{2,3,5,7,11,13,17\}\right\}.
\]
If $b=crW^2\ne0$ satisfies the original pole predicate $\Psi(b)$,
then
\[
 M_5(\beta)^{32}\mid a,\qquad
 h(\beta)\le H_*|a|^{3/4}.
\]
Here $M_5$ and the exact constant $H_*$ are those of
Theorem~\ref{hyp:o8-height}; in particular
$1\le H_*<2^{4783227}$.

For every integral $X$, the scalar certificate is complete with
successively chosen positive indices $h_0,a$. One may require $h_0$
to be divisible by any prescribed positive integer and $a$ to be
arbitrarily divisible and arbitrarily large. After $S$ has been chosen,
the normalized-leaf conditions and the nonvanishing of $W$ persist when
$h_0$ is replaced by every integer in one suitable congruence class
$h\equiv h_0\pmod L$. The modulus $L$ may also be required to be
divisible by any prescribed positive integer.
\end{lemma}

\begin{proof}
The value $r$ is nonzero: $\beta\ge2$, and $\beta=2401$ would
give $X^2=2399$, which has no rational solution. At a prime
$p>76176$, the multiplier $c$ is a unit. Every prime in $M_5(\beta)$
has odd valuation in $r$, by the contact computation in the proof of
Theorem~\ref{hyp:o8-height}. One can also observe that a negative
valuation of $\beta=X^2+2$ is even, so only its finite contacts
contribute here.

Write $R=v_p(r)$. The pole obstruction in Lemma~\ref{hyp:o8-pole}
and $b\ne0$ make the odd valuation $R+2v_p(W)$ positive. In a
norm node $u^2+rv^2$ the two nonzero summands have opposite valuation
parity, so there is no cancellation. The constant leaf on the
all-left path first excludes $R<0$. Then $v_p(W)>-R/2$.
Descending the tree, the implication
\[
 v_p(N_r(u,v))>-tR,\quad t<1
 \quad\Longrightarrow\quad
 v_p(u)>-tR/2,\quad v_p(v)>-(t+1)R/2
\]
shows that every leaf $l_i$ has valuation greater than $-R$.
Thus $v_p(X_i)>(m_i-1)R\ge m_i-1$, and integrality of the
valuation gives $v_p(X_i)\ge m_i$. This is the original norm
argument with the two assigned orders changed to thirty-two.

The order-one tests put $J_0$ and $S$ in the formal group at each
place above $p$. Such a completion is unramified, and the formal
logarithm is an isometry on that group. Since
$[h_0]S=[a]J_0$, it follows that
\[
 v_{\mathfrak p}\!\left(\frac{z(S)}{z(J_0)}\right)
       =v_p(a/h_0).
\]
Both rational coefficients of this ratio have valuation at least
thirty-two, so $v_p(a/h_0)\ge32$. As $h_0$ is a nonzero integer,
$v_p(a)\ge32$. This proves $M_5(\beta)^{32}\mid a$.
Theorem~\ref{hyp:o8-height} now gives
$h(\beta)\le H_*M_5(\beta)^{24}\le H_*|a|^{3/4}$.
No exponent-nineteen height estimate is used.

For completeness, let $X\in\Z$. The integer $r$ is nonzero and
$|r|>1$. The original seven-prime squareclass basis, proved in the
finite-correction argument following
Lemma~\ref{lem:original-seven-parity-completeness}, permits a choice
of $c$ such that $cr>0$ and $cr$ is square in
$\Q_2,\Q_3,\Q_{17}$. Let $\mathcal S$ be the finite prime support
of the integer $cr$.

At the places above $\mathcal S$, choose $h_0$ sufficiently
divisible that $J_0$ is in arbitrarily deep formal subgroups.
This remains possible after requiring any fixed divisibility of
$h_0$, because these subgroups have finite index in the local
elliptic groups. Having fixed this nonzero $J_0$, choose $a$
sufficiently divisible that $S=aP_*$ is still deeper. The
coefficients of $f(J_0),f(S)$ and $z(S)/z(J_0)$ then attain every
finite lower valuation bound needed here. This includes ramified
and bad primes, by working in sufficiently deep formal subgroups.
We arrange that all normalized leaves are integral when $p\mid r$
and divisible by $p$ when $p\mid c$ but $p\nmid r$.
The index $a$ may simultaneously exceed any fixed bound.

At $p\mid r$, the root of the displayed tree is congruent to one
modulo $p$, because $r\equiv0$ and the all-left leaf is one.
At $p\mid c$, $p\nmid r$, all other leaves vanish modulo $p$,
so the root is again congruent to one. Thus $W$ is a unit on
$\mathcal S$, and in particular is nonzero. At these primes
$v_p(crW^2)=v_p(cr)\ge0$; outside $\mathcal S$ its valuation is
even. Hence $b$ has square reduced denominator, is positive, and
has the same local squareclasses as $cr$ at two, three, and
seventeen. Lemma~\ref{lem:original-seven-parity-completeness}
supplies the scalar pole witnesses.

Finally, with $S$ fixed, these scalar tests are rational functions
of $J_0$ defined at the chosen point. Their coefficient maps are
continuous in the product of the completions above each prime in
$\mathcal S$. The finite lower valuation bounds persist in
sufficiently small neighborhoods. A common finite-index subgroup
of the integer orbit gives a congruence $h\equiv h_0\pmod L$
contained in all those neighborhoods. Enlarging $L$ preserves
that property. The same unit calculation proves the asserted
stability for the final shared point.
\end{proof}

\begin{lemma}[The fixed multiplier seventeen with preparation zero]
\label{v5:preparation-zero-bridge}
Let $b\in\Q^\times$, put $B=17b$, and define $\Psi_0(b)$ by
the existence of $e,p,q,s\in\Q$ satisfying
\[
 p^2=B(1-e)(1+3e),\qquad h=e(B-b^2),\qquad
 Bq^2=(s^2-Bh)(s^2+3Bh),
\]
where $h$ is the displayed expression, not a further free choice.
Then $\Psi_0(b)$ excludes negative odd valuations of $b$ at every
prime outside $\{2,3,17\}$. Conversely, it holds whenever $b$ has
square reduced denominator, is positive and square in
$\Q_2,\Q_3,\Q_{17}$, and
\[
 b<17,\qquad v_2(b)\ge3,\qquad v_3(b)\ge1,
 \qquad v_{17}(b)\ge2.
\]
No additional nonzero condition on $e,p,q,s,h$ is imposed.
\end{lemma}

\begin{proof}
Taking the preparation witnesses $t=0$ and $U=b$ in
\eqref{eq:o8-pole} gives exactly the displayed equations.
The obstruction therefore follows from
Lemma~\ref{lem:original-seven-parity-soundness}, including its
degenerate witness cases.

For completeness put
\[
 V=B(B-b^2)=B^2(1-b/17).
\]
The real inequality makes $V$ positive. The valuation bounds make
$1-b/17$ a square at two, three, and seventeen: its difference
from one has valuation at least three at two and at least one
at the two odd primes. Thus $V$ is square at these three places.

Let $\mu$ and $\theta$ be the positive squarefree representatives
of $B$ and $V$. The prescribed squareclasses of $b$ give
$\mu\equiv17\pmod{24}$; in particular $\mu>1$ and its support
contains seventeen but neither two nor three. Also
$\theta\equiv1\pmod{24}$, and seventeen is absent from its
support. At a prime $\ell\mid\mu$ outside the three exceptions,
$v_\ell(b)$ is positive and odd, since the reduced denominator
of $b$ is square. Hence $1-b/17$ is a square unit there, so $V$
and $\theta$ are local squares. The same assertion at seventeen
was already arranged. These facts make the supports disjoint.

Conversely, if $\ell\mid\theta$, then
$\ell\notin\{2,3,17\}$. If $v_\ell(b)>0$, the valuation
of $V$ is $2v_\ell(b)$; if $v_\ell(b)<0$, it is
$3v_\ell(b)$, again even because the denominator is square.
Thus an odd valuation of $V$ forces $v_\ell(b)=0$ and
$1-b/17\equiv0\pmod\ell$. It follows that
$B=17b\equiv17^2\pmod\ell$ is a nonzero square unit.
This proves the other mutual-squareness condition.

Lemma~\ref{lem:original-seven-simultaneous-realization} now gives
the two rational points in the squareclass of $B$. To make the
conversion explicit, write $V=\theta c^2$ and scale the second
point to the curve with parameter $nV$. If the first abscissa is
$Ba^2$, scale both points and parameters by $a$ so that their
abscissas are $B$ and $Bs^2$, with $s\ne0$, and their curve
parameters are $e_0$ and $e_0V$. For their ordinates $y_1,y_2$, set
\[
 e=e_0/B,\qquad p=y_1/B,\qquad q=y_2/(B^2s),\qquad
 h=e(B-b^2).
\]
The first curve equation gives $p^2=B(1-e)(1+3e)$.
Since $e_0V=B^2h$, the second gives
$Bq^2=(s^2-Bh)(s^2+3Bh)$. These are the required witnesses.
The only arithmetic existence theorem used in this specialization
is the existing simultaneous-realization lemma, with its stated
restricted-converse and descent dependencies.
\end{proof}

\begin{lemma}[Real bounds independent of later progression moduli]
\label{v5:real-progression-bridge}
Let $G$ be the closure of the integer orbit
\[
 \{([n]P_*,[n]P_*^\sigma):n\in\Z\}
       \subset E(\mathbb R)\times E^\sigma(\mathbb R),
\]
and let $G^0$ be its identity component. Fix a finite collection of
rational functions formed from nonzero integral multiples and fixed
translates of either elliptic coordinate, with all denominators
nonzero as rational functions on the relevant curve. Suppose their
pole set is contained in a finite union of inverse images of finite
sets under these nonzero multiplication-and-translation maps on
either factor. There is a bound $C$ such that every nonconstant
integer progression contains arbitrarily large positive indices at
which all these functions are defined and their real absolute values
are at most $C$. The bound is chosen from the functions before the
progression modulus is specified.

If $m=|G/G^0|$, $m\mid L$, and $m\mid h_0$, the progression
$h\equiv h_0\pmod L$ is dense in $G^0$, and its positive tail
meets every nonempty open subset of $G^0$ infinitely often.
\end{lemma}

\begin{proof}
The real elliptic groups are compact Lie groups. Their closed
subgroup $G$ is a compact abelian Lie group with finitely many
components, and $G^0$ is a torus. The projection of $G^0$ onto
each elliptic identity circle is surjective. Indeed the projected
integer orbit is infinite, since $P_*$ has infinite order and each
real embedding is injective on the group of algebraic points.
If the connected projection were trivial, the projected $G$ would
be finite; a nontrivial connected closed subgroup of a circle is
the whole circle.

Each nonzero multiplication map on a circle has finite fibers.
Consequently each of the prescribed inverse images of a finite
set has Haar measure zero in every component of $G$. The finite
union of pole sets cannot cover a component. Choose in each
component a nonempty open neighborhood whose compact closure
avoids the poles. Continuity bounds all the fixed rational
functions on these finitely many compact sets by one constant $C$.

For a progression $h=r+kn$, its closure is
$[r]P_*+[k]G$. Multiplication by the positive integer $k$ maps
$G^0$ onto itself, so this closure contains an entire component
of $G$. The positive tail of the progression is dense in its
closure: compactness gives unbounded positive multiples tending
to the identity, so each negative multiple is also a limit of
positive multiples. It therefore meets the chosen neighborhood
in that component at arbitrarily large positive indices.

If $m\mid k$, then $[k]G=G^0$. If also $m\mid r$, the coset
is $G^0$ itself. Apply this with $k=L$ and $r=h_0$ for the final
assertion.
\end{proof}

\section{The shared-point quartic in 962 variables}
\label{sec:v6-962}

This section gives a new complete quadratic compilation of the shared-point
construction in the public agentic notes~\cite{agentic-improvements}.
The shared comparison points, reciprocal-trace detectors, fixed height
scales, direct Lucas test, and affine anchor change are arithmetic
and algebraic improvements recorded in those notes. The simultaneous prime theorem used to choose
the shared points is Sun's Theorem~4.1, in the precise form of
Theorem~\ref{thm:v5-simultaneous-primes} below. Our contribution here is an
independently specified gate program, its source verification, and
further changes of coordinates and circuit identities leading to a
962-variable polynomial after 32 certified eliminations. The
definition does not depend on an unavailable coefficient table.
The relevant passages in the public notes are
\href{https://github.com/JoseBrox/diophantine-z-q_improvements/blob/f4d1b5e6c23d4adeb3ae1f986ab0bf87884bdff6/paper/diophantine_optimization_20261106_v3.tex#L2354-L3044}{shared-point construction},
\href{https://github.com/JoseBrox/diophantine-z-q_improvements/blob/f4d1b5e6c23d4adeb3ae1f986ab0bf87884bdff6/paper/diophantine_optimization_20261106_v3.tex#L3352-L3405}{fixed height scales},
and
\href{https://github.com/JoseBrox/diophantine-z-q_improvements/blob/f4d1b5e6c23d4adeb3ae1f986ab0bf87884bdff6/paper/diophantine_optimization_20261106_v3.tex#L4135-L4758}{quadratic compilation and anchor translation}.

We retain the arithmetic conclusions and conventions of
Section~\ref{sec:degree8-original}, including the restricted-converse
input to Lemma~\ref{hyp:o8-pole}. In this section the height estimate is
Theorem~\ref{hyp:o8-height}, with exponent 24 and its original constant
$H_*$; it is not the exponent-19 estimate used for the preceding quartic.
Thus
\[
  1\le H_*<2^{4783227}.
\]
The exact finite compilation and coefficient assertions below require no
arithmetic conjecture. Identifying the rational input projection with
$\Z$ uses the stated arithmetic results and the simultaneous prime theorem.

The point, pole, and correction variables below specify the
arithmetic source predicate. The final compilation subsection
gives explicit maps to its implementation coordinates.

\subsection{Points, local parameters, and quotient domains}
\label{subsec:v5-point-source}

Use the normalized curve and generator already established in
Section~\ref{sec:degree8-original}:
\[
 K=\Q(\sqrt2),\quad d=3-\sqrt2,\quad D=d^2,
 \quad E:y^2=x^3-Dx,
\]
\[
 P_*=(57/8-7\sqrt2/2,-165/8+417\sqrt2/32),
 \qquad 2E(K)=\Z P_*.
\]
The last assertion follows from Proposition~\ref{prop:generator} by the
unit change of variables recorded there. Conjugation of $K/\Q$ is
denoted by $\sigma$. All field equations below mean both rational
component equations in the basis $1,\sqrt2$.

The source uses nine nonidentity points
\[
 A,S,T,Z,H,U_A,U_S,U_T,U_Z\in2E(K).
\]
Each is represented as the double of an affine halfpoint
$(\alpha,\eta)\in E(K)$, using four rational halfpoint coordinates.
The equations for a double are
\begin{equation}\label{eq:v5-double}
 \eta^2=\alpha^3-D\alpha,\qquad
 2\eta m=3\alpha^2-D,\qquad
 x=m^2-2\alpha,\qquad y=m(\alpha-x)-\eta.
\end{equation}
Products of degree greater than two in this display are evaluated by
the quadratic gates specified below. Set
\[
 Q_0=P_*,\quad(Q_1,Q_2,Q_3,Q_4)=(A,S,T,Z),
 \quad(U_0,U_1,U_2,U_3,U_4)=(H,U_A,U_S,U_T,U_Z).
\]
Compute $[4]Q_i$ by two doublings, or by fixed rational arithmetic for
$Q_0$, and compute
\[
 V_i=U_i-[4]Q_i\quad(0\le i\le4).
\]
For the subtraction of $(x',y')$ from $(x,y)$ we retain an inverse of
$x-x'$ in $K$ and use
\[
 m=(y+y')/(x-x'),\quad x''=m^2-x-x',\quad
 y''=m(x-x'')-y.
\]

\begin{lemma}[Domains of the point quotients]\label{lem:v5-point-domains}
The equations \eqref{eq:v5-double} force $\eta\ne0$ and describe a
finite nonidentity point of $2E(K)$. Every such point has nonzero
abscissa and ordinate. The guarded differences $V_i$ have the same
properties. Consequently $f(Q)=1/x(Q)$ and $z(Q)=-x(Q)/y(Q)$ are
defined and nonzero at every point where they are used here.
\end{lemma}
\begin{proof}
If $\eta=0$, the curve equation gives $\alpha=0$ or $\alpha^2=D$.
The right side of the slope equation is then $-D$ or $2D$, respectively,
a contradiction. The ordinary duplication formulas therefore apply.
They give
\[
 x([2](\alpha,\eta))=
 \left(\frac{\alpha^2+D}{2\eta}\right)^2\ne0,
\]
because $D=d^2$ and the numerator is positive at either real embedding.
An ordinate-zero point has abscissa $0,d$, or $-d$.
The latter two are not squares in $K$: their rational norm is 7,
whereas the norm of a square is a rational square. Thus the doubled
ordinate is nonzero. A guarded affine difference is finite and is
again in $2E(K)$, so the same argument applies. Conversely any
nonidentity point of $2E(K)$ has an affine halfpoint of nonzero
ordinate and is represented by these formulas.
\end{proof}

Write
\[
 a_0=1,\quad a_1=-2/5,\quad
 a_j=\frac{4j-5}{4j+1}a_{j-1}\ (j\ge2),\qquad
 G_J(w)=\sum_{j=0}^J a_jw^j,
\]
and, for $e=16,52$, put
\begin{equation}\label{eq:v5-log}
 \ell_e(Q)=z(Q)G_{e/4-1}(Df(Q)^2).
\end{equation}
Lemma~\ref{lem:quartic-reciprocal-nonzero} applies to both truncation
degrees 3 and 12. Thus all these values are nonzero globally.

\begin{lemma}[The logarithm estimate at unramified places]
\label{lem:v5-unramified-log}
Let $p>76176$ be a rational prime and $\mathfrak p\mid p$ a place
of $K$. If a nonidentity point $Q$ is in the formal group at
$\mathfrak p$, normalize $v(p)=1$ and put $r=v(z(Q))$. Then $r\ge1$,
and for $e\in\{16,52,64\}$,
\[
 v(\log_E Q)=r,\qquad
 v(\log_E Q-\ell_e(Q))\ge(e+1)r,\qquad
 \ell_e(Q)/\log_E Q\in1+p^e\mathcal O_{K_{\mathfrak p}}.
\]
Here $\ell_{64}$ is defined by the same formula with $G_{15}$.
Ratios with numerator precision $e$ and denominator precision $e'$
in this set have the valuation of their full logarithm ratio and
relative error of valuation at least $\min(e,e')$.
\end{lemma}
\begin{proof}
The local extension is unramified, including when $p$ is inert.
The curve has good reduction and $D$ is a unit. On the formal branch,
\[
 f=z^2(1-Df^2),\qquad w=Df^2,\qquad v(w)=4r.
\]
The recurrence gives
$4w(1-w)G'(w)+(1+w)G(w)=1-w$ and hence the same identity
$\log_E Q=zG(w)$ as in Lemma~\ref{lem:quartic-reciprocal-log}.
For $j\ge1$ the denominator of $a_j$ divides $(4j+1)!$, so
\[
 v(a_j)\ge-\frac{4j+1}{p-1}.
\]
Set $m=e/4\le16$. If $j\ge m$ and $4j+1<p$, then
$v(a_jw^j)\ge4mr$. Otherwise $j\ge19044$ and
\[
 v(a_jw^j)\ge4jr-j\ge3jr\ge4mr.
\]
The lower bounds tend to infinity, proving the estimate for the
entire tail. Every nonconstant term of $G(w)$ has positive valuation,
so the leading term gives $v(\log_E Q)=r$. Division proves the
relative estimate, and division of two relative estimates proves
the ratio assertion. No step uses residue degree one.
\end{proof}

\subsection{The complete rational source}
\label{subsec:v5-complete-source}

The ordinary input is $X$, and $\beta=X^2+2$. Define
\[
 (k_1,k_2,k_3,k_4)=(16,52,52,52),\qquad
 L_i=\frac{\ell_{k_i}(U_i)}{\ell_{52}(H)},\qquad
 L'_i=\sigma\!\left(\frac{\ell_{k_i}(V_i)}{\ell_{52}(V_0)}\right).
\]
The sixteen field tests $D_j$, with their orders $m_j$, are exactly
\begin{equation}\label{eq:v5-test-list}
\begin{array}{c|c}
 f(U_0),\ldots,f(U_4),\ \sigma f(V_0),\ldots,\sigma f(V_4)&1\quad(10)\\
 L_i-L'_i\ (1\le i\le4)&16,52,52,52\\
 L_1-\beta&16\\
 F(L_2,L_3,L_4)&52,
\end{array}
\end{equation}
where
\begin{equation}\label{eq:v5-lucas-polynomial}
 F(U,V,W)=V^2-U(U-2)VW+(U-2)^2W^2-1.
\end{equation}
There is no rational proxy for an integer point label.

For $j=1,\ldots,4$ write $x(Q_j)=x_{j,0}+x_{j,1}\sqrt2$ and set
\[
 \mathcal A_0=1,\quad
 \mathcal A_{2j-1}=x_{j,0}^2,\quad
 \mathcal A_{2j}=(1+x_{j,1}^2)^2.
\]
These are positive nonzero rational squares. In fact $x(Q_j)$ is
positive at both real embeddings by Lemma~\ref{lem:v5-point-domains},
so $x_{j,0}>0$. Retain one rational coordinate $\lambda$ and put
\[
 \Theta=\{0,-1/2,-1,1/2399,-32/407\},\qquad
 F_5(T)=\prod_{\theta\in\Theta}(T-\theta),\qquad
 t_r=\mathcal A_r\lambda,\quad\delta_r=t_r-1.
\]
For every $0\le r\le8$, require $t_r\notin\Theta\cup\{1\}$ and define
\begin{align}
 R_r&=\mathcal A_r^{-5}F_5(t_r)/\delta_r^6,
 &\widetilde D_{r,j}&=D_j/\delta_r^{m_j+1},\label{eq:v5-detector-data}\\
 W_r&=\prod_{j=1}^{16}
 \left(1+R_r(\widetilde D_{r,j}/R_r^{m_j})^2\right),
 &T_r&=\Tr_{K/\Q}(W_r^{-1}),\qquad B_r=R_r/T_r^2.
 \label{eq:v5-detector-product}
\end{align}
Require $W_r,T_r\ne0$. For $B=B_r$ retain four rational witnesses
$e_r,s_r,p_r,q_r$ and impose
\begin{equation}\label{eq:v5-zero-pole}
 \begin{aligned}
 C_r&=17B_r,\qquad h_r=e_r(C_r-B_r^2),\\
 p_r^2&=C_r(1-e_r)(1+3e_r),\\
 C_rq_r^2&=(s_r^2-C_rh_r)(s_r^2+3C_rh_r).
 \end{aligned}
\end{equation}
The first two expressions are arithmetic definitions. The last two
are the preparation-zero specialization of the original pole predicate.
There is no nonzero guard on $e_r,s_r,p_r,q_r$, or $h_r$.

The scalar block shares the point $H$. Define
\[
 r_\beta=\beta(\beta-1)(\beta-2401)(32\beta+343)(2\beta-1).
\]
It is nonzero for rational $X$: the only positive possible root not
excluded by $\beta\ge2$ would require $X^2=2399$.
Let $Y_1,\ldots,Y_6$ be, in order, the two rational components of
$f(H),f(S),z(S)/z(H)$, and set
\[
 (e_1,\ldots,e_6)=(1,1,1,1,32,32),\qquad
 y_i=Y_i/r_\beta^{e_i},\qquad N_r(u,v)=u^2+rv^2.
\]
The scalar norm tree has no padding leaf. It is
\begin{equation}\label{eq:v5-scalar-tree}
 \begin{aligned}
 n_0&=N_{r_\beta}(1,y_1),\quad n_1=N_{r_\beta}(y_2,y_3),\quad
 n_3=N_{r_\beta}(y_5,y_6),\\
 W_\beta&=N_{r_\beta}\bigl(N_{r_\beta}(n_0,n_1),
                          N_{r_\beta}(y_4,n_3)\bigr).
 \end{aligned}
\end{equation}
Use eight Boolean coordinates $\epsilon_-,\epsilon_p$, for
$p\in\{2,3,5,7,11,13,17\}$, impose $\epsilon^2-\epsilon=0$, and set
\[
 c=(1-2\epsilon_-)\prod_p(1+(p-1)\epsilon_p),\qquad
 b_\beta=c r_\beta W_\beta^2,\qquad C_\beta=17b_\beta.
\]
Require $b_\beta\ne0$. Retain five rational witnesses
$t_\beta,e_\beta,s_\beta,p_\beta,q_\beta$ and impose
\begin{align}
 (1+b_\beta t_\beta^2)U_\beta
    &=b_\beta(1-b_\beta t_\beta^2),
 &h_\beta&=e_\beta(C_\beta-U_\beta^2),\label{eq:v5-scalar-preparation}\\
 p_\beta^2&=C_\beta(1-e_\beta)(1+3e_\beta),
 &C_\beta q_\beta^2&=(s_\beta^2-C_\beta h_\beta)
                     (s_\beta^2+3C_\beta h_\beta).
 \label{eq:v5-scalar-pole}
\end{align}
Here $U_\beta,h_\beta$ are arithmetic values. The preparation equation
forces $1+b_\beta t_\beta^2\ne0$: otherwise $b_\beta t_\beta^2=-1$
but its right side is $2b_\beta\ne0$.

Finally fix the integers
\begin{equation}\label{eq:v5-fixed-constant}
 C=3^{2^{26}},\qquad \mathsf S=C^2=3^{2^{27}},\qquad
 \kappa=\mathsf S^2=3^{2^{28}}.
\end{equation}
The symbol \texttt{S} in the coefficient files means $\mathsf S$;
it is not the point $S$ and is not a quantified coordinate.
Write $f(S)=f_0+f_1\sqrt2$, retain three private rational witnesses
$u,v,w$, and impose
\begin{equation}\label{eq:v5-physical-anchor}
 u^2+6v^2-2w^2-3(f_0^2+f_0f_1-f_1^2)^2+2\mathsf S^2=0.
\end{equation}
This completes the rational source. Every rational expression has the
displayed guard or a domain implication already proved. Equations
\eqref{eq:v5-zero-pole} and \eqref{eq:v5-scalar-pole} retain their
degenerate witness values.

\begin{lemma}[The three-witness anchor]\label{lem:v5-three-witness-anchor}
Equation~\eqref{eq:v5-physical-anchor} is equivalent to
$f_0/C,f_1/C\in\Z_{(2)}\cap\Z_{(3)}$. For the nonzero integer label
$a$ of $S=[a]P_*$ it implies
\[
 0\ne a\in2\,3^{2^{25}-1}\Z,\qquad
 |a|>(21H_*)^{13/2}>(20H_*)^4.
\]
Its permitted labels contain arbitrarily divisible positive integers.
\end{lemma}
\begin{proof}
Put $R=\Z_{(2)}\cap\Z_{(3)}$. The quaternion algebra $(2,3)_\Q$
ramifies exactly at 2 and 3. The quadratic embedding criterion
\cite[Proposition 14.6.7]{voight2021} shows, for rational $t\ne0$,
that $1-2/(3t^2)$ is represented by $2x^2+3y^2-6z^2$ precisely
when it is nonsquare in both $\Q_2$ and $\Q_3$. This last condition
is exactly $t\in R$. Indeed at 2 a positive valuation of $t$ gives
odd valuation of the expression, valuation zero gives residue 3
modulo 8, and negative valuation gives a square in $1+8\Z_2$.
At 3, nonnegative valuation of $t$ gives odd negative valuation,
while negative valuation gives a square in $1+3\Z_3$.
The expression is nonzero because $t^2=2/3$ has no rational solution.
Multiplying by $3t^2$ gives
\[
 t\in R\quad\Longleftrightarrow\quad
 \exists u,v,w\in\Q:\ u^2+6v^2-2w^2-3t^2+2=0.
\]
The case $t=0$ has witnesses $(0,0,1)$. A rational-square value of
$1-2/(3t^2)$ is excluded by the division-algebra criterion as well:
a nonzero trace-zero quaternion cannot square to a nonzero rational
square in a division algebra.

The form $A^2+AB-B^2$ has no nontrivial zero over $\F_2$ or $\F_3$.
Scaling a pair therefore proves
$v_p(a^2+ab-b^2)=2\min(v_p(a),v_p(b))$ for $p=2,3$.
Apply the preceding three-witness formula to
$(f_0^2+f_0f_1-f_1^2)/C^2$ and multiply its witnesses by $C^2$.
This gives exactly \eqref{eq:v5-physical-anchor}.

The local calculation in Section~\ref{subsec:o8-anchor-subgroup}
applies with exponent $2^{26}$: the reduction of $P_*$ at 3 has
order two and $v_3(z(2P_*))=1$. For even nonzero $a$ it gives
$v_3(f([a]P_*))=2+2v_3(a)$. At 2 the rational components of
$f([a]P_*)$ are integral. The anchor is consequently equivalent to
$0\ne a\in2\,3^{2^{25}-1}\Z$. Its minimum possible absolute value
has binary logarithm greater than
\[
 1+\tfrac32(2^{25}-1)=50331647.5,
\]
whereas
$\log_2((21H_*)^{13/2})<\tfrac{13}{2}(5+4783227)=31091008$.
The last comparison with $(20H_*)^4$ follows from
$21^{13}>20^8$ and $H_*\ge1$. The subgroup description proves the
last assertion.
\end{proof}

\subsection{Local tests and the height implication}
\label{subsec:v5-soundness}

\begin{lemma}[One place for every test]\label{lem:v5-one-place}
Let $R\in\Q^*$, $D_j\in K$, and positive integers $m_j$ define
\[
 W=\prod_j(1+R(D_j/R^{m_j})^2),\quad
 T=\Tr(W^{-1}),\quad B=R/T^2,
\]
with $W,T\ne0$ and a preparation-zero pole certificate for $B$.
If $p>76176$ and $v_p(R)$ is odd, then $v_p(R)>0$ and some
$\mathfrak p\mid p$ has $v_{\mathfrak p}(D_j)\ge m_j$ for every $j$.
\end{lemma}
\begin{proof}
Set $\alpha=v_p(R)$. The nonconstant term in each nontrivial factor
has odd valuation $\alpha+2(v(D_j)-m_j\alpha)$; unramifiedness
makes these valuations integral. The factor valuation is its minimum
with zero. Thus $\lambda_{\mathfrak p}=v_{\mathfrak p}(W)\le0$.
Put $\lambda=\max_{\mathfrak p\mid p}\lambda_{\mathfrak p}$.
Integrality of trace gives $v_p(T)\ge-\lambda$. Hence
\[
 v_p(B)\le\alpha+2\lambda.
\]
The left side is odd and, by Lemma~\ref{hyp:o8-pole}, not negative.
It is positive, so $\alpha>0$ and $\lambda>-\alpha/2$.
At a place attaining $\lambda$, each factor has valuation at least
$\lambda$, because all factor valuations are nonpositive. For
nonzero $D_j$ this implies
\[
 v(D_j)>(m_j-3/4)\alpha\ge m_j-3/4.
\]
Integrality gives $v(D_j)\ge m_j$; a zero $D_j$ satisfies it as well.
\end{proof}

\begin{lemma}[The two-anchor comparison]\label{lem:v5-two-anchors}
Write the integer labels of $(A,S,T,Z)$ as $(n,a,t,b)$.
If the first fourteen tests in \eqref{eq:v5-test-list} have their
specified orders at a place $\mathfrak p\mid p$, $p>76176$, then
\[
 L_i\in\mathcal O_{K_{\mathfrak p}},\qquad
 L_i\equiv n_i\pmod{p^{k_i}},\qquad
 (n_1,n_2,n_3,n_4)=(n,a,t,b).
\]
This holds at split and inert primes.
\end{lemma}
\begin{proof}
Write $H=[h]P_*$ and $U_i=[u_i]P_*$. The order-one tests put
$H,U_i$ in the formal group at $\mathfrak p$ and $V_0,V_i$ in the
formal group at the conjugate place. Positive valuation of $f=1/x$
gives formal membership on this good model. The guarded points give
$h\ne0,4$. Their full logarithm ratios, after conjugating the second,
are
\[
 r_i=u_i/h,\qquad r'_i=(u_i-4n_i)/(h-4).
\]
These identities follow from the homomorphism property on equal
formal multiples, and do not require $P_*$ itself to be formal.
At least one of $h,h-4$ is a $p$-unit. Its ratio is integral, and
Lemma~\ref{lem:v5-unramified-log} gives the same valuation to its
truncated ratio. The cross test makes the other truncated ratio
integral too. Both relative errors are therefore absolute errors,
and $L_i\equiv r_i\equiv r'_i$ modulo $p^{k_i}$. Multiplication
by $h,h-4$ and subtraction give $4L_i\equiv4n_i$. Divide by four.
At an inert prime conjugation preserves the single local valuation,
so the same argument applies.
\end{proof}

\begin{lemma}[Height transfer]\label{lem:v5-height-transfer}
Let $M_r$ be the odd-contact radical of $t_r$ for the five contacts
$\Theta$, omitting primes at most 76176. If $M_r^k$ divides a
nonzero integer $N$ for every $r$, then
\[
 h(x(Q_i))<7H_*|N|^{24/k}\qquad(1\le i\le4).
\]
At every prime counted in $M_r$, the source imposes all sixteen
field tests with their specified orders at one place above that prime.
\end{lemma}
\begin{proof}
The fractional linear map $\phi(t)=2+1/t$ carries $\Theta$ to
$\infty,0,1,2401,-343/32$. Its integral matrix has determinant $-1$.
All primes dividing the relevant fixed coefficients and contact
differences, including differences from 1, lie in
$\{2,3,5,7,11,37,109,439,2399\}$. Thus outside the cutoff
$M_r=M_5(\phi(t_r))$, and
\[
 h(t_r)\le h(\phi(t_r))+\log3
       \le(H_*+\log3)M_r^{24}.
\]
At a prime of $M_r$, $\delta_r$ is a unit. Exactly one contact
factor has odd positive valuation, while $\mathcal A_r^{-5}$ is
a rational square. Therefore $v_p(R_r)$ is odd.
Lemma~\ref{lem:v5-one-place} applied to $\widetilde D_{r,j}$ gives
the asserted orders for $D_j$.

Set $T_0=(H_*+\log3)|N|^{24/k}$. The case $r=0$ gives
$h(\lambda)\le T_0$, and the other cases give
$h(\mathcal A_r\lambda)\le T_0$. Hence
$h(\mathcal A_r)\le2T_0$. For the first component,
$h(x_{i,0})\le T_0$. For $x_{i,1}=a/b$ in lowest terms,
\[
 h(1+x_{i,1}^2)=\log(a^2+b^2)\ge2h(x_{i,1}),
\]
so $h(x_{i,1})\le T_0/2$. Finally
\[
 h(x(Q_i))\le h(x_{i,0})+h(x_{i,1})+\tfrac32\log2
 <7H_*|N|^{24/k}.
\]
\end{proof}

\begin{lemma}[The two reduced group cardinalities]
\label{lem:v5-two-reductions}
Let $p>76176$ be prime and let $\pi\in\Z[i]$ have norm $p$, with
$\pi\equiv9\pmod{16}$ and $\pi\equiv1+3i\pmod7$.
The two reductions of $E/K$ at the places above $p$ have group
cardinalities $N(\pi-1)$ and $N(\pi+1)$, in some order.
\end{lemma}
\begin{proof}
The congruences give $p\equiv1\pmod8$ and $p\equiv3\pmod7$.
Thus $p$ splits in $K$ and $(7/p)=-1$. The reductions of
$d=3-\sqrt2$ have product 7 and opposite square classes. The two
curves $y^2=x^3-d^2x$ are consequently opposite quadratic twists
of $y^2=x^3-x$ and have opposite traces.

The Frobenius calculation already given in
Section~\ref{subsec:o8-pattern-proof} applies: the order-four
automorphism is defined over $\F_p$, its rational centralizer is
$\Q(i)$, and integrality puts Frobenius in $\Z[i]$.
Write it as $a+bi$, so $a^2+b^2=p$ and the two cardinalities are
$p+1-2a$ and $p+1+2a$. Full rational two-torsion makes both
cardinalities divisible by four. Since $p\equiv1\pmod4$, this
forces $a$ odd and $b$ even. The real part of $\pi$ is odd and its
imaginary part is even. Unique factorization in $\Z[i]$ therefore
makes $a+bi$ one of $\pm\pi,\pm\overline\pi$; the associates
with factor $i$ have the wrong parity. Hence the unordered
cardinalities are exactly $N(\pi-1),N(\pi+1)$.
\end{proof}

\begin{theorem}[Arithmetic interpretation of the source]
\label{thm:v5-source-projection}
With the inherited arithmetic results and
Theorem~\ref{thm:v5-simultaneous-primes}, the rational source of
Section~\ref{subsec:v5-complete-source} has input projection exactly $\Z$.
\end{theorem}
\begin{proof}[Soundness]
Let $\mathcal B=\max(|n|,|a|,|t|,|b|)$ and
$\Delta=F(a,t,b)\in\Z$. Polynomial substitution preserves
congruences between integral local elements. Lemmas
\ref{lem:v5-two-anchors} and \ref{lem:v5-height-transfer} therefore
give $M_r^{52}\mid\Delta$ whenever $\Delta\ne0$.
The index-height estimate \eqref{eq:o8-indexheight} yields
\[
 \mathcal B^2-6<7H_*|\Delta|^{24/52},\qquad
 |\Delta|<5\mathcal B^4,
\]
the latter because $\mathcal B\ge|a|>7$. Consequently
\[
 \mathcal B^{2/13}<(7\cdot5^{6/13}+6)H_*<21H_*.
\]
For the last strict inequality, $5^6 7^{13}<15^{13}$.
This contradicts Lemma~\ref{lem:v5-three-witness-anchor}. Thus
$\Delta=0$, and Lemma~\ref{lem:o8-lucas}, with $b\ne0$, gives
$\log\mathcal B>|a|-3$.

Lemma~\ref{v5:scalar32-bridge} applies to the exact tree
\eqref{eq:v5-scalar-tree} and gives
\[
 h(\beta)\le H_*|a|^{3/4}<|a|/20.
\]
If $\beta\ne n$, put $N=|\operatorname{num}(\beta-n)|$ in lowest
terms. The order-16 test gives $M_r^{16}\mid N$. Therefore
\[
 \mathcal B^2<7H_*N^{3/2}+6,\qquad
 N\le2e^{h(\beta)}\mathcal B,
\]
and $\log\mathcal B<2\log(32H_*)+3h(\beta)$.
The anchor also gives $\log(32H_*)<|a|/20$, so the right side
is less than $|a|/4<|a|-3$, a contradiction. Hence
$\beta=n\in\Z$. Since $X\in\Q$ and $X^2=n-2\in\Z$,
the reduced denominator of $X$ is one.
\end{proof}

\begin{proof}[Completeness]
Fix $X\in\Z$, set $n=X^2+2$, and choose the Boolean correction
$c$ so that $cr_\beta>0$ is square in $\Q_2,\Q_3,\Q_{17}$,
as in the finite squareclass computation accompanying
Lemma~\ref{hyp:o8-pole}. Let $\mathcal S$ be the prime support of
$cr_\beta$. The scalar completeness construction in
Lemma~\ref{v5:scalar32-bridge} first chooses $h_0>0$ divisible by
64 with $[h_0]P_*$ sufficiently deep at these primes. It then
chooses a positive label $a$ sufficiently divisible relative to
$h_0$ and satisfying the anchor. The normalized scalar leaves are
integral at primes dividing $r_\beta$ and divisible at primes
dividing $c$ but not $r_\beta$. Choose positive $t,b$ with
$F(a,t,b)=0$ by Lemma~\ref{lem:o8-lucas}.
This fixes the four root points and all scales $\mathcal A_r$.

Temporarily put $H=[h]P_*$ and $U_i=[hn_i]P_*$.
All comparison tests are rational functions of $[h]P_*$ with
nonzero denominators at $h=h_0$. Their poles lie above finite
sets under fixed nonzero multiplication and translation maps
on the two real factors, as required by the real progression lemma.
Continuity preserves the scalar
conditions at every place over $\mathcal S$ and bounds below the
valuations of all sixteen comparison tests at the places over
$2,3,17$, throughout a progression
\[
 h\equiv h_0\pmod L,\qquad64\mid L.
\]
One obtains $L$ by putting sufficiently small open formal subgroups
inside the finitely many translated neighborhoods and taking a
common multiple of their indices. Lemma~\ref{v5:real-progression-bridge}
gives, independently of future progression moduli, real bounds for
the tests attained by arbitrarily large positive $h$ in each
nonconstant progression. The direct polynomial $F$ introduces
no denominator and has the same boundedness property.

Choose a positive integer $D_0$ such that every
$D_0^2\theta/\mathcal A_r$ is integral, and put
\[
 v=D_0^2w^2,\qquad\lambda=u/v.
\]
The contact positions are the integer forms
$u-(D_0^2\theta/\mathcal A_r)w^2$. Delete repeated forms, but retain
their incidences in the nine detectors. The five forms within one
detector are distinct, and the zero contact is the common form $u$.
Apply Theorem~\ref{thm:v5-simultaneous-primes} to these forms.
Prescribe arbitrarily high divisibility of $w$, a sufficiently
large lower bound for $u/w^2$, and primes larger than 76176 and
all prime factors of the fixed scales and contact differences.
In its odd squarefree avoidance modulus include every odd prime
factor of $3\cdot17L$ except 7.
Writing this modulus as $M$, the prescribed Gaussian classes are
\[
 \pi_p\equiv9\pmod{16},\qquad
 \pi_p\equiv1+3i\pmod7,\qquad
 \pi_p\equiv i\pmod M.
\]

Let the resulting Gaussian generators be $\pi_p$, and write
$d_-(p)=N(\pi_p-1)$ and $d_+(p)=N(\pi_p+1)$ for the cardinalities
of the two reduced elliptic groups, orienting the places by
Lemma~\ref{lem:v5-two-reductions}. These cardinalities annihilate
the reductions of $P_*$; their exact point orders are not needed.
They satisfy
\[
 v_2(d_-)=6,\quad v_2(d_+)=2,\quad
 \gcd(d_-(p),d_+(q))=4
\]
for every pair of selected primes. No odd prime factor of $L$
divides either order. For 7 this follows separately from
$d_+\equiv6$ and $d_-\equiv2\pmod7$.
Set $M_-=\operatorname{lcm}_p d_-(p)$ and
$M_+=\operatorname{lcm}_p d_+(p)$. The congruences
\[
 h\equiv0\pmod{M_-},\quad h\equiv4\pmod{M_+},\quad
 h\equiv h_0\pmod L
\]
are compatible: their pairwise gcds are $4,64,4$, respectively.
Choose a sufficiently large positive $h$ in this progression
attaining the real bounds. It avoids $0,4,-4$, so all point and
subtraction guards hold. At a selected prime, $H,U_i$ are formal
at the place of order $d_-$, while $V_0,V_i$ are formal at its
conjugate place. Their exact ratios are $n_i$. Lemma
\ref{lem:v5-unramified-log} gives all sixteen required orders,
including the order-52 polynomial test because $F(a,t,b)=0$.

We check that the nine preparation-zero poles can be filled.
For each detector,
\[
 R_r=\frac{\prod_{\theta\in\Theta}
                   (u-(\theta/\mathcal A_r)v)}
                 {v^5\delta_r^6}.
\]
The denominator is a rational square. At the five distinct contact
primes of this detector $R_r$ has valuation one, and elsewhere
its valuation is even. The primes do not divide $v$: otherwise
two distinct forms in the same detector would give the same
prime value, impossible since $w>0$ and their coefficients are
distinct. The lower bound for the primes excludes the fixed scale
and difference denominators. At either place over a
contact prime a factor of $W_r$ has valuation zero or negative,
and its inverse has residue one or zero, respectively. At the
successful place every inverse factor has residue one.
Consequently $T_r$ has residue one or two and is a unit.
It follows that $B_r$ has positive odd valuation at its contact
primes and even valuation elsewhere. Its reduced denominator is
a square.

The prescribed Gaussian classes make every contact prime square
in $\Q_2,\Q_3,\Q_{17}$. The last product formula therefore makes
$B_r$ square in each of those fields. The uniform real bounds
and $\lambda\longrightarrow+\infty$ give
\[
 R_r\sim\mathcal A_r^{-6}\lambda^{-1},\quad
 W_r=1+O(\lambda^{-3}),\quad T_r=2+O(\lambda^{-3}),\quad
 B_r\sim(4\mathcal A_r^6\lambda)^{-1}.
\]
In particular $0<B_r<17$. At a small prime $p=2,3,17$, put
$\nu_p=v_p(v)$ and $\eta_{p,r}=v_p(\mathcal A_r)$.
The common form $u$ is a selected prime greater than these small
primes, so $v_p(u)=0$ and $v_p(\lambda)=-\nu_p$.
The fixed neighborhoods give $v(D_j)\ge-c_p$. For sufficiently
large prescribed $\nu_p$ one has
\[
 v_p(R_r)=\nu_p-6\eta_{p,r},\qquad
 v(\widetilde D_{r,j})\ge-c_p+(m_j+1)(\nu_p-\eta_{p,r}).
\]
The nonconstant factor term has valuation at least
$3\nu_p+(10m_j-8)\eta_{p,r}-2c_p$, tending to infinity.
Thus $T_r$ is as close to 2 as needed, including at the ramified
prime 2, and
\[
 v_p(B_r)=\nu_p-6\eta_{p,r}-2v_p(2).
\]
Prescribing the divisibility of $w$ before applying the prime theorem
therefore ensures $v_2(B_r)\ge3$, $v_3(B_r)\ge1$,
$v_{17}(B_r)\ge2$. Lemma~\ref{v5:preparation-zero-bridge} supplies
the four witnesses in every coordinate pole.

Finally the progression preserved all scalar leaf conditions.
The constant leaf on the all-left path makes $W_\beta$ a unit at
every prime dividing $r_\beta$; the divisibility of the other
leaves makes it a unit at primes dividing $c$ but not $r_\beta$.
Since $|r_\beta|>1$, it is nonzero. Therefore
$b_\beta=cr_\beta W_\beta^2$ is positive, has square reduced
denominator, and is locally square at $2,3,17$.
Lemma~\ref{hyp:o8-pole} supplies the scalar preparation and pole
witnesses for this final $H$. The anchor already holds for $S$.
All witnesses have now been supplied without changing the root points.
\end{proof}

% Replace the former exact-compilation subsection with this fragment.
% Retain the preceding shared-point arithmetic source and its hypotheses.
% No claim here removes or proves its simultaneous-prime input.
\subsection{The quadratic compilation in 962 variables}
\label{subsec:v6-962-compilation}

We compile the shared-point source above using different point charts,
pole coordinates, and arithmetic circuits. Every change has explicit
rational witness maps. The arithmetic source, including its simultaneous
prime input, is unchanged. Throughout this subsection
\[
 K=\Q(\sqrt2),\qquad D=(3-\sqrt2)^2=11-6\sqrt2,
 \qquad \mathsf S=3^{2^{27}}.
\]
An equation in $K$ means its two rational component equations. The fixed
integer $\mathsf S$ is a coefficient, not a coordinate. An assignment of
a quadratic expression to one rational coordinate adds one quadratic
equation. Affine expressions require no coordinate.

\begin{theorem}[The 962-variable quartic]
\label{thm:v6-962-polynomial}
The complete table specified below defines
\[
 P_{962}\in\Z[X_0,\ldots,X_{961}]
\]
with degree four, exactly $55409$ nonzero coordinate monomials, integer
coefficient content one, and coefficient height
\[
 H(P_{962})=64\mathsf S^2=64\cdot3^{2^{28}}<2^{429496736}.
\]
It is the sum of squares of $919$ quadratic residuals and has the same
rational input projection as the shared-point source of
Theorem~\ref{thm:v5-source-projection}. Consequently, under that theorem's
arithmetic hypotheses, for every $x\in\Q$,
\[
 x\in\Z\quad\Longleftrightarrow\quad
 \exists X_1,\ldots,X_{961}\in\Q,\quad
 P_{962}(x,X_1,\ldots,X_{961})=0.
\]
The $962$ variables include the input and $961$ existential witnesses.
\end{theorem}

\paragraph{A quadratic chart for the nine points.}
For each point introduce $t,h\in K$ and set
\begin{equation}\label{eq:v6-point-chart}
 u=t^2,\qquad X=h^2,\qquad
 u(X-u/4)=D,\qquad Y=h(u/2-X).
\end{equation}
The three assignments use six rational coordinates. The displayed
constraint contributes two quadratic residuals. The four rational
components of $t,h$ replace the four halfpoint components.

\begin{lemma}[Exact point-chart maps]\label{lem:v6-point-chart}
The chart \eqref{eq:v6-point-chart} describes precisely the finite doubles
used in the shared-point source, with the same rational witnesses up to
the following invertible changes. It preserves all required nonzero
point-coordinate domains. The rational component $X_0$ of $X$ is positive.
\end{lemma}
\begin{proof}
For an old halfpoint $(a,b)$ with $b^2=a^3-Da$ and its tangent equation,
both $a,b$ are nonzero. Indeed $b=0$ would give $a=0$ or $a^2=D$,
whereas the tangent numerator $3a^2-D$ would be $-D$ or $2D$. Then
$a=0$ is excluded by the curve equation. Define
\[
 t=b/a,\qquad h=(a^2+D)/(2b)=a/t-t/2.
\]
The curve equation gives $t^2=a-D/a$. Substitution gives
$t^2(h^2-t^2/4)=D$. The doubling formulas give exactly $X=h^2$
and $Y=h(t^2/2-h^2)$.

Conversely, a chart solution has $t\ne0$ because $D\ne0$. Set
\[
 a=th+t^2/2,\qquad b=at,\qquad m=h+t.
\]
If $a=0$, then $h=-t/2$, contradicting the chart equation. Thus $b\ne0$.
Direct expansion gives
\[
 a^2-t^2a=D,\quad b^2=a^3-Da,\quad
 2bm=3a^2-D,\quad m^2-2a=h^2,
\]
and $m(a-h^2)-b=h(t^2/2-h^2)$. The substitutions are inverse, so they
recover exactly the old finite double. Lemma~\ref{lem:v5-point-domains}
therefore supplies its nonzero abscissa and ordinate, and the same remains
true for the unchanged guarded differences and subsequent doublings.

Both real embeddings of $D$ are positive. Since $t\ne0$, the chart
equation implies $h^2=t^2/4+D/t^2>0$ under both embeddings. Hence
$X_0$, the average of these two positive values, is positive.
\end{proof}

The five affine subtraction formulas keep their original inverses of
the differences of abscissas. Neither a subtraction guard nor a point
occurrence is deleted. Relative to the direct halfpoint doubling circuit,
the chart saves two rational arithmetic coordinates for each of the nine
points.

\paragraph{The logarithm polynomial.}
Retain the source's coefficients
\[
 a_0=1,\quad a_1=-2/5,\quad
 a_j=a_{j-1}(4j-5)/(4j+1).
\]
For $P(w)=\sum_{j=0}^{12}a_jw^j$ write
\begin{equation}\label{eq:v6-log-square-completion}
 P(w)=a_{12}Q(w)^2+R(w),\qquad
 \deg Q=6,\quad Q\text{ monic},\quad\deg R\le5.
\end{equation}
This identity is constructed without any witness division. Starting
with the leading coefficient of $Q$ equal to one, match degrees eleven
down to six successively. At each step the new coefficient is determined
by division by $2a_{12}\ne0$. The remainder then has degree at most five.
After computing $w^2,w^3$, every term through $w^6$ is quadratic in
$1,w,w^2,w^3$. Assigning $Q$ and then $a_{12}Q^2+R$ therefore uses
four field assignments, rather than five. The eight order-$52$
logarithms save sixteen rational coordinates. The order-$16$ logarithms
keep their exact degree-three polynomials. No truncation order changes.

Usually the final logarithm is $\ell=zP(Df^2)$, with
$f=1/x$ and $z=-x/y$. For $U_0$ and $V_0$ only, omit the separate $z$
coordinates and impose the direct quotient equations
\begin{equation}\label{eq:v6-log-quotient}
 y\ell=-xP(Df^2).
\end{equation}
The polynomial output is an affine expression in assigned coordinates,
so these are quadratic component equations. The ordinate is nonzero by
Lemma~\ref{lem:v5-point-domains}; hence \eqref{eq:v6-log-quotient}
is exactly the old logarithm. The source's logarithm-nonvanishing result
and every logarithm-ratio domain remain unchanged. In the scalar block
replace $z(H)\rho=z(S)$ by
\begin{equation}\label{eq:v6-scalar-ratio}
 x(H)\rho=-y(H)z(S).
\end{equation}
Both $x(H),y(H)$ are nonzero, so this has the same unique value
$\rho=z(S)/z(H)$.

\paragraph{Products, tests, and detector evaluation.}
For each of the three orders, compute the elementary symmetric
coefficients of $\prod_j(1+qD_j^2)$ by a balanced product tree.
At an internal node with coefficient vectors $a_i,b_j$, assign each
component of $\sum_{i+j=k}a_ib_j$ as one quadratic expression.
This is the ordinary coefficient-convolution identity. It preserves
every coefficient and every test, including its assigned order.

The degree-ten detector polynomial is evaluated in two blocks:
\begin{equation}\label{eq:v6-detector-block-evaluation}
 \sum_{j=0}^{10}c_jq^j
 =\sum_{j=0}^{4}c_jq^j
   +q^5\left(\sum_{j=0}^{5}c_{j+5}q^j\right).
\end{equation}
Compute $q^2,q^3,q^4,q^5$ using four scalar assignments, assign the
two components of the parenthesized block, and assign the two final
components directly. The low block is not assigned separately.
Every final summand is quadratic in the available coordinates. This
uses eight rational coordinates. The degree-two and degree-four
polynomials use their direct power sums, with their two component
sums assigned once. These formulas are identities over $K$.

The Lucas expression is regrouped as
\[
 t^2-a(a-2)tb+((a-2)b)^2-1
 =t^2-(at)C+C^2-1,\qquad C=(a-2)b.
\]
Thus one field multiplication is removed without changing the tested
polynomial. The fixed quintic is evaluated by its exact coefficients
using $t^2$; it remains
$F_5(t)=t(t+1/2)(t+1)(t-1/2399)(t+32/407)$.

\paragraph{Weight quotients and their domains.}
Fix a detector and write $A$ for its scale, $\delta=A\lambda-1$,
and $e$ for its weight. The scales are
\[
 1,\quad x(Q)_0^2,\quad (1+x(Q)_1^2)^2
 \qquad(Q=A,S,T,Z).
\]
They are all positive by Lemma~\ref{lem:v6-point-chart}. Retain the
inverse $\delta i=1$, and replace the separate fifth-power product
and inverse of $F_5$ by
\begin{equation}\label{eq:v6-weight-domain}
 F_5(A\lambda)e=(A\delta)^4(A\delta).
\end{equation}
Here $(A\delta)^4$ denotes the output of its retained power-assignment
gates; the displayed product is quadratic in the available coordinates.
The right side is nonzero, so this equation forces both $F_5(A\lambda)$
and $e$ nonzero. In particular the removed inverse of $F_5$ is recovered
uniquely. Define the low-order weight by
\[
 \delta q_1=e i^2,
\]
where $i^2$ is likewise a retained square-assignment output. This gives
the unchanged $q_1=e/\delta^3$. Retain the repeated-squaring
values $e^6,e^{24},e^{96}$ and set
\[
 c=q_1e^6,\qquad q_{16}=ce^{24},\qquad q_{52}=ce^{96}.
\]
These are the old $q_1e^{30}$ and $q_1e^{102}$, with one shared product.

If the detector product is $W=w_0+w_1\sqrt2$, retain $w_0j=1$ and use
\begin{equation}\label{eq:v6-norm-final-weight}
 w_0n=(w_0^2-2w_1^2)/2,\qquad ev=in,\qquad B=nv.
\end{equation}
The trace guard is essential and is retained. It gives $W\ne0$ in $K$,
so $\operatorname{Norm}(W)\ne0$ and
$n=\operatorname{Norm}(W)/(2w_0)\ne0$. Equations
\eqref{eq:v6-weight-domain}--\eqref{eq:v6-norm-final-weight} therefore
give
\[
 B=\frac{i}{e}\left(\frac{\operatorname{Norm}(W)}{2w_0}\right)^2\ne0,
\]
the original detector input. The reverse map restores $R=i/e$, its
old inverse intermediates, and $n^2$ by their unique rational values.
No inverse is inferred from an unguarded zero-over-zero quotient.

\paragraph{The scalar norm and finite corrections.}
The scalar contact polynomial $r$ and its inverse $ri=1$ remain.
Writing $f=f(H)_0$, $u=f(H)_1$, and $g=f(S)_0$, retain the leaf
$\ell=ui$ and replace two norm-tree
nodes by
\[
 rn_0=r+f^2,\qquad rn_1=u\ell+g^2.
\]
with the existing leaf $ui$ retained. Since $r\ne0$, these equal
$n_0=1+r(fi)^2$ and $n_1=(ui)^2+r(gi)^2$. All other nodes and all
orders in the scalar norm tree are unchanged.

Replace the eight Boolean correction choices and their seven product
coordinates by $c_0,\ldots,c_7\in\Q$ satisfying
\begin{equation}\label{eq:v6-correction-chain}
 c_0^2=1,\qquad
 (c_j-c_{j-1})(c_j-p_jc_{j-1})=0\quad(1\le j\le7),
\end{equation}
where $(p_1,\ldots,p_7)=(2,3,5,7,11,13,17)$. Induction shows
$c_j\ne0$ and $c_j/c_{j-1}\in\{1,p_j\}$. Thus $c_7$ is exactly an
allowed signed squarefree correction. Conversely every original
choice determines such a chain. Recover the original Booleans by
\[
 \epsilon_-=(1-c_0)/2,\qquad
 \epsilon_j=\frac{c_j/c_{j-1}-1}{p_j-1}.
\]
All denominators are nonzero. The scalar preparation coordinate,
its defining quotient, and the explicit inverse of its pole input
are retained.

\paragraph{The ten pole blocks.}
Write $b$ for a pole input, $C=17b$, and $U$ for its prepared value.
For the nine coordinate poles $U=b$. In the scalar pole $U$ retains
the source's preparation formula and its denominator domain.
Every $C$ is nonzero: this was proved above for coordinate detectors,
and the scalar input has its retained inverse.

Replace the old witness $e_0$ by $z=Ce_0$. Introduce the five
intermediates
\[
 P=p^2,\quad F=U^2,\quad H=z(C-F),\quad V=s^2,\quad Q=q^2,
\]
and impose
\begin{equation}\label{eq:v6-scaled-poles}
 CP=(C-z)(C+3z),\qquad CQ=(V-H)(V+3H).
\end{equation}
Every equation is quadratic. Under $z=Ce_0$, the first residual is
$C$ times the old residual $p^2-C(1-e_0)(1+3e_0)$, and the second
is unchanged because $H=Ce_0(C-U^2)$. The inverse map $e_0=z/C$
is defined. This saves one arithmetic coordinate at each pole,
without changing the pole predicate or any preparation condition.

\paragraph{Counts and the final polynomial.}
The initial allocation is
\[
\begin{array}{l|r|r}
 &\text{coordinates}&\text{structural residuals}\\\hline
 \text{input}&1&0\\
 \text{nine point charts}&36&18\\
 \text{common slope}&1&0\\
 \text{nine coordinate poles}&36&18\\
 \text{scalar pole with preparation}&5&2\\
 \text{private anchor}&3&1\\
 \text{signed correction chain}&8&8\\\hline
 \text{total}&90&47
\end{array}
\]
The complete program has $904$ arithmetic coordinates and their
$904$ defining equations. Hence the raw counts are $994$ coordinates
and $951$ quadratic residuals.

The private anchor uses exactly the old affine change
\[
 u=\mathsf S+(U-3V+3W)/4,\quad
 v=\mathsf S/2+V/4,\quad w=3\mathsf S/2+W/4.
\]
Its determinant is $1/64$. If $T=f_0^2+f_0f_1-f_1^2$, sixteen times
the old anchor after substitution is
\[
 -48T^2+U^2-6UV+6UW+15V^2-18VW+7W^2+8\mathsf S U.
\]
No anchor condition or witness is removed.

Each certified compiler elimination solves a row $aY+q=0$ with
$a\in\Q^*$ and $q$ independent of $Y$, substitutes $Y=-q/a$
everywhere, and removes that row. It is accepted only if every
remaining residual is still quadratic. The reverse map reinserts
the unique displayed value of $Y$. There are exactly $32$ such
eliminations, no row additions, and no omission of an unused coordinate.
The input is retained. Consequently
\[
 90+904-32=962,\qquad47+904-32=919.
\]
Normalize the remaining rows to primitive integer polynomials
$q_1,\ldots,q_{919}$ and define
\begin{equation}\label{eq:v6-final-polynomial}
 P_{962}(X_0,\ldots,X_{961})=\sum_{\nu=1}^{919}q_\nu^2.
\end{equation}
The complete coefficient table is
\nolinkurl{quartics/quartic-962/polynomial.json}, with SHA-256
\begin{center}\small
\nolinkurl{6b798de7d00ac3c7d725ca0a69d39f1eac1ddb02fee2756cb6323770b952af32}.
\end{center}
An entry with coefficient $c$, exponent \texttt{sPower} $k$, and
coordinate list $[i_1,\ldots,i_d]$ means
$c\mathsf S^kX_{i_1}\cdots X_{i_d}$. Every entry is included and repeated
indices give powers. This defines an ordinary integer polynomial.

\begin{proof}[Proof of Theorem~\ref{thm:v6-962-polynomial}]
For a source solution apply the point-chart, correction-chain and
pole substitutions, and evaluate all new assignment and quotient
coordinates in dependency order. The identities and domain proofs
above show that every structural residual vanishes. Conversely apply
the inverse substitutions, restore every deleted quotient by its
proved nonzero denominator, and evaluate the source assignments.
Every source point, logarithm, test, detector input, scalar value and
anchor value is recovered. These maps fix $X_0$. The affine anchor
change, certified eliminations and nonzero row normalizations preserve
the same projection. Over $\Q$, \eqref{eq:v6-final-polynomial} vanishes
exactly when all $919$ residuals vanish. This proves the projection
statement, and Theorem~\ref{thm:v5-source-projection} gives its conditional
arithmetic interpretation.

Independent exact convolution gives $55409$ distinct monomials.
Of their coefficients, $55401$ are $\mathsf S$-independent integers,
seven are $c\mathsf S$ with $0<|c|\le768$, and one is
$64\mathsf S^2$. The largest absolute $\mathsf S$-independent coefficient
has bit length $1526$. Thus no coefficient vanishes on specialization
and $64\mathsf S^2$ is the exact height. It occurs at $X_{79}^2$,
where $X_{79}=U$. The coefficient of $X_0^4$ is one, proving degree
four and integer content one. Finally
$\log_2(64\mathsf S^2)<6+(8/5)2^{28}<429496736$ because $3^5<2^8$.
\end{proof}

The frozen package includes the builder, raw systems, all $904$ gate
definitions, the exact elimination and normalization certificate,
compact coordinate map, final residuals, complete coefficient table,
source-equivalence proofs and independent reviews. Its manifest has
SHA-256
\begin{center}\small
\nolinkurl{36fca72cb0eb126193ffa43d890ff09d5e29cbf1274561d2632ed1e227163787}.
\end{center}
The independent checker verifies every gate, all $894$ certificate
operations, the complete sum-of-squares expansion, support,
integrality, content and height. A fresh isolated reproduction gives
the same seven deterministic artifacts byte for byte. These finite
checks establish the stated algebraic facts; they do not prove any
inherited arithmetic input.


% End Version 5 main additions
\section{Why a single defining polynomial must have degree at least four}
\label{sec:degree-lower-bound}

The argument in this section comes from the agentic notes
\cite[section on unconditional lower bounds]{agentic-improvements}.
We give its specialization to $\Q$, including the treatment of singular
cubics, cones and geometric reducibility.  It is independent of the
arithmetic inputs to our constructions.

For $F\in\Q[T,Y_1,\ldots,Y_r]$, write
\[
 S_F=\{t\in\Q:\exists\,y\in\Q^r,\ F(t,y)=0\}.
\]
All degrees below are total degrees, including the input variable $T$.

\begin{theorem}[Degree obstruction]
\label{thm:degree-four-necessary}
If $S_F=\Z$, then $\deg F\geq4$.
\end{theorem}

We use two published results.  The integral form of Hilbert's
irreducibility theorem says that an irreducible polynomial
$g(T,Y)\in\Q(T)[Y]$ admits infinitely many specializations $T=t\in\Z$
which preserve its degree in $Y$ and are irreducible over $\Q$.
Deleting finitely many values also avoids its coefficient poles and
any prescribed nonzero polynomial in $t$.
This integral formulation is recalled in the introduction of
M\"uller~\cite{Muller2002HIT}; merely obtaining specializations in
$\Q$ would not suffice here.
The second result is Koll\'ar's theorem
\cite[Theorem~2]{Kollar2002Cubic}: a geometrically integral cubic
hypersurface of dimension at least two over a perfect field is
unirational over that field if it has a rational point and is not a cone.
This is a special case of his statement for irreducible cubics.
We apply it only over $\Q$, and remove every vertex direction before
applying it.  In particular, we do not assume that a cone over a plane
cubic is unirational.

\begin{lemma}\label{lem:nonconstant-rational-not-integral}
A nonconstant rational function $R\in\Q(U_1,\ldots,U_m)$ cannot take
only integer values at the rational points of a nonempty Zariski-open
subset of $\mathbb A^m_{\Q}$ on which it is defined.
\end{lemma}
\begin{proof}
Shrink the open subset to a nonempty principal open subset and choose
a connected real open ball in it, avoiding the denominator of $R$.
Rational points are dense in this ball.  If their images belong to the
closed set $\Z$, continuity puts the image of the entire ball in $\Z$.
Connectedness then makes $R$ constant on the ball.  Its numerator minus
that constant times its denominator vanishes on a real open set, hence
is the zero polynomial.  This contradicts nonconstancy.
\end{proof}

Consequently, if an integral variety $X/\Q$ is unirational and
$f\in\Q(X)$ is nonconstant, the rational values of $f$ cannot all be
integers.  Indeed, compose $f$ with a dominant rational map from
projective space to $X$, and restrict to a nonempty affine open subset
of its domain.  Dominance keeps this composite nonconstant, so
Lemma~\ref{lem:nonconstant-rational-not-integral} applies.

\begin{lemma}\label{lem:quadratic-projection-finite}
If $\deg F\leq2$ and $S_F\subseteq\Z$, then $|S_F|\leq2$.
\end{lemma}
\begin{proof}
Otherwise choose rational zeros $a,b,c$ with distinct input
coordinates.  If they lie on an affine line, the restriction of $F$
to that line has three zeros and degree at most two, so vanishes
identically.  The input is nonconstant on the line and takes every
rational value, a contradiction.

If the points are not collinear, let $L$ be their rational affine
plane.  If $F|_L=0$, a line in $L$ with nonconstant input gives the same
contradiction.  If $F|_L$ factors over $\Q$ into affine linear
polynomials, its zero set consists of at most two rational lines.
One contains two of the three selected points and again has
nonconstant input.  A nonzero polynomial of degree at most one
cannot vanish at three noncollinear points.

It remains to consider an irreducible quadratic on $L$.  Its
projective conic has a smooth rational point: a reduced irreducible
quadratic in characteristic zero has at most one geometric singular
point, while the three chosen affine points are distinct.  A conic
with a smooth rational point is rational.  Explicitly, intersecting
it with the pencil of lines through that point gives a rational
parametrization by taking the second intersection point.  The input
function is nonconstant because it takes three different values.
Composing it with this parametrization contradicts
Lemma~\ref{lem:nonconstant-rational-not-integral}.
\end{proof}

\begin{proof}[Proof of Theorem~\ref{thm:degree-four-necessary}]
Suppose $S_F=\Z$ and $\deg F\leq3$.  The polynomial is neither zero
nor a nonzero constant.  Lemma~\ref{lem:quadratic-projection-finite}
excludes degree at most two.  It also excludes a reducible cubic:
every nonconstant irreducible factor then has degree at most two,
and the finite union of their projections cannot equal $\Z$.
Thus $F$ is irreducible of degree three.

Homogenize $F$ to $H$ on a rational vector space $E$, using $X_0$ as
the homogenizing coordinate and $X_1$ for the input coordinate.
Thus on $X_0\ne0$ the input is $X_1/X_0$.
The homogeneous polynomial $H$ is irreducible over $\Q$.
If it is geometrically reducible, its geometric irreducible factors
form a single Galois orbit with equal degrees.  In characteristic
zero they occur without repetition.  Since the total degree is
three, there are exactly three conjugate linear factors.
Every rational zero lies on each of their hyperplanes.
Take rational zeros normalized by $X_0=1$ with inputs zero and one.
Their whole affine joining line lies on all three hyperplanes and
its input takes every value in $\Q$, a contradiction.
We may therefore assume $H$ is geometrically integral.

Define the rational vertex subspace
\[
 W=\{v\in E:D_vH=0\text{ identically}\},\qquad
 D_vH(x)=\left.\frac{d}{ds}H(x+sv)\right|_{s=0}.
\]
These are linear conditions on $v$.  In characteristic zero,
$D_vH=0$ is equivalent to $H(x+sv)=H(x)$ for all $x,s$.
Thus $H$ descends to a cubic $\overline H$ on $E/W$.
After a basis of $W$ is extended to one of $E$, this simply deletes
the coordinates on which $H$ does not depend.
The quotient has no remaining vertex direction.  This assertion
holds over the algebraic closure as well, since the kernel of a
matrix over $\Q$ commutes with extension of scalars.  Hence its
projective hypersurface is not a cone.

For $v\in W$ and a rational zero $x$ with $x_0=1$, translation by
$sv$ gives rational zeros whose inputs, whenever defined, are
\[
 \frac{x_1+s v_1}{1+s v_0}.
\]
Unless $v_1=x_1v_0$, this is a nonconstant fractional linear function
of $s$, with cofinite image in $\Q$.  Its image cannot be contained
in $\Z$.  Applying the resulting equality first at input zero and
then at input one gives $v_0=v_1=0$ for every $v\in W$.
Therefore $X_0$ and $X_1$ descend to independent linear forms on
$E/W$.

Choose quotient coordinates beginning with these forms.
The affine chart $X_0=1$ of $\overline H=0$ has exactly the same
input projection as $F=0$.  To see both inclusions, project rational
points from $E$ to $E/W$, and in the reverse direction use a rational
linear section of that projection.  This section preserves $X_0$
and $X_1$, and $H$ is constant along each fiber.

If the quotient projective hypersurface has dimension at least two,
Koll\'ar's theorem applies.  It is geometrically integral, has a
rational point in $X_0=1$, is not a cone, and the ground field is
perfect.  It is therefore unirational.  Its input function is
nonconstant, since its rational values include all integers.
The consequence of Lemma~\ref{lem:nonconstant-rational-not-integral}
then contradicts that its input projection is contained in $\Z$.

If the quotient has dimension one, its affine chart is defined by
an irreducible polynomial $G(T,Y)\in\Q[T,Y]$.  If its degree in $Y$
is zero, its input projection is finite.  If that degree is one,
solving the linear equation for $Y$ gives a rational point above
every rational $T$ outside a finite set, also impossible.
If its degree in $Y$ is at least two, Gauss's lemma makes it
irreducible in $\Q(T)[Y]$.  Integral Hilbert irreducibility supplies
an integer $t$ for which $G(t,Y)$ remains irreducible of that same
degree.  It has no rational root, contradicting $S_G=\Z$.
Finally, a zero-dimensional projective quotient has only finitely
many rational points and cannot have input projection $\Z$.
These exhaust the possibilities and prove the theorem.
\end{proof}

\begin{remark}
This lower bound concerns a single polynomial equation and counts
degree in every variable.  It allows systems of quadratic equations;
over $\Q$, their sum of squares is quartic.  It puts no lower bound
here on the number of variables in a degree-four definition.
\end{remark}


\appendix
\section{Dependency order in the arithmetic argument}\label{app:obligations}
The even coefficient bound uses stationary-base selection as well as
integrality. The odd minimum uses a fixed nonzero witness and the prior
genus-index lower bound. Support removal is applied only to the actual
unramified specialization. The functional bound and determinant clearing
are proved separately, and the same corrected coordinate is used at every
address. These constructions are given in Sections~\ref{sec:rc}--\ref{sec:limits}
and Appendices~\ref{app:even-source}--\ref{app:endpoints}.

The conversion to ordinary integers uses the canonical model and total
multiplication, followed by the exact computation relation. Only the
shared-shift separation argument in Section~\ref{sec:bridge} identifies
the external rational variable with the ordinary value of its integer
index. The literal substitution and size calculation then apply to the
same expression $P_*$.

\section{Source identities and scope}\label{app:sources}
The companion \texttt{SOURCE\_MAP.json} records exact paths, hashes and
locators, distinguishing mathematical and construction sources, artifacts,
and version records. The base collection contains 83 sources; parity, construction
and uniformity supplements are in separate addenda. The principal local
sources are listed below; full paths
are in the source map.
\begin{longtable}{p{.07\textwidth}p{.58\textwidth}p{.25\textwidth}}
\toprule ID & Source basename & Use\\\midrule\endhead
S01 & \texttt{ordinary-integer-polynomial.js.txt} & Outer construction\\
S02 & \texttt{rational-circuit-translator.js.txt} & Substitution semantics\\
S03 & \texttt{finite-polynomial-expressions.js.txt} & Fixed constructors\\
S04 & \texttt{variable-shift-polynomial-library.js.txt} & Close and Test\\
S05 & \texttt{elliptic-generator-proof.txt} & Generator certificate\\
S06 & \texttt{explicit-elliptic-data.txt} & Rank, torsion, codes\\
S07 & \texttt{verify-elliptic-generator.js.txt} & Exact finite algorithm\\
S08 & \texttt{six-contact-height-proof.txt} & Contacts and constants\\
S09 & \texttt{z-in-q-height-independent-audit.txt} & Height proof\\
S10 & \texttt{explicit-pole-predicate.txt} & Pole polynomial\\
S11 & \texttt{z-in-q-pole-completeness-audit.txt} & Simultaneous descent\\
S12 & \texttt{parity-frozen-4eaee875.tex} & Raw local matrices\\
S13 & \nolinkurl{multiplication-certificate-correctness-proof.txt} & Multiplication proof\\
S14 & \texttt{twenty-two-prime-pattern-proof.txt} & Selected primes\\
S15 & \texttt{z-in-q-variable-shift-bridge.txt} & Ordinary-value graph\\
S16 & \texttt{z-in-q-ordinary-integer-assembly-proof.txt} & Outer composition\\
S17 & \texttt{z-in-q-restricted-converse-assembly.txt} & RC assembly inputs\\
S18 & \texttt{ring-limits-frozen-4eaee875.tex} & Finite models and limits\\
S19 & \texttt{paired-functional-and-uniformity-proof.txt} & Paired functional\\
S20 & \texttt{determinant-transfer-audit.txt} & Clearing and transfer\\
S21 & \texttt{variable-shift-precision.py.txt} & Exact decision program\\
S22 & \texttt{precision-register.js.txt} & Ideal IR and lowering\\
S23 & \texttt{specialize-precision.js.txt} & Wrapper specialization\\
S24 & \texttt{integer-halting-correctness-proof.txt} & Natural compiler\\
S25 & \texttt{z-in-q-explicit-cutoff-bound.txt} & Literal cutoff\\
S26 & \texttt{z-in-q-polynomial-size-bound-proof.txt} & Syntax/norm argument\\
S27 & \texttt{z-in-q-numerical-polynomial-bounds.txt} & Tower comparison\\
S28 & \texttt{determinant-switch-audit.txt} & Switches and outer rank\\
S29 & \texttt{ring-determinants-frozen-4eaee875.tex} & Ring data and derivatives\\
S30 & \texttt{preliminaries-frozen-4eaee875.tex} & Selmer coranks and lengths\\
S31 & \texttt{fixed-certificate-constants.js.txt} & Exact least-multiple search\\
S32 & \texttt{pointwise-frozen-4eaee875.tex} & Distinct point indices\\
\bottomrule
\end{longtable}
In particular S30, S18, S29 and S32 are, respectively, Sections
2, 4, 5 and 9 of OpenAI's \emph{A pointwise 2-converse for elliptic
curves with rational two-torsion}~\cite{openaipointwise}, frozen at
\texttt{openai/math} commit
\nolinkurl{4eaee87548250b2682affda7d05d947c05554bae}.
The source filenames are provenance labels for those sections, not
independent publications or independent arithmetic assumptions.

The supplement manifest also records the revised fixed-degree CM proof,
the complete natural compiler proof, actual arithmetic and analytic
producer arguments, finite-Artin and determinant-count clarifications,
the fixed BFH witness application, exact substitution checks, and the
new fixed-form, weighted-theta, operator and normalization derivations.
For the height estimate, the derivation in Section~\ref{sec:height}
supersedes the historical arithmetic-height arguments in S08--S09 and
the earlier CM supplements. Those historical alternatives are not
additional premises. Primary theorem applications remain explicit in
the mathematical text.
The version records identify the source files used in the construction.

% End 05-bounds-appendices.tex


% Begin 06-literal-construction.tex
\section{Literal finite-constructor specification}\label{app:literal}
The supplementary file \nolinkurl{ordinary-integer-definition.json}
is part of this specification. It is a finite directed acyclic description,
not an expanded monomial list. The following rules determine one polynomial
from its bytes. Diagnostic labels and intended-theorem strings add no equations.

The root is \texttt{integers}. Variable zero is the external input and
variables one through $292890231904$ are its rational witnesses.
A nonnegative point reference $r$ means $X_r,\ldots,X_{r+4}$.
References $-1,-2,-3$ respectively mean
\[
 (1,0,0,0,0),\quad(0,g_0,-145,g_1,g_2),\quad
                      (0,g_0,-145,-g_1,-g_2).
\]
Every model recipe conjoins
$8g_0-1641=0$, $4g_1+15465=0$, $32g_2-87483=0$.

\subsection{Scalar nodes and fixed expressions}
Read nodes in their stored topological order.
A \texttt{const} is its decimal integer; \texttt{var} is its local variable;
\texttt{add} and \texttt{mul} are ordinary polynomial operations.
For a \texttt{call} to a template with $p$ parameters, substitute the
specified node polynomials for variables $0,\ldots,p-1$, and $X_{w+j}$
for variable $p+j$, where $w$ is the witness start. The value is the
callee's entire root polynomial.

A \texttt{fixed\_integer} evaluates a finite expression using literals,
references to preceding definitions, addition, multiplication,
nonnegative powers, maximum, factorial, the catalogue product or maximum,
and the least formal multiple. The catalogue uses the first 22 primes
in increasing order in each unit class. The searches are deterministic
and terminate by Section~\ref{sec:size}. They produce constants, not
quantified witnesses.

A \texttt{psi\_component} uses the division-polynomial recurrence in
Section~\ref{sec:syntax}, pair multiplication
$(a,b)(c,d)=(ac+2bd,ad+bc)$ and $A=(-11243,7950)$.
It returns component zero or one; every index is fixed.
The positive fixed point codes use
\[
 x_n=x(P_0)-\frac{\psi_{n+1}(P_0)\psi_{n-1}(P_0)}{\psi_n(P_0)^2},
 \qquad y_n=\frac{\psi_{2n}(P_0)}{2\psi_n(P_0)^4}.
\]
Nontorsion makes these fixed denominators nonzero.
Negative indices negate both ordinate components.

A \texttt{finite\_guard\_product} ranges over
$(a,b,c)\in[-T,T]^3\cap\Z^3$ with $ab\ne c$, with fixed $T$.
Write the fifteen coordinates of the three codes as reduced
$A_j/B_j$, $B_j>0$, and put $D=\prod_jB_j$. Its factor is
\[
                 \sum_{j=1}^{15}(DX_j-(D/B_j)A_j)^2.
\]
Multiply these factors, with empty product one. This is an integer
polynomial with precisely the finite false-triple zero set used above.

\subsection{Recipe constraints and witness scope}
A model-evaluation recipe sums squares of the generator equations,
listed natural calls, scalar-constant calls, add/multiply/power-value calls,
and the five coordinate differences of \texttt{rootEquality}.
A source-variable row with a non-null natural witness start contributes
exactly that natural call. Operation kinds zero, one and two select
addition, multiplication and power value. Source-node and provenance
fields add no equations.

A sum-of-squares composition sums squares of its callee polynomials;
a product composition multiplies them. Parameter \texttt{[var,i]} means
$X_i$. Each private block uses its recorded start and length.
The checked ranges are disjoint from parameters and other private blocks.
These rules include globally conjoined fixed equations in active callees.

\subsection{The saved templates and the mathematical predicates}
\label{subsec:v10-template-dictionary}
The following dictionary identifies the relations encoded by the saved
templates, in addition to the decoding rules above. Point arguments are
canonical five-coordinate codes. Field values have their two rational
components in the basis $(1,\sqrt2)$; conjugation negates the second
component. An inverse equation guards each displayed quotient.
Conjunction is encoded by a sum of squares, and disjunction by a
product, with the private witness blocks specified in the file.
The identities are identities of zero sets over $\Q$; squaring a
callee polynomial can change the literal polynomial without changing
the relation it defines.

\begin{longtable}{p{.29\textwidth}p{.65\textwidth}}
\toprule Saved template & Relation and constraints\\\midrule\endhead
\nolinkurl{member} & Membership in the canonical code domain $D=2E(K)$
of Section~\ref{sec:model}.\\
\nolinkurl{raw_add}, \nolinkurl{raw_z}, \nolinkurl{raw_inverse_x} &
The group-law relation on typed inputs, $z=-x/y$, and $w=1/x$,
respectively, with $z(O)=w(O)=0$.\\
\nolinkurl{log25} & The two-component equality to the displayed
scaled degree-25 polynomial $L$ in Section~\ref{sec:mult}.\\
\nolinkurl{pole} & $Q_{\rm pole}(B)^2$, where $Q_{\rm pole}$ is the
eleven-witness polynomial of Section~\ref{sec:pole}. The outer square
comes from the singleton conjunction added by the rational builder.\\
\nolinkurl{test_group} & The twelve field tests of
Section~\ref{sec:mult}, including $A_i-\sigma(A'_i)$ for $i=1,2,3$.
Its four input points are $[c_*]P_0,[c_*]P_1,[c_*]P_2,[c_*]P_3$.
It types the first four auxiliary points and obtains the second four
by subtraction; its two denominator points and logarithms are nonzero.\\
\nolinkurl{slope} & The normalized twelve-test products in six private
groups, their traced inverses, and $\Phi(R/\Theta^2)$ from
Section~\ref{sec:mult}, with all contact and denominator guards.\\
\nolinkurl{height_certificate} & The 27 slopes summing in triples to
the nine homogeneous coordinate quantities in that section, together
with membership of the three input points.\\
\nolinkurl{multiply.scalar}, \nolinkurl{multiply.guard},
\nolinkurl{multiply} & The fixed $[c_*]$ graph; the inverse of the
finite false-triple product; and their conjunction with the height
certificate at the fixed $\ell_0$.\\
\nolinkurl{natural} & Four model squares, formed by multiplication,
whose sum is the input code, as in Section~\ref{sec:compiler}.\\
\nolinkurl{value.signed_encode} & $N=\iota(n)$ and a natural code
$C=\iota(c)$ with $n=c/2$ for even $c$, and $n=-(c+1)/2$ for odd
$c$. The two branches use $c=2q$ and $c=2q+1$ with $q\ge0$.\\
\nolinkurl{value.abs} & $a=|x|$: $a$ is a sum of four rational
squares and $a^2=x^2$.\\
\nolinkurl{value.close_coordinates} & The strict coordinate window
of Section~\ref{sec:bridge}, using the typed points $A,L,J$,
$r=\operatorname{Tr}(z(J)/z(A))/2$ and
$\epsilon=|(z(L))_0|+|(z(L))_1|>0$, with $|t-r|<\epsilon$.
This template alone does not assert the precision relation.\\
\nolinkurl{value.test} & Exactly $\Test(N,\kappa,t)$ of
Section~\ref{sec:bridge}: $C_1=[\kappa n]P_0$,
$V_0=H-\kappa P_0$, $V_1=U-C_1$, the six displayed raw-$z$ tests,
their normalization by $F(t)$, and the single guarded pole input.
These ratios use $z$, not $L$.\\
\nolinkurl{value.exceptions} & The six branches
$N=\iota(c)$ and $t=c$, $c\in\Cset$, with the constant points
computed by the fixed group-law chains.\\
\nolinkurl{scalar.integer_constant_}\,$i$ & The fixed integer
scalar graph requested by the translated precision computation.\\
\bottomrule
\end{longtable}

The composition recipes make the remaining identifications explicit.
\texttt{power\_value} translates the source power template with its
three input codes and output code; \texttt{precision} translates
the five-natural-input precision polynomial, in the order
$(n,\kappa,a,\ell,j)$ after signed encoding. The recipe
\texttt{close} introduces $A,L,J$, forms those five signed encodings,
conjoins \texttt{precision}, and calls
\nolinkurl{value.close_coordinates}. Thus it is the full $\Close$.
The recipe \texttt{main} introduces one shared $\kappa$ and conjoins
\texttt{close} and \nolinkurl{value.test} at exactly that same code.
The product recipe \texttt{value\_union} joins this branch with
\nolinkurl{value.exceptions}; \texttt{value} also imposes
\texttt{member} on $N$. Finally \texttt{integers} existentially
introduces $N$ and calls \texttt{value} with external coordinate
$t=X_0$. These are precisely $\Val$ and its input projection.
The ordered calls and private ranges are checked by
\nolinkurl{auditOrdinaryIntegerPolynomial} in
\nolinkurl{construction/ordinary-integer-polynomial.js}.
\nolinkurl{scripts/check_construction.py} regenerates the saved
expression byte for byte. The dictionary supplies the mathematical
identification; those finite checks verify the stored syntax and
wiring, not the arithmetic theorems used to prove the predicates.

The following parameter and witness counts are extracted from the exact
specification, without expanding its polynomial.
\begin{longtable}{p{.47\textwidth}rrr}
\toprule Template & Parameters & Witnesses & Root node\\\midrule\endhead
\nolinkurl{member} & 5 & 7 & 149\\
\nolinkurl{raw_add} & 15 & 6 & 308\\
\nolinkurl{raw_z} & 7 & 1 & 68\\
\nolinkurl{raw_inverse_x} & 7 & 0 & 54\\
\nolinkurl{log25} & 4 & 0 & 229\\
\nolinkurl{pole} & 1 & 11 & 106\\
\nolinkurl{test_group} & 44 & 162 & 515\\
\nolinkurl{slope} & 22 & 1,298 & 2681\\
\nolinkurl{height_certificate} & 36 & 35,106 & 260\\
\nolinkurl{multiply.scalar} & 10 & 1 & 136\\
\nolinkurl{multiply.guard} & 15 & 1 & 22\\
\nolinkurl{multiply} & 15 & 35,134 & 83\\
\nolinkurl{natural} & 5 & 140,604 & 79\\
\nolinkurl{value.signed_encode} & 10 & 140,654 & 138\\
\nolinkurl{value.abs} & 2 & 4 & 23\\
\nolinkurl{value.close_coordinates} & 16 & 53 & 271\\
\nolinkurl{value.test} & 11 & 35,282 & 775\\
\nolinkurl{value.exceptions} & 6 & 406 & 487\\
\nolinkurl{scalar.integer_constant_0} & 10 & 1 & 136\\
\nolinkurl{scalar.integer_constant_1} & 10 & 1 & 136\\
\nolinkurl{scalar.integer_constant_2} & 10 & 1 & 136\\
\nolinkurl{scalar.integer_constant_3} & 10 & 1 & 136\\
\nolinkurl{scalar.integer_constant_4} & 10 & 1 & 136\\
\bottomrule\end{longtable}

\begin{longtable}{lrrr}
\toprule Recipe & Parameters & Witnesses & Total variables\\\midrule\endhead
\nolinkurl{power_value} & 20 & 17,786,698 & 17,786,718\\
\nolinkurl{precision} & 25 & 292,889,492,836 & 292,889,492,861\\
\nolinkurl{close} & 11 & 292,890,196,199 & 292,890,196,210\\
\nolinkurl{main} & 6 & 292,890,231,486 & 292,890,231,492\\
\nolinkurl{value_union} & 6 & 292,890,231,892 & 292,890,231,898\\
\nolinkurl{value} & 6 & 292,890,231,899 & 292,890,231,905\\
\nolinkurl{integers} & 1 & 292,890,231,904 & 292,890,231,905\\
\bottomrule\end{longtable}


The five scalar-constant templates use the values in
\texttt{library.fixedConstants}. Large program literals are retained
byte-for-byte in the supplementary specification. Replacing them by
unspecified sufficiently large constants would change this definition.

